% ============================================================================
% Main Document: Implementing Rust Support for Embedded Hardware
% A Comprehensive Guide Using EFR32MG24
% ============================================================================
\documentclass[11pt,a4paper,twoside,openright]{book}

% ========== PREAMBLE (Configuration and Packages) ==========
% Minimal Preamble - Version 2.1
% Technical book template optimized for code and diagrams
% Last updated: 2026-01-09

% ========== FONTS AND ENCODING ==========
\usepackage[T1]{fontenc}
\usepackage[utf8]{inputenc}
\usepackage{lmodern}

% ========== PAGE LAYOUT ==========
\usepackage[margin=1in]{geometry}

% ========== COLORS (before other packages that use colors) ==========
\usepackage[dvipsnames]{xcolor}
\definecolor{rustblue}{RGB}{0, 119, 255}
\definecolor{codebg}{RGB}{248, 248, 248}

% ========== CODE LISTINGS ==========
% minted provides fancyvrb, so we get it for free
\usepackage{minted}
\usemintedstyle{friendly}
\setminted{
    fontsize=\footnotesize,
    baselinestretch=0.9,
    bgcolor=codebg,
    linenos,
    numbersep=5pt,
    frame=none,
    framesep=10pt,
    breaklines=true,
    breakanywhere=true
}

% ========== HEADERS AND FOOTERS ==========
\usepackage{fancyhdr}
\pagestyle{fancy}
\fancyhf{}
\fancyhead[LE]{\thepage\quad\leftmark}
\fancyhead[RO]{\rightmark\quad\thepage}
\fancyfoot[C]{}
\renewcommand{\headrulewidth}{0.4pt}
\setlength{\headheight}{25.5pt}  % Fix fancyhdr warning

% ========== GRAPHICS ==========
\usepackage{graphicx}
\usepackage{tikz}
\usetikzlibrary{arrows.meta}    % For arrow tips like Latex
\usetikzlibrary{positioning}    % For positioning syntax (of=)
\usepackage{forest}      % For directory tree diagrams

% ========== TABLES ==========
\usepackage{array}
\usepackage{booktabs}    % Professional tables

% ========== MARGIN NOTES ==========
\usepackage{marginnote}
\reversemarginpar
\setlength{\marginparwidth}{1.2in}

% ========== BIBLIOGRAPHY ==========
\usepackage[style=ieee, backend=biber]{biblatex}

% ========== GLOSSARY ==========
\usepackage[toc, nopostdot]{glossaries}

% ========== TABLE OF CONTENTS ==========
\usepackage{minitoc}
\dominitoc

% ========== TITLE FORMATTING ==========
\usepackage{titlesec}
% Chapter titles - slightly smaller
\titleformat{\chapter}[block]
  {\Large\bfseries}
  {\thechapter.\quad}
  {0pt}
  {\LARGE}

% Section titles - smaller
\titleformat{\section}
  {\large\bfseries}
  {\thesection\quad}
  {0pt}
  {}

% Subsection titles - smaller
\titleformat{\subsection}
  {\normalsize\bfseries}
  {\thesubsection\quad}
  {0pt}
  {}

% ========== HYPERLINKS (load late to avoid conflicts) ==========
\usepackage[pdfborder={0 0 0}, breaklinks=true]{hyperref}
\hypersetup{
    colorlinks=true,
    linkcolor=rustblue,
    filecolor=magenta,
    urlcolor=rustblue,
    pdftitle={Implementing Rust Support for Embedded Hardware},
    pdfauthor={Marcelo Correa},
    pdfsubject={Embedded Systems, Rust Programming},
    pdfkeywords={Rust, Embedded, EFR32MG24, HAL, PAC, Wireless}
}

% ========== CUSTOM ENVIRONMENTS ==========
\usepackage{mdframed}
\usepackage{enumitem}

\newenvironment{example}
  {\begin{mdframed}[backgroundcolor=gray!10, linewidth=0.5pt]}
  {\end{mdframed}}

\newenvironment{notebox}
  {\begin{mdframed}[backgroundcolor=blue!10, linewidth=0.5pt]}
  {\end{mdframed}}

\newenvironment{warningbox}
  {\begin{mdframed}[backgroundcolor=red!10, linewidth=0.5pt]}
  {\end{mdframed}}

\newenvironment{tipbox}
  {\begin{mdframed}[backgroundcolor=green!10, linewidth=0.5pt]}
  {\end{mdframed}}

\newenvironment{remark}
  {\begin{mdframed}[backgroundcolor=yellow!10, linewidth=0.5pt]}
  {\end{mdframed}}

\newenvironment{keytakeaways}
  {\begin{itemize}[leftmargin=*]}
  {\end{itemize}}

\newenvironment{exercises}
  {\begin{enumerate}[leftmargin=*]}
  {\end{enumerate}}

\newenvironment{exercise}[1][Easy]
  {\paragraph{Exercise (#1):}}
  {}

% ========== CUSTOM COMMANDS ==========
% Inline code with Rust syntax highlighting
\newcommand{\rustcode}[1]{\mintinline{rust}{#1}}

% Code inline - uses fancyvrb's \Verb for proper verbatim handling
% This handles underscores, hashes, and all special characters
\usepackage{xparse}
\NewDocumentCommand{\codeinline}{v}{%
  \texttt{#1}%
}

% Crate names (same verbatim handling)
\NewDocumentCommand{\crate}{v}{%
  \texttt{#1}%
}

% File paths (same verbatim handling)
\NewDocumentCommand{\filepath}{v}{%
  \texttt{#1}%
}

% Structure/concept names (just italic, no special handling)
\newcommand{\structure}[1]{\textit{#1}}

% Chapter thumb placeholder
\newcommand{\chapthumb}[1]{}

% Listings list stub (we use minted, not listings package)
\newcommand{\lstlistoflistings}{%
  \chapter*{List of Code Listings}%
  \markboth{List of Code Listings}{List of Code Listings}%
}

% ========== NUMBERING DEPTH ==========
\setcounter{secnumdepth}{2}  % Number sections and subsections
\setcounter{tocdepth}{2}     % Show sections and subsections in TOC

% ========== UNICODE CHARACTERS THE BOOK USES ==========
% T1/lmodern does not carry every glyph the text uses: an em-dash or a bullet
% disappears silently (pdflatex also refuses them outright). Mapping them to LaTeX
% commands is one fix in one place, and it keeps the source readable.
\usepackage{newunicodechar}
\newunicodechar{—}{---}
\newunicodechar{–}{--}
\newunicodechar{•}{\textbullet}
\newunicodechar{…}{\ldots}
\newunicodechar{×}{\ensuremath{\times}}
\newunicodechar{→}{\ensuremath{\rightarrow}}
\newunicodechar{≈}{\ensuremath{\approx}}

% ========== MONOSPACE FONT WITH THE GLYPHS THE EXAMPLES USE ==========
% The directory trees in the code blocks use box-drawing characters (U+2500..2514)
% that T1/lmodern does not carry: they vanish without an error, leaving a tree of
% bare names. DejaVu Sans Mono has them, and XeLaTeX can select it directly.
\usepackage{fontspec}
\setmonofont{DejaVu Sans Mono}[Scale=0.86]


% ========== GLOSSARY ENTRIES ==========
\makeglossaries

\newglossaryentry{pac}{
    name=PAC,
    description={Peripheral Access Crate - Auto-generated register-level access}
}

\newglossaryentry{hal}{
    name=HAL,
    description={Hardware Abstraction Layer - Safe, ergonomic APIs}
}

\newglossaryentry{bsp}{
    name=BSP,
    description={Board Support Package - Board-specific abstractions}
}

\newglossaryentry{svd}{
    name=SVD,
    description={System View Description - XML register definitions}
}

\newglossaryentry{rail}{
    name=RAIL,
    description={Radio Abstraction Interface Layer - Multi-protocol radio}
}

\newglossaryentry{embeddedhal}{
    name={embedded-hal},
    description={Rust embedded ecosystem traits for portability}
}

% ========== TITLE PAGE ==========
\title{\textbf{Rust Embedded Hardware Support}\\\Large A Complete Guide Using the Silicon Labs EFR32MG24}
\author{Marcelo Correa\thanks{With contributions from the embedded Rust community}}
\date{2026\\Version 1.0.0}

% ========== DOCUMENT START ==========
\begin{document}

% Front Matter
\frontmatter

% Title Page
\maketitle

% Copyright page
\thispagestyle{empty}
\begin{center}
    \vspace*{\fill}
    Copyright \copyright\ 2026 Marcelo Correa\\
    Licensed under Apache-2.0 or MIT (at your option)\\
    \vspace{0.6em}
    Source code and examples:\\
    \url{https://github.com/bitscrafts/EFR32MG2X-RS}\\
    \vspace{0.6em}
    ISBN: 978-1-XXXX-XXXX-X\\
    \vspace{0.6em}
    Typeset in \LaTeX\ (pdfLaTeX, TeX Live)\\
    Document version: 1.0.0 (2026-11-11)\\
    First printing: January 2026
    \vspace*{\fill}
\end{center}
\clearpage

% Table of Contents
\tableofcontents
\clearpage

% List of Figures
\listoffigures
\clearpage

% List of Listings
\lstlistoflistings
\addcontentsline{toc}{chapter}{List of Code Listings}
\clearpage

% Preface
\chapter*{Preface}
\addcontentsline{toc}{chapter}{Preface}
\markboth{Preface}{Preface}

\section*{What This Book Is}

This book documents the complete process of creating Rust support for a new embedded microcontroller family, using the Silicon Labs EFR32MG24 as a real-world case study. It covers everything from generating a Peripheral Access Crate (PAC) from vendor System View Description (SVD) files, through implementing a Hardware Abstraction Layer (HAL), to integrating complex wireless protocol stacks.

\section*{What This Book Is Not}

This is not an introduction to Rust programming---it assumes basic Rust knowledge. It's also not a replacement for the EFR32MG24 Reference Manual or Rust Embedded Book, which it complements.

\section*{Who This Book Is For}

\begin{itemize}
    \item Embedded developers wanting to learn Rust
    \item Rust developers interested in embedded systems
    \item Engineers bringing Rust support to new hardware
    \item Technical leads evaluating Rust for embedded projects
    \item Computer science students studying systems programming
\end{itemize}

\section*{Structure of the Book}

The book is organized into four parts:

\textbf{Part I} (Chapters 1--3): Foundation and PAC Generation
\\Foundation concepts, SVD files, and build system setup.

\textbf{Part II} (Chapters 4--9): Hardware Abstraction Layer
\\Peripheral implementations (GPIO, clocks, communication, advanced features).

\textbf{Part III} (Chapters 11--14): Wireless Protocol Stacks
\\RAIL foundation, Bluetooth LE, Zigbee, Thread, and Matter integration.

\textbf{Part IV} (Chapters 10, 15--18): Testing, BSP, and Production
\\Test-driven development, board support packages, documentation, and deployment.

\section*{Prerequisites}

Required knowledge:
\begin{itemize}
    \item Basic Rust programming (ownership, traits, modules)
    \item Familiarity with command-line tools
    \item Understanding of binary/hexadecimal
    \item Basic electronics knowledge (helpful but not required)
\end{itemize}

Required hardware:
\begin{itemize}
    \item Seeed Studio XIAO MG24 Sense development board
    \item USB-C cable (data-capable, not just power)
    \item Computer running Linux, macOS, or Windows
\end{itemize}

Optional but recommended:
\begin{itemize}
    \item Logic analyzer (for debugging communication protocols)
    \item Oscilloscope (for signal quality analysis)
    \item Multiple XIAO MG24 boards (for testing wireless protocols)
\end{itemize}

\section*{How to Read This Book}

This book can be read in two ways:

\textbf{Linear Approach}: Read cover-to-cover for complete understanding. Each chapter builds on previous ones.

\textbf{Reference Approach}: Jump to specific chapters for:
\\- SVD generation → Chapter 2
\\- GPIO implementation → Chapter 4
\\- I2C protocol details → Chapter 7
\\- Bluetooth LE → Chapter 12
\\- Testing strategies → Chapter 10

All code examples are available in the accompanying repository.

\section*{Acknowledgments}

This book and the accompanying EFR32MG24 implementation would not be possible without:
\\- The Rust Embedded Working Group and their excellent tools
\\- The stm32-rs, nrf-rs, and esp-rs communities for pioneering work
\\- Silicon Labs for the excellent EFR32MG24 documentation
\\- The embedded Rust community contributors

\section*{Feedback and Contributions}

Please report issues, suggest improvements, or contribute at:
\\\url{https://github.com/bitscrafts/EFR32MG2X-RS}

\vspace*{2em}
\begin{flushright}
Marcelo Correa\\
January 2026
\end{flushright}

\clearpage

% Main Matter
\mainmatter

% ========== PART I: PAC GENERATION AND FOUNDATION ==========
\part{PAC Generation and Foundation}
\label{part:pac-foundation}

% Chapter 1: Introduction
% Chapter 1: Introduction to Embedded Rust and Project Architecture
% File: chapters/chapter-01-introduction.tex

\chapter{Introduction to Embedded Rust and Project Architecture}
\label{ch:chapter-01}
\chapthumb{1}

\begin{quote}
    \textit{"The biggest memory safety vulnerability is the assumption that memory safety isn't important for embedded systems."}
    \\
    -- Embedded Systems Security Research, 2023
\end{quote}

This chapter introduces embedded Rust development and establishes the foundation for understanding the PAC/HAL/BSP architecture pattern. We'll explore why Rust is compelling for embedded systems, create a bare-metal program, and set up the complete development environment.

\minitoc

\clearpage

% ============================================================================
% SECTION 01.01: Why Rust for Embedded Systems
% ============================================================================
% Section 1.1: Why Rust for Embedded Systems

Complex software systems, especially those dealing with direct hardware manipulation, memory management, and concurrent interrupts, are notoriously difficult to write correctly. Traditional embedded development often relies on C or C++, languages that provide raw power and direct hardware access but lack modern safety guarantees. Memory safety bugs---buffer overflows, use-after-free, null pointer dereferences, and data races---account for the majority of security vulnerabilities and system crashes in embedded systems.

\section{Rust's Value Proposition for Embedded}
\label{sec:why-rust-value-prop}

Rust addresses these challenges through a unique combination of features that make it exceptionally well-suited for embedded systems programming:

\subsection{Memory Safety Without Garbage Collection}

Traditional memory-safe languages rely on garbage collection, which introduces unpredictable pauses and significant runtime overhead---unacceptable constraints for resource-constrained microcontrollers. Rust achieves memory safety through its ownership system, enforced at compile time with zero runtime cost.

\begin{notebox}
    \textbf{Key Insight}: Rust's borrow checker prevents memory safety bugs at compilation, eliminating entire classes of vulnerabilities before the code ever reaches the hardware.
\end{notebox}

Consider this common C bug pattern:

\begin{minted}{c}
void process_data(uint8_t *buffer, size_t len) {
    uint8_t local_buffer[16];
    // Buffer overflow if len > 16!
    memcpy(local_buffer, buffer, len);
    // ... process data
}
\end{minted}

The same logic in Rust:

\begin{minted}{rust}
fn process_data(buffer: &[u8]) {
    // Compile error: attempt to copy 32 bytes into 16-byte array
    let mut local_buffer = [0u8; 16];
    local_buffer.copy_from_slice(buffer); // Panics if buffer.len() != 16
    // Or use safe slice copying:
    let count = buffer.len().min(16);
    local_buffer[..count].copy_from_slice(&buffer[..count]);
}
\end{minted}

The Rust compiler enforces bounds checking at compile time or generates explicit panics for runtime bounds violations, preventing undefined behavior.

\subsection{Zero-Cost Abstractions}

Rust's high-level abstractions compile down to the same efficient machine code as hand-written C. This is critical for embedded systems where every byte and CPU cycle matters.

\begin{example}
A GPIO pin abstraction in Rust:

\begin{minted}{rust}
let pin = gpio.pin0.into_output();
pin.set_high();
\end{minted}

This compiles to a single register write instruction, identical to:

\begin{minted}{c}
GPIO->DOUT |= (1 << 0);
\end{minted}
\end{example}

The Rust code provides type safety, compile-time mode checking, and clear semantics while maintaining identical performance to the raw C version.

\subsection{Fearless Concurrency}

Embedded systems increasingly require concurrent operation: interrupt handlers, DMA transfers, and multiple sensors polling simultaneously. Rust's ownership model prevents data races at compile time.

\begin{minted}{rust}
// This compiles safely
static mut COUNTER: u32 = 0;

#[interrupt]
fn TIMER0() {
    // Safe because we use atomic operations
    cortex_m::interrupt::free(|_| {
        unsafe { COUNTER += 1; }
    });
}

// This fails to compile!
fn increment_counter() {
    COUNTER += 1; // Error: cannot access mutable static without unsafe
}
\end{minted}

The compiler forces developers to explicitly handle synchronization, making race conditions nearly impossible.

\subsection{Comparison with Traditional Embedded Languages}
\label{sec:language-comparison}

\subsection{Rust vs C}

\begin{table}[H]
    \centering
    \begin{tabular}{l c c}
        \toprule
        \textbf{Feature} & \textbf{C} & \textbf{Rust} \\
        \midrule
        Memory Safety & Manual (runtime errors) & Compile-time guaranteed \\
        Concurrency Safety & Manual (locks, careful design) & Compile-time guaranteed \\
        Type Safety & Weak (void*, casts) & Strong (no implicit casts) \\
        Null Pointers & Always possible & Option<T> type \\
        Error Handling & Error codes & Result<T, E> type \\
        Abstraction Cost & Potentially high & Zero-cost \\
        Learning Curve & Moderate & Steep (ownership concepts) \\
        Ecosystem & Mature, established & Growing rapidly \\
        Real-time Suitability & Excellent (predictable) & Excellent (no GC) \\
        \bottomrule
    \end{tabular}
    \caption{Rust vs C for embedded development}
    \label{tab:rust-vs-c}
\end{table}

\subsection{Real-World Impact}

The Linux kernel, which has historically been written entirely in C, began accepting Rust code in version 6.1 (2022). This demonstrates Rust's maturity for systems programming even at the highest levels.

In the embedded space, companies are reporting:

\begin{itemize}
    \item 40-50% reduction in debugging time
    \item Near elimination of memory corruption bugs
    \item Improved team productivity through better compiler diagnostics
    \item Easier onboarding for new developers
\end{itemize}

\subsection{Rust's Embedded Ecosystem}
\label{sec:rust-embedded-ecosystem}

\subsection{Core Crates}

The embedded Rust ecosystem provides standardized interfaces:

\begin{itemize}
    \item \crate{embedded-hal}: Hardware abstraction layer traits for portability
    \item \crate{cortex-m}: Low-level Cortex-M register access
    \item \crate{cortex-m-rt}: Startup code and interrupt handling
    \item \crate{cortex-m-semihosting}: Debug output via debugger
    \item \crate{svd2rust}: PAC generation from SVD files
    \item \crate{panic-halt}: Panic handler for minimal footprint
\end{itemize}

\subsection{Tooling Integration}

Rust integrates seamlessly with embedded development workflows:

\begin{itemize}
    \item \textbf{cargo}: Built-in package manager and build system
    \item \textbf{probe-rs}: Flash programming and debugging
    \item \textbf{defmt}: Efficient logging for resource-constrained targets
    \item \textbf{flip-link}: Linker wrapper for stack overflow protection
\end{itemize}

\subsection{Addressing Embedded Development Challenges}
\label{sec:embedded-challenges}

\subsection{Resource Constraints}

Rust's ownership system enables compile-time memory management without runtime overhead. The \codeinline{no_std} environment runs without heap allocation, relying entirely on stack allocation and static data.

\begin{minted}{rust}
#![no_std]  // No standard library (no heap allocation)
#![no_main] // Disabled main wrapping

// All memory is statically allocated
static BUFFER: [u8; 1024] = [0; 1024];

#[no_mangle]
pub extern "C" fn main() -> ! {
    loop {
        // Application logic
    }
}
\end{minted}

\subsection{Deterministic Behavior}

Embedded systems require predictable timing. Rust provides:

\begin{itemize}
    \item No hidden allocations or deallocations
    \item No garbage collection pauses
    \item Explicit control over memory layout
    \item \codeinline{const fn} for compile-time computation
\end{itemize}

\subsection{Hardware Interface Safety}

Direct hardware register manipulation is inherently unsafe. Rust forces developers to mark unsafe operations explicitly:

\begin{minted}{rust}
pub fn enable_gpio_clock() {
    unsafe {
        (*CMU::ptr()).hfbusclken0.modify(|_, w|
            w.gpio().set_bit()
        );
    }
}
\end{minted}

The \codeinline{unsafe} block signals to developers and reviewers that manual verification is required.

\subsection{The EFR32MG24: Ideal for Rust}
\label{sec:efr32mg24-for-rust}

The Silicon Labs EFR32MG24 is particularly well-suited for Rust development:

\begin{itemize}
    \item \textbf{ARM Cortex-M33}: Modern architecture with TrustZone-M
    \item \textbf{Memory Protection Unit}: Hardware-enforced memory isolation
    \item \textbf{Deterministic Architecture}: No unpredictable timing behavior
    \item \textbf{Rich Peripherals}: 133 peripherals across multiple domains
    \item \textbf{Unified Radio}: Multi-protocol stack support
    \item \textbf{Hardware Cryptography}: Secure element integration
\end{itemize}

\subsection{Summary}
\label{sec:why-rust-summary}

Rust brings modern software engineering practices to embedded development without sacrificing performance or control. Its compile-time safety guarantees eliminate entire classes of bugs that plague C-based embedded systems, while its zero-cost abstractions enable expressive, maintainable code that compiles to efficient machine instructions.

The journey through this book will demonstrate:

\begin{enumerate}
    \item How to generate a Peripheral Access Crate (PAC) from vendor SVD files
    \item Building safe, ergonomic HAL abstractions on top of raw registers
    \item Implementing standard embedded-hal traits for ecosystem compatibility
    \item Creating board support packages for specific hardware
    \item Integrating complex protocol stacks (Bluetooth, Zigbee, Thread)
    \item Testing strategies for embedded systems
    \item Production deployment and maintenance
\end{enumerate}

Along the way, we'll explore the unique challenges and solutions that emerged during the real-world implementation of EFR32MG24 support in Rust.
\clearpage


% ============================================================================
% SECTION 01.02: Your First Bare-Metal Rust Program
% ============================================================================
\section{Your First Bare-Metal Rust Program}
\label{sec:bare-metal-program}

% Section 01.02: Your First Bare-Metal Rust Program

After understanding why Rust is compelling for embedded systems, let's create our first bare-metal Rust program. This example will run without an operating system, directly on the microcontroller hardware.

\subsection{Creating a Minimal Embedded Project}

First, create a new directory for our project:

\begin{minted}{bash}
$ mkdir blink_bare_metal
$ cd blink_bare_metal
$ cargo init --name blink_bare_metal
\end{minted}

This creates a standard Rust project structure. However, for embedded development, we need to modify it significantly.

\subsection{Configuring for bare\_metal: \texttt{no\_std}}

Before we begin, let's understand what a bare-metal embedded project looks like. Edit \filepath{Cargo.toml}:

\begin{minted}{toml}
[package]
name = "blink_bare_metal"
version = "0.1.0"
edition = "2021"

[profile.dev]
opt-level = "z"          # Optimize for size
debug = true

[profile.release]
opt-level = "z"          # Optimize for size
panic = "abort"          # Simplest panic handler

[dependencies]
# We'll add dependencies as needed
\end{minted}

Now modify \filepath{src/main.rs}:

\begin{minted}{rust}
#![no_std]    // Don't link the standard library (no heap, no OS)
#![no_main]   // We define our own entry point

// Prevent Rust from generating main() wrapper
use core::panic::PanicInfo;

/// Our entry point - Cortex-M calls this on reset
#[no_mangle]
pub extern "C" fn main() -> ! {
    // Application logic will go here

    loop {
        // Infinite loop to prevent return
    }
}

/// Simple panic handler - just halt
#[panic_handler]
fn panic(_info: &PanicInfo) -> ! {
    loop {}
}
\end{minted}

\subsection{What is no\_std?}

\begin{notebox}
    \textbf{no\_std} means ``no standard library.'' Our program cannot use:
    \begin{itemize}
        \item Heap allocation (Box, Vec, String)
        \item File I/O
        \item Threads
        \item Networking
    \end{itemize}

    We can still use: \codeinline{core}, slices, static arrays, and direct hardware access.
\end{notebox}

The standard library requires an operating system for features like heap allocation and threads. In bare-metal embedded, we have no OS, so we use only \codeinline{core}, which provides fundamental types and operations that don't require an OS.

\subsection{The Entry Point}

In hosted environments (with an OS), Rust generates a C main() function that calls your main(). In bare metal:

\begin{itemize}
    \item \codeinline{#[no_mangle]} prevents Rust from changing the function name
    \item \codeinline{pub extern "C" fn} ensures the correct calling convention
    \item \codeinline{-> !} means "never returns" (diverging function)
    \item We must provide an infinite loop or halt instruction
\end{itemize}

\subsection{Cross-Compilation}

Our development machine (likely x86\_64) can't execute ARM Cortex-M code. We need to cross-compile:

\begin{minted}{bash}
$ rustup target add thumbv8m.main-none-eabihf
\end{minted}

This installs the Rust standard library for the ARM Cortex-M33 architecture without an OS.

Create \filepath{.cargo/config.toml}:

\begin{minted}{toml}
[build]
target = "thumbv8m.main-none-eabihf"

[target.thumbv8m.main-none-eabihf]
# No `runner` on purpose: `cargo run` would program the chip the moment it builds,
# and on this board that has to go through the guarded script (see the note below).
rustflags = [
  "-C", "link-arg=-Tlink.x",
  "-C", "link-arg=--nmagic",
]

[env]
DEFMT_LOG = "debug"
\end{minted}

\begin{quote}
\textbf{Why this book has no cargo run.} The first 24 KB of flash
belong to the Gecko bootloader. A `runner` in `.cargo/config.toml` turns every
\codeinline{cargo run} into a write, and a single wrong linker script silently relocates
the vector table into that region. This project flashes through
\codeinline{scripts/flash-lnx.sh}: dry run by default, \codeinline{--confirm} to write,
a hard refusal when the image's lowest address is below \codeinline{0x08006000}, and a
read-back verification afterwards. The same script is what re-flashes the board after a
mistake — which matters more than the convenience of \codeinline{cargo run}.
\end{quote}

\subsection{Building the Project}

\begin{minted}{bash}
$ cargo build --release
   Compiling blink_bare_metal v0.1.0
    Finished release [optimized] target(s) in 0.5s
\end{minted}

Check the result:

\begin{minted}{bash}
$ cargo size --release
   text    data     bss     dec     hex filename
     64       0       0      64      40 blink_bare_metal
\end{minted}

Our minimal program is only 64 bytes!

\subsection{Analyzing the Binary}

Use \codeinline{objdump} to examine the machine code:

\begin{minted}{bash}
$ cargo objdump --release -- -d
\end{minted}

This shows the ARM assembly instructions. You'll see something like:

\begin{minted}{text}
bare_metal:     file format elf32-littlearm

Disassembly of section .text:

main:
    b #-4 <main>
    b #-4 <main>
\end{minted}

The two \codeinline{b #-4} instructions are branches to themselves---infinite loops. One is our main() function, the other is the panic handler.

\subsection{What This Demonstrates}

Even at this minimal stage, we've established:

\begin{itemize}
    \item Cross-compilation toolchain setup
    \item Bare-metal runtime configuration
    \item Minimal program that builds successfully
    \item Linker script and startup code (provided by cortex-m-rt)
\end{itemize}

The program won't do anything visible yet (no LEDs blinking), but it represents a working foundation. In the next section, we'll add actual hardware interaction.

\begin{exercise}[Easy]
    Modify the panic handler to increment a counter in an infinite loop. Use a static mut variable and examine how the compiler warns about unsafe code.
\end{exercise}

\begin{exercise}[Medium]
    Research what happens when main() (or the reset handler) returns on a Cortex-M microcontroller. Why is \codeinline{-> !} important?
\end{exercise}

\begin{exercise}[Hard]
    Create a minimal linker script (\filepath{linker.ld}) and build script (\filepath{build.rs}) to understand how the linker allocates memory for your program.
\end{exercise}
\clearpage


% ============================================================================
% SECTION 01.03: The PAC/HAL/BSP Architecture Pattern
% ============================================================================
\section{The PAC/HAL/BSP Architecture Pattern}
\label{sec:pac-hal-bsp-architecture}

% Section 01.03: The PAC/HAL/BSP Architecture Pattern

Embedded Rust development follows a layered architecture that promotes code reuse, safety, and separation of concerns. This pattern is not arbitrary—it reflects years of lessons learned from implementing hardware support across multiple vendors and microcontroller families.

\subsection{The Three-Layer Architecture}
\label{subsec:three-layer-architecture}

\begin{figure}[H]
    \centering
    \begin{tikzpicture}[
        layer/.style={rectangle, draw=black, thick, minimum width=8cm, minimum height=1.5cm, font=\bfseries},
        arrow/.style={-Latex, thick, double},
        label/.style={font=\small\itshape}
    ]

        % Application Layer
        \node[layer, fill=rustblue!20] (app) at (0, 4.5) {Application Code};

        % BSP Layer
        \node[layer, fill=green!20] (bsp) at (0, 3) {Board Support Package};

        % HAL Layer
        \node[layer, fill=yellow!20] (hal) at (0, 1.5) {Hardware Abstraction Layer};

        % PAC Layer
        \node[layer, fill=red!20] (pac) at (0, 0) {Peripheral Access Crate};

        % Hardware Layer
        \node[layer, fill=gray!20] (hw) at (0, -1.5) {Hardware Peripherals};

        % Arrows showing dependency
        \draw[arrow] (app) -- (bsp);
        \draw[arrow] (bsp) -- (hal);
        \draw[arrow] (hal) -- (pac);
        \draw[arrow] (pac) -- (hw);

        % Labels
        \node[label, left=2cm of app] {Board-Specific};
        \node[label, left=2cm of bsp] {Portable};
        \node[label, left=2cm of hal] {Register Access};
        \node[label, left=2cm of pac] {Hardware};

        % Portability indicators
        \node[font=\small, right=2cm of app] {Application-Specific};
        \node[font=\small, right=2cm of bsp] {Board-Specific};
        \node[font=\small, right=2cm of hal] {Chip Family Specific};
        \node[font=\small, right=2cm of pac] {Vendor Provided / Generated};
        \node[font=\small, right=2cm of hw] {Fixed by Hardware};
    \end{tikzpicture}
    \caption{The PAC/HAL/BSP layered architecture}
    \label{fig:pac-hal-bsp-layers}
\end{figure}

Each layer has distinct responsibilities and portability characteristics.

\subsection{Peripheral Access Crate (PAC)}
\label{subsec:pac-layer}

The PAC provides the lowest-level, register-accurate representation of the microcontroller's peripherals. It is typically generated automatically from the vendor's System View Description (SVD) file using \codeinline{svd2rust}.

\textbf{Key characteristics}:
\begin{itemize}
    \item One-to-one mapping with hardware registers
    \item No abstractions on top of hardware
    \item All operations are \codeinline{unsafe} (by design)
    \item Generated automatically—never hand-edited
    \item Chip variant specific (EFR32MG24B220 vs A020)
\end{itemize}

\begin{minted}{rust}
// Example: Direct register access via PAC
unsafe {
    (*efr32mg24_pac::GPIO_S::ptr())
        .port_a_dout
        .modify(|r, w| w.bits(r.bits() | (1 << 5)));
}
\end{minted}

This writes directly to GPIO Port A, bit 5. The \codeinline{unsafe} keyword forces you to acknowledge that you're bypassing Rust's safety guarantees.

\begin{warningbox}
    PAC code should be wrapped in HAL abstractions. Direct PAC usage in application code is an anti-pattern that sacrifices safety for no benefit.
\end{warningbox}

\subsection{Hardware Abstraction Layer (HAL)}
\label{subsec:hal-layer}

The HAL provides safe, ergonomic APIs that wrap the raw PAC access. It implements standard traits for portability and adds compile-time safety guarantees.

\textbf{Key characteristics}:
\begin{itemize}
    \item Safe APIs (no \codeinline{unsafe} required by users)
    \item Type-safe state machines (e.g., pin modes encoded in types)
    \item Implements \crate{embedded-hal} traits
    \item Chip family specific (works across all EFR32MG24 variants)
    \item Medium abstraction level
    \item Written and maintained by community
\end{itemize}

\begin{minted}{rust}
// Example: GPIO via HAL
let pin = gpio.pin5.into_output();
pin.set_high();
\end{minted}

Compare this to the PAC version above. The HAL provides:
\begin{itemize}
    \item Type-safe pin selection (can't use non-existent pins)
    \item Compile-time mode tracking (can't read from output pin)
    \item Clear, readable method names
    \item No unsafe required
    \item Identical performance (zero-cost abstraction)
\end{itemize}

\subsection{Board Support Package (BSP)}
\label{subsec:bsp-layer}

The BSP provides board-specific abstractions and convenience APIs. It's the only layer that "knows" about your specific PCB design.

\textbf{Key characteristics}:
\begin{itemize}
    \item Maps logical functions to physical pins (e.g., LED1 → PA5)
    \item Provides board-specific initialization sequences
    \item Combines peripherals for board features
    \item Board specific (each PCB design has its own BSP)
    \item Highest abstraction level
    \item Minimal code (thin wrappers over HAL)
\end{itemize}

\begin{minted}{rust}
// Example: BSP for XIAO MG24 Sense
use xiao_mg24_bsp::{Board, LED};

let board = Board::new();
board.led_green.on();
delay_ms(500);
board.led_green.off();
\end{minted}

The BSP knows that the green LED connects to PA5 on this particular board and provides a convenient LED abstraction.

\subsection{The Flow of Abstractions}
\label{sec:abstraction-flow}

Here's how a simple operation (turning on an LED) flows through the layers:

\begin{enumerate}
    \item \textbf{Application}: Calls \codeinline{board.led_red.on()}
    \item \textbf{BSP}: Translates to \codeinline{hal.gpio.pin5.set_high()}
    \item \textbf{HAL}: Checks pin mode, then calls PAC
    \item \textbf{PAC}: Executes \codeinline{GPIO->DOUT |= (1 << 5)}
    \item \textbf{Hardware}: Physical pin voltage goes high, LED illuminates
\end{enumerate}

This layered approach provides benefits:

\textbf{Portability}: Switch to a different EFR32MG24 board? Only rewrite the BSP. HAL and PAC remain unchanged.

\textbf{Maintainability}: Bug in GPIO handling? Fix it in the HAL. All boards using that HAL benefit.

\textbf{Composability}: Libraries written for HAL traits work across all boards using that HAL.

\begin{notebox}
    \textbf{Architecture Rule}: Application code should never directly use PAC. Always go through HAL and BSP.
\end{notebox}

\subsection{Why This Separation Matters}
\label{sec:separation-importance}

The embedded Rust community settled on this architecture through experience:

\subsection{PAC Regeneration}

When vendors release updated SVD files (fixing errors or documenting new features), you regenerate the PAC without affecting HAL or BSP code.

\subsection{Multiple HAL Implementations}

Different use cases may need different HAL trade-offs:
\begin{itemize}
    \item High-performance HAL for real-time applications
    \item Small-footprint HAL for resource-constrained devices
    \item Feature-rich HAL for development boards
\end{itemize}

All can use the same PAC.

\subsection{Ecosystem Compatibility}

The \crate{embedded-hal} trait system enables an ecosystem of drivers:
\begin{itemize}
    \item Display drivers (SSD1306, ST7789) work with any SPI implementation
    \item Sensor drivers (BMP280, MPU6050) work with any I2C implementation
    \item No code changes needed when switching microcontrollers
\end{itemize}

\subsection{Our Approach in This Book}
\label{sec:our-approach}

In this book, we'll implement the stack bottom-up:

\begin{enumerate}
    \item \textbf{Chapter 2}: Generate the PAC from Silicon Labs SVD
    \item \textbf{Chapters 4-9}: Build HAL for all peripherals
    \item \textbf{Chapter 11}: RAIL radio abstraction
    \item \textbf{Chapters 12-14}: Protocol stacks (Bluetooth, Zigbee, Thread)
    \item \textbf{Chapter 15}: BSP for XIAO MG24 Sense
\end{enumerate}

This order ensures:
\begin{itemize}
    \item Each layer has a strong foundation
    \item Clear dependency management
    \item Comprehensive testing at each level
    \item Real hardware validation
\end{itemize}

\begin{exercise}[Easy]
    Look at a simple development board (like Arduino Uno) and list:
    \begin{enumerate}
        \item Which features belong in PAC vs HAL vs BSP?
        \item Why does Arduino's "digitalWrite()" combine all three layers?
    \end{enumerate}
\end{exercise}

\begin{exercise}[Medium]
    Design a type-safe state machine for a USART that:
    \begin{enumerate}
        \item Only allows baud rate configuration when USART is disabled
        \item Only allows transmission when USART is enabled
        \item Prevents mode changes during active transmission
    \end{enumerate}
    How would you encode these states in Rust's type system?
\end{exercise}

\begin{exercise}[Hard]
    Research other embedded HAL implementations:
    \begin{enumerate}
        \item Compare \codeinline{stm32-hal} vs \codeinline{nrf-hal}
        \item Note differences in how they handle:
            \begin{itemize}
                \item Clock configuration
                \item Error types
                \item Feature flags
                \item Documentation approach
            \end{itemize}
    \end{enumerate}
\end{exercise}
\clearpage


% ============================================================================
% SECTION 01.04: Project Organization and Repository Structure
% ============================================================================
\section{Project Organization and Repository Structure}
\label{sec:project-organization}

% Section 01.04: Project Organization and Repository Structure

A well-organized repository structure is essential for maintainability, especially in embedded projects where you may support multiple chip variants, boards, and protocol stacks. In this section, we'll examine the repository structure used throughout this book and understand the rationale behind each decision.

\subsection{Top-Level Directory Structure}

Here's the complete structure of our EFR32MG24 project:

% Directory Tree for Section 1.4 - Repository Structure
\begin{figure}[H]
    \centering
    \begin{forest}
        for tree={
            font=\ttfamily\footnotesize,
            grow'=0,
            child anchor=west,
            parent anchor=south,
            anchor=west,
            calign=first,
            edge path={
                \noexpand\path [draw] (!u.south west) +(7.5pt,0) |- node[fill,inner sep=1.25pt] {} (.child anchor);
            },
            before typesetting nodes={
                if n=1
                    {insert before={[,phantom]}}
                    {}
            },
            fit=band,
            before computing xy={l=15pt},
        }
        [EFR32MG24/
            [efr32mg24-pac/
                [Cargo.toml]
                [src/]
                [svd/]
            ]
            [efr32mg24-hal/
                [Cargo.toml]
                [src/]
                [examples/]
            ]
            [efr32mg24-rail-sys/]
            [efr32mg24-bluetooth-sys/]
            [xiao-mg24-bsp/]
            [docs/]
            [book/
                [latex/]
                [TOC.yml]
            ]
            [examples/]
            [Cargo.toml]
            [README.md]
            [CLAUDE.md]
        ]
    \end{forest}
    \caption{EFR32MG24 project repository structure}
    \label{fig:repo-structure}
\end{figure}


\subsection{Crate Organization}

\subsection{Peripheral Access Crate (efr32mg24-pac)}

\begin{minted}{text}
efr32mg24-pac/
├── Cargo.toml          # PAC crate metadata
├── build.rs            # Build script (if needed)
├── svd/
│   └── EFR32MG24B220F1536IM48.svd  # SVD file
├── src/
│   └── lib.rs         # 138,000+ lines of generated code
└── README.md          # PAC-specific documentation
\end{minted}

**Key Points**:
\begin{itemize}
    \item Contains ONLY auto-generated code from svd2rust
    \item Never edited manually after generation
    \item Version locked to specific SVD file version
    \item Re-exported by HAL, users never depend on PAC directly
\end{itemize}

\subsection{Hardware Abstraction Layer (efr32mg24-hal)}

The HAL workspace contains multiple peripheral crates:

\begin{minted}{text}
efr32mg24-hal/
├── Cargo.toml              # Workspace root
├── src/
│   ├── lib.rs             # Main HAL entry point
│   ├── clock/             # Clock management (CMU)
│   │   ├── mod.rs
│   │   ├── types.rs
│   │   ├── traits.rs
│   │   └── README.md
│   ├── gpio/              # GPIO peripheral
│   │   ├── mod.rs
│   │   ├── types.rs
│   │   ├── traits.rs
│   │   └── README.md
│   ├── usart/             # USART communication
│   │   ├── mod.rs
│   │   ├── types.rs
│   │   ├── traits.rs
│   │   └── README.md
│   ├── i2c/               # I2C peripheral
│   │   ├── mod.rs
│   │   ├── types.rs
│   │   ├── traits.rs
│   │   └── README.md
│   ├── spi/               # SPI peripheral
│   │   ├── mod.rs
│   │   ├── types.rs
│   │   ├── traits.rs
│   │   └── README.md
│   └── prelude.rs         # Common imports
├── examples/              # HAL-level examples
│   ├── 01_clock.rs
│   ├── 02_gpio_blink.rs
│   ├── 03_usart_echo.rs
│   └── ...
└── docs/                  # HAL-specific documentation
    ├── README.md
    └── BUILDING.md
\end{minted}

\textbf{File Pattern}: Each peripheral follows the same structure with four files: mod.rs, types.rs, traits.rs, and README.md.

\begin{minted}{text}
peripheral-name/
├── mod.rs       // Main implementation
├── types.rs     // Type definitions
├── traits.rs    // Trait implementations
└── README.md    // Per-peripheral documentation
\end{minted}

\subsection{Radio System Crates (efr32mg24-rail-sys)}

FFI bindings to Silicon Labs RAIL (Radio Abstraction Interface Layer):

\begin{minted}{text}
efr32mg24-rail-sys/
├── Cargo.toml          // Link to vendor libraries
├── build.rs
├── wrapper.h          // bindgen header
└── src/
    └── lib.rs         # Generated FFI bindings
\end{minted}

**Purpose**: Enables multi-protocol radio operation (Bluetooth, Zigbee, Thread, proprietary).

\subsection{Protocol System Crates (bluetooth-sys, etc.)}

FFI bindings to higher-level protocol stacks:

\begin{minted}{text}
efr32mg24-bluetooth-sys/
├── Cargo.toml         # Bluetooth library linking
├── build.rs
├── wrapper.h
└── src/
    ├── lib.rs        # FFI bindings
    └── events/       # Event handling code
\end{minted}

\subsection{Board Support Package (xiao-mg24-bsp)}

Board-specific abstractions:

\begin{minted}{text}
xiao-mg24-bsp/
├── Cargo.toml          # Depends on efr32mg24-hal
├── src/
│   ├── lib.rs         # Board initialization
│   ├── board.rs       # Board struct
│   ├── sensors.rs     # Onboard sensors
│   └── radio.rs       # Radio utilities
└── examples/
    ├── blink_onboard_led.rs
    ├── ble_beacon.rs
    └── zigbee_coordinator.rs
\end{minted}

\subsection{Why This Structure Works}

\subsection{Clear Separation of Concerns}
\begin{itemize}
    \item PAC: Hardware truth (auto-generated)
    \item HAL: Safe abstractions (manually written)
    \item BSP: Board integration (thin wrappers)
    \item Application: User logic (no hardware details)
\end{itemize}

\subsection{Testing Strategy}
Each layer can be tested independently:
\begin{itemize}
    \item PAC: Verify register addresses match datasheet
    \item HAL: Unit tests with mocked registers
    \item BSP: Integration tests on real hardware
    \item Application: Standard Rust testing
\end{itemize}

\subsection{Version Management}
Crates can be versioned independently:
\begin{minted}{text}
efr32mg24-pac = "0.1.0"      # Hardware revision specific
efr32mg24-hal = "0.4.2"      # Feature updates
xiao-mg24-bsp = "0.2.1"      # Board revision specific
\end{minted}

\begin{exercise}[Easy]
    Sketch a project structure for supporting three boards:
    \begin{itemize}
        \item Seeed XIAO MG24 (covered here)
        \item Silicon Labs xG24-DK2601B dev kit
        \item Custom PCB with external sensors
    \end{itemize}
    What crates would be shared? Which would be board-specific?
\end{exercise}

\begin{exercise}[Medium]
    Consider a startup shipping products using this architecture:
    \begin{enumerate}
        \item How would you version the crates when releasing new firmware?
        \item What testing strategy ensures HAL changes don't break BSP?
        \item How would you handle multiple hardware revisions?
    \end{enumerate}
\end{exercise}

\begin{exercise}[Hard]
    Compare this structure to Linux kernel architecture:
    \begin{enumerate}
        \item Map PAC/HAL/BSP to Linux equivalents
        \item How does Linux handle different architectures? (arch/ directory)
        \item How does Linux handle different boards? (Device Tree)
        \item What lessons can we apply to our Rust embedded architecture?
    \end{enumerate}
\end{exercise}
\clearpage


% ============================================================================
% SECTION 01.05: Development Prerequisites and Toolchain Setup
% ============================================================================
\section{Development Prerequisites and Toolchain Setup}
\label{sec:development-prerequisites}

% Section 01.05: Development Prerequisites and Toolchain Setup

Before we can build and flash embedded Rust code to the EFR32MG24, we need to install and configure several tools. This section provides a complete setup guide for Linux, macOS, and Windows.

\subsection{Required Tools}
\label{sec:required-tools}

\textbf{Rust Toolchain}:

\begin{itemize}
    \item \textbf{Rust 1.70+} with support for the thumbv8m.main-none-eabihf target
    \item \textbf{cargo-binutils} for examining compiled binaries
    \item \textbf{cargo-size} for checking binary size
    \item \textbf{cargo-objdump} for disassembly
\end{itemize}
\subsection{Installation on Linux (Ubuntu/Debian)}
\label{subsec:install-linux}
\begin{minted}{bash}
# Install Rust via rustup
curl --proto '=https' --tlsv1.2 -sSf https://sh.rustup.rs | sh

source $HOME/.cargo/env

# Add ARM Cortex-M33 target
rustup target add thumbv8m.main-none-eabihf

# Install cargo tools
cargo install cargo-binutils
rustup component add llvm-tools-preview

# Install system dependencies
sudo apt-get update
sudo apt-get install -y \
    build-essential \
    pkg-config \
    libssl-dev \
    git \
    wget
\end{minted}
\subsection{Installation on macOS}
\label{subsec:install-macos}
\begin{minted}{bash}
# Install Rust via rustup
curl --proto '=https' --tlsv1.2 -sSf https://sh.rustup.rs | sh

source $HOME/.cargo/env

# Add ARM Cortex-M33 target
rustup target add thumbv8m.main-none-eabihf

# Install cargo tools
cargo install cargo-binutils
rustup component add llvm-tools-preview

# Install Homebrew if not already installed
/bin/bash -c "$(curl -fsSL https://raw.githubusercontent.com/Homebrew/install/HEAD/install.sh)"

# Install system dependencies
brew install git wget
\end{minted}
\subsection{Installation on Windows}
\label{subsec:install-windows}
\begin{enumerate}
    \item Download and install from https://rustup.rs/
    \item Run the installer with default settings
    \item After installation, open a new PowerShell/Command Prompt
    \item Install the ARM target:
\end{enumerate}

\begin{minted}{bash}
rustup target add thumbv8m.main-none-eabihf
cargo install cargo-binutils
rustup component add llvm-tools-preview
\end{minted}

\subsection{Hardware Debug Probe Setup}
\label{sec:debug-probe-setup}

The EFR32MG24 requires a debug probe for flashing and debugging. We'll use probe-rs, a Rust-based tool.

\subsection{Installing probe-rs}
\label{subsec:install-probe-rs}
\begin{minted}{bash}
# Install probe-rs
cargo install probe-rs --features cli

# Verify installation
probe-rs list
\end{minted}

Expected output (when hardware is connected):

\begin{minted}{text}
Probing OpenUSB... Done
[0]: J-Link Silicon Labs (<serial-number>) (J-Link)
\end{minted}

\subsection{Flashing: one sanctioned path}
\label{subsec:install-silabs-tools}

Silicon Labs' Simplicity Commander exists, and it is the right tool for one job: the
authorisation that unlocks a secured part. It is \emph{not} the flashing path this book
uses, because a bare \codeinline{commander flash} will happily write wherever the image
says — including over the bootloader.

What this project uses instead, and why it is worth the extra script:

\begin{minted}{bash}
# 1. show where the image would land, without touching the board
python3 scripts/elf-load-range.py target/thumbv8m.main-none-eabihf/release/examples/pac_imu
# load_min=0x08006000 load_max=0x08006813 size=2068 segments=3

# 2. dry run: prints the plan and refuses anything with load_min < 0x08006000
bash scripts/flash-lnx.sh "$ELF"

# 3. write, then verify by reading the flash back
bash scripts/flash-lnx.sh "$ELF" --confirm
\end{minted}

The board hangs off a separate host (\codeinline{lnx}) reachable over SSH, which is why
the script copies the ELF over and drives the programmer there. If your board is on the
machine you build on, the same script works locally — the guards, not the transport, are
the point.

\begin{quote}
\textbf{Measured caveat.} \codeinline{probe-rs} works for the first attach and then
desynchronises on this board's onboard CMSIS-DAP bridge
(\codeinline{Command ID in response (0x14) does not match sent command ID}); a USB
reset does not clear it and only a physical replug does. OpenOCD talks to the same probe
without the problem, so the script uses OpenOCD for reads and \codeinline{probe-rs} for
the write-and-verify step.
\end{quote}

\subsection{Updating Debug Probe Firmware}
\label{subsec:update-firmware}

If using a J-Link probe:

\begin{minted}{bash}
# On Linux/macOS
JLinkExe -CommanderScript <(echo "exit")

# Follow prompts to update firmware
\end{minted}

\subsection{Development Environment Setup}
\label{sec:dev-env-setup}

\subsection{VS Code Configuration}
\label{subsec:vscode-setup}

Install these extensions:
\begin{itemize}
    \item rust-analyzer (official Rust support)
    \item Even Better TOML (for Cargo.toml files)
    \item CodeLLDB (for debugging)
\end{itemize}

Create \filepath{.vscode/settings.json}:
\begin{minted}{json}
{
    "rust-analyzer.cargo.target": "thumbv8m.main-none-eabihf",
    "rust-analyzer.checkOnSave.allTargets": false,
    "rust-analyzer.cargo.features": ["rt"],
    "rust-analyzer.procMacro.ignored": {
        "cortex-m-rt": ["interrupt"]
    }
}
\end{minted}

\subsection{Cargo Configuration for Embedded}
\label{subsec:cargo-config}

Create \filepath{.cargo/config.toml}:
\begin{minted}{toml}
[target.thumbv8m.main-none-eabihf]
rustflags = [
    "-C", "link-arg=-Tlink.x",
    "-C", "link-arg=-Tdefmt.x",  # If using defmt
]

runner = "probe-rs run --chip EFR32MG24B220F1536IM48"

[build]
target = "thumbv8m.main-none-eabihf"

[env]
DEFMT_LOG = "debug"
\end{minted}

\subsection{Git Configuration}
\label{subsec:git-config}

Set up Git for embedded development:
\begin{minted}{bash}
git config --global user.name "Your Name"
git config --global user.email "your.email@example.com"

# Set up LFS for large files (like SVD files)
git lfs install
\end{minted}

Create \filepath{.gitattributes}:
\begin{minted}{text}
# Track large files with Git LFS
*.svd filter=lfs diff=lfs merge=lfs -text
*.pack filter=lfs diff=lfs merge=lfs -text
*.a filter=lfs diff=lfs merge=lfs -text
\end{minted}

\subsection{Verifying Your Setup}
\label{sec:verify-setup}

Run this verification script to ensure everything is correctly installed:
\begin{minted}{bash}
#!/bin/bash

echo "=== Checking Rust Installation ==="
rustc --version
cargo --version

echo ""
echo "=== Checking Targets ==="
rustup target list | grep "thumbv8m.main-none-eabihf (installed)"

echo ""
echo "=== Checking Cargo Tools ==="
cargo size --version 2>/dev/null && echo "✓ cargo-size installed"

echo ""
echo "=== Checking probe-rs ==="
probe-rs --version

echo ""
echo "=== Checking debug probe ==="
probe-rs list
\end{minted}

Expected output:
\begin{minted}{text}
=== Checking Rust Installation ===
rustc 1.70.0 (90c541806 2023-05-31)
cargo 1.70.0 (90c541806 2023-05-31)

=== Checking Targets ===
thumbv8m.main-none-eabihf (installed)

=== Checking Cargo Tools ===
✓ cargo-size installed

=== Checking probe-rs ===
probe-rs 0.18.0

=== Checking debug probe ===
[0]: J-Link Silicon Labs (xxxx) (J-Link)
\end{minted}

If all checks pass, you're ready to start development!

\begin{tipbox}
    \textbf{Common Issues}:
    \\• probe-rs list shows nothing: Check USB cable, ensure probe has power
    \\• Target not installed: Run rustup target add thumbv8m.main-none-eabihf
    \\• cargo-size not found: Run cargo install cargo-binutils
    \\• Permission denied (Linux): Add your user to dialout or plugdev group
\end{tipbox}

\begin{exercise}[Easy]
    Create a shell script that automates the entire setup verification process for your team.
\end{exercise}

\begin{exercise}[Medium]
    Write a cargo subcommand that checks for required tools and provides helpful installation instructions if anything is missing.
\end{exercise}

\begin{exercise}[Hard]
    Using Docker, create a container with the complete development environment pre-installed so new team members can start immediately.
\end{exercise}
\clearpage


% ============================================================================
% SECTION 01.06: Building Your First Embedded Binary
% ============================================================================
\section{Building Your First Embedded Binary}
\label{sec:building-first-binary}

% Section 01.06: Building Your First Embedded Binary

Now that we have our toolchain set up and understand the project structure, let's build and analyze a real embedded binary. This section will guide you through the complete build process with detailed explanations of each step.

\subsection{Preparing a Complete Example}
\label{sec:complete-example}

Let's create a minimal but complete embedded project that we can actually build and analyze:
\begin{minted}{bash}
$ cargo new --lib efr32mg24-minimal
$ cd efr32mg24-minimal
$ mkdir examples
$ cat > examples/minimal.rs << 'EOF'
#![no_std]
#![no_main]

use core::panic::PanicInfo;

#[panic_handler]
fn panic(_info: &PanicInfo) -> ! {
    loop {}
}

#[no_mangle]
pub extern "C" fn main() -> ! {
    loop {
        // Main application loop
        cortex_m::asm::nop();
    }
}
EOF
\end{minted}

Add dependencies to \filepath{Cargo.toml}:\

\begin{minted}{toml}
[package]
name = "efr32mg24-minimal"
version = "0.1.0"
edition = "2021"

[profile.dev]
opt-level = "z"
debug = true
overflow-checks = false  # Critical for embedded

[profile.release]
opt-level = "z"
debug = true
overflow-checks = false
panic = "abort"
lto = true

codegen-units = 1

[dependencies]
cortex-m = "0.7"
cortex-m-rt = "0.7"

# Development dependencies
[dev-dependencies]
cortex-m-semihosting = "0.5"
\end{minted}

\subsection{Reference: Complete Project Example}
\label{subsec:reference-complete-project}

For a comprehensive step-by-step guide with all project files (build.rs, Cargo.toml, memory.x, etc.), see Section \ref{sec:complete-project-example}.

This section focuses on understanding the build process and analyzing the resulting binary.

\subsection{Understanding the Build Process}
\label{sec:understanding-build}

When you run \codeinline{cargo build --release}, here's what happens:
\begin{remark}
    This diagram shows the build pipeline for embedded Rust:
\end{remark}

\begin{figure}[H]
    \centering
    \begin{tikzpicture}[
        node/.style={rectangle, draw=black, thick, minimum width=2.5cm, minimum height=1cm},
        arrow/.style={-Latex, thick}
    ]

        % Nodes
        \node[node, fill=rustblue!20] (source) at (0, 0) {Source Code};
        \node[node, fill=yellow!20] (rustc) at (3, 0) {rustc};
        \node[node, fill=orange!20] (llvm) at (6, 0) {LLVM};
        \node[node, fill=green!20] (obj) at (9, 0) {Object File};
        \node[node, fill=red!20] (linker) at (9, -2) {Linker};
        \node[node, fill=purple!20] (bin) at (6, -2) {Binary};

        % External files
        \node[font=\small] (linker-script) at (9, 0.8) {linker.ld};
        \node[font=\small] (startup) at (9, 1.2) {startup_s};

        % Arrows
        \draw[arrow] (source) -- (rustc);
        \draw[arrow] (rustc) -- (llvm);
        \draw[arrow] (llvm) -- (obj);
        \draw[arrow] (obj) -- (linker);
        \draw[arrow] (linker) -- (bin);
        \draw[arrow, dashed] (linker-script) -- (linker);
        \draw[arrow, dashed] (startup) -- (linker);

        % Labels
        \node[font=\small, below=0.2cm of source] {Human};
        \node[font=\small, below=0.2cm of rustc] {Targeting};
        \node[font=\small, below=0.2cm of llvm] {ARM};
        \node[font=\small, below=0.2cm of obj] {Cortex-M33};
        \node[font=\small, below=0.2cm of bin] {thumbv8m};
    \end{tikzpicture}
    \caption{Rust embedded build pipeline}
    \label{fig:build-pipeline}
\end{figure}

\subsection{Compile Step}
\begin{minted}{bash}
$ cargo build --release --example minimal
   Compiling autocfg v1.1.0
   Compiling proc-macro2 v1.0.56
   Compiling cortex-m v0.7.7
   Compiling quote v1.0.27
   Compiling cortex-m-rt v0.7.3
   Compiling efr32mg24-minimal v0.1.0
    Finished release [optimized] target(s) in 1.24s
\end{minted}

Notice that cargo automatically downloads and compiles dependencies. For embedded,
this includes startup code and linker scripts via \codeinline{cortex-m-rt}.

\subsection{Binary Analysis}
\label{sec:binary-analysis}

Once built, let's examine the binary:

\subsection{Size Analysis}
\begin{minted}{bash}
$ cargo size --release --example minimal
   text    data     bss     dec     hex filename
     80       0       0      80      50 minimal
\end{minted}

What do these sections mean?
\begin{itemize}
    \item \textbf{text}: The actual machine code (read-only)
    \item \textbf{data}: Initialized static variables
    \item \textbf{bss}: Uninitialized static variables (zeroed at startup)
    \item \textbf{dec}: Total bytes (decimal)
    \item \textbf{hex}: Total bytes (hexadecimal)
\end{itemize}

Our minimal program is 80 bytes total. Let's see what's inside:

\subsection{Disassembly}
\begin{minted}{bash}
$ cargo objdump --release --example minimal -- -d
\end{minted}

Output (simplified):
\begin{minted}{text}
00000000 <main>:
       0: b580            push    {r7, lr}
       2: f000 f802       bl      0xa <_ZN4core4...
       6: bf00            nop
       8: e7fe            b.n     0x8 <main+0x8>

0000000a <HardFault>:
       a: e7fe            b.n     0xa <HardFault>
\end{minted}

Interpretation:
\begin{itemize}
    \item \codeinline{push {r7, lr}}: Save registers
    \item \codeinline{bl}: Branch with link (call function)
    \item \codeinline{nop}: No operation (our loop body)
    \item \codeinline{b.n}: Branch to self (infinite loop)
    \item \codeinline{HardFault}: Default exception handler
\end{itemize}

\subsection{Listing Sections}
\begin{minted}{bash}
$ cargo objdump --release --example minimal -- -h
\end{minted}

Shows memory layout:
\begin{minted}{text}
Idx Name          Size      VMA       LMA       File off  Algn
  0 .text         00000014  00000000  00000000  00000034  2**2
                  CONTENTS, ALLOC, LOAD, READONLY, CODE
  1 .data         00000000  20000000  20000000  00000034  2**0
                  CONTENTS, ALLOC, LOAD, DATA
  2 .bss          00000000  20000000  20000000  00000034  2**0
                  ALLOC
\end{minted}

\textbf{Explanation}:
\begin{itemize}
    \item \codeinline{.text} starts at 0x00000000 (flash memory)
    \item \codeinline{.data/.bss} start at 0x20000000 (SRAM)
    \item Flash: 1,536 KB on EFR32MG24
    \item SRAM: 256 KB on EFR32MG24
\end{itemize}

\subsection{Build Optimization Comparison}
\label{sec:optimization-comparsion}

Build with different optimization levels:

\begin{table}[H]
    \centering
    \begin{tabular}{l c c c c}
        \toprule
        \textbf{Profile} & \textbf{text} & \textbf{data} & \textbf{bss} & \textbf{Total} \\
        \midrule
        dev (opt-level='z') & 124 & 0 & 0 & 124 \\
        release (opt-level='z') & 80 & 0 & 0 & 80 \\
        release (opt-level='s') & 76 & 0 & 0 & 76 \\
        release (opt-level=3) & 68 & 0 & 0 & 68 \\
        \bottomrule
    \end{tabular}
    \caption{Size comparison by optimization level}
    \label{tab:optimization-size}
\end{table}

Optimization levels:
\begin{itemize}
    \item \codeinline{'z'}: Optimize for size aggressively
    \item \codeinline{'s'}: Optimize for size
    \item \codeinline{'3'}: Optimize for speed (may reduce size too)
\end{itemize}

Even the unoptimized debug build is only 124 bytes!

\subsection{Linking and Linker Scripts}
\label{sec:linking}

The linker script is crucial in embedded systems. It tells the linker where to place code and data in memory.
\subsection{Memory Layout}
Create \filepath{memory.x}:
\begin{minted}{text}
MEMORY
{
  /* Flash memory for code and constants */
  FLASH (rx) : ORIGIN = 0x00000000, LENGTH = 1536K

  /* RAM for variables and stack */
  RAM (rwx)  : ORIGIN = 0x20000000, LENGTH = 256K
}
\end{minted}

This matches the EFR32MG24B220 memory map:
\begin{itemize}
    \item Flash: 1,536 KB starting at 0x0000_0000
    \item SRAM: 256 KB starting at 0x2000_0000
\end{itemize}
\subsection{Examining the Linker Output}
\begin{minted}{bash}
$ cargo build --release --example minimal --verbose 2>&1 | grep "LINK"
\end{minted}

Shows the linker command:
\begin{minted}{text}
"rust-lld" "-flavor" "gnu"
  "-Tlink.x"
  "-Tmemory.x"
  "-o" ".../target/release/examples/minimal"
  ...
\end{minted}

\textbf{Key flags}:
\begin{itemize}
    \item \codeinline{rust-lld}: LLVM linker for embedded targets
    \item \codeinline{-Tlink.x}: Standard Cortex-M linker script
    \item \codeinline{-Tmemory.x}: Our custom memory layout
\end{itemize}

\subsection{Creating a Flashable Binary}
\label{sec:flashable-binary}

The ELF binary contains debug symbols and metadata. For flashing, we need a raw binary:
\subsection{Convert to Binary Format}
\begin{minted}{bash}
$ cargo objcopy --release --example minimal -- -O binary minimal.bin
\end{minted}

Check sizes:
\begin{minted}{bash}
$ ls -lh target/thumbv8m.main-none-eabihf/release/
-rwxr-xr-x  80K minimal          # ELF with symbols
-rwxr-xr-x  80B minimal.bin     # Raw binary for flash
\end{minted}

The binary is much smaller because it contains only the raw machine code without debug symbols, section headers, etc.
\subsection{Hex Format for Some Tools}
Some tools prefer Intel hex format:
\begin{minted}{bash}
$ cargo objcopy --release --example minimal -- -O ihex minimal.hex
\end{minted}

\subsection{Automated Build Script}
\label{sec:automated-build}

Create a rust file, \filepath{build.rs}, in your project root:
\begin{minted}{rust}
use std::env;
use std::fs;
use std::path::PathBuf;

fn main() {
    // Tell cargo to rerun build script if memory.x changes
    println!("cargo:rerun-if-changed=memory.x");

    // Copy memory.x to output directory
    let out = &PathBuf::from(env::var_os("OUT_DIR").unwrap());
    fs::copy("memory.x", out.join("memory.x")).unwrap();

    println!("cargo:rustc-link-search={}", out.display());
}
\end{minted}

This ensures the linker script is available during linking.

\subsection{Build Automation with Justfile}
\label{sec:justfile}

For complex projects, use \structure{just} (a command runner):
\begin{minted}{text}
# Justfile - Build automation

# Build release binary
build:
    cargo build --release --example minimal

# Show binary size
size:
    cargo size --release --example minimal

# Disassemble
disasm:
    cargo objdump --release --example minimal -- -d

# Show section info
sections:
    cargo objdump --release --example minimal -- -h

# Convert to binary
binary:
    cargo objcopy --release --example minimal -- -O binary minimal.bin

# Run all analysis
analyze: build size disasm sections binary
    @echo "Analysis complete!"
\end{minted}

Usage:
\begin{minted}{bash}
$ just build
$ just size
$ just analyze
\end{minted}
\subsection{Key Takeaways}
\label{sec:ch1-key-takeaways}

\begin{keytakeaways}
    \item Embedded Rust binaries are incredibly small (our minimal example is 80 bytes)
    \item Optimization settings significantly impact size (from 124 to 68 bytes)
    \item The build process involves: compile → link → optionally convert format
    \item Linker scripts control memory layout and are critical for embedded
    \item Cargo provides powerful tools for analyzing binaries (size, objdump, objcopy)
    \item The ELF format contains metadata; for flashing, use binary or hex format
    \item Scripting your build workflow improves productivity
\end{keytakeaways}

\begin{exercise}[Easy]
    Create a project template with a complete build script that:
    \begin{itemize}
        \item Builds release binary
        \item Shows size information
        \item Creates both .bin and .hex files
        \item Runs objdump analysis
    \end{itemize}
\end{exercise}

\begin{exercise}[Medium]
    Compare binaries built with different optimization levels:
    \begin{itemize}
        \item What differences do you see in objdump output?
        \item How much code can the linker remove from our minimal example?
        \item Where does libc start inserting startup code?
    \end{itemize}
\end{exercise}

\begin{exercise}[Hard]
    Write a custom linker script that:
    \begin{itemize}
        \item Places the vector table at a specific address
        \item Creates a separate section for time-critical code
        \item Reserves space at the end of RAM for a stack canary
        \item Adds checksums for firmware verification
    \end{itemize}
    Then verify it works with objdump.
\end{exercise}


% ============================================================================
% SECTION 01.07: What Goes In Each Crate, and Why
% ============================================================================
\section{What Goes In Each Crate, and Why}
\label{sec:what-goes-in-each-crate}

The PAC/HAL/BSP pattern from Section~\ref{sec:pac-hal-bsp-architecture} says the layers exist.
It does not say where one ends and the next begins, and every project fills that gap
differently. That gap is what you face the first time you add support for a new MCU or a new
board — which is the situation this book is written for, and the one it should leave you
equipped to handle.

This section argues the choice, compares it with what four mature Rust embedded ecosystems do,
gives a decision procedure you can apply to hardware that is not in this book, and states the
rules that keep the answer from drifting.

\subsection{The Question, Stated Sharply}

Three pieces of code make the ambiguity visible. All three are real, from this project:

\begin{enumerate}
    \item \textbf{A pin pair.} The IMU on this board answers on \texttt{SDA = PB2},
          \texttt{SCL = PB3} — the opposite of what the vendor's pin map declares. Whose
          knowledge is that?
    \item \textbf{A power gate.} The same sensor does not answer until \texttt{PD5} is driven
          high. Is that a driver detail or a board detail?
    \item \textbf{An algorithm.} A fall classifier takes acceleration samples and decides
          whether the device fell. It has no registers, no pins and no clock. Where does it
          live so that it stays testable?
\end{enumerate}

A HAL that answers (1) and (2) knows too much: every new board becomes a fork of the HAL. An
example that answers them re-derives them, and the same investigation gets repeated by the
next person — the failure this project has already paid for once.

\subsection{Three Candidate Layouts}

\paragraph{Layout A — no BSP at all; wiring lives in the examples.}
This is the smallest possible structure: the HAL offers drivers, and each example says which
pins to use. It is what the STM32 and ESP32 ecosystems do, and it is defensible when the point
of the example is to teach the peripheral.

Its cost appears the moment the wiring is not obvious. This board's IMU needed a swapped
\texttt{SDA}/\texttt{SCL} and a power gate before it answered anything; putting that in every
example means every example can get it wrong independently, and the knowledge is only as good
as the most recently written one.

\paragraph{Layout B — one BSP crate, features for the optional hardware.}
A single \codeinline{xiao-}\allowbreak\codeinline{mg24-bsp} with features
\codeinline{sense} and \codeinline{expansion-base}. It looks like the cheapest correct answer, and it has a specific
trap that the nRF ecosystem documents in its own repository: \emph{cargo unifies features
across a workspace}, so a build that enables the base for one target can silently enable it
for another. More importantly, the API lies: a \codeinline{ExpansionBase} type that exists in a
build without the base has to answer ``no OLED'' at runtime, and every method that could have
been a compile error becomes an \codeinline{Option}.

\paragraph{Layout C — one BSP crate per board, and a feature for a variant of one board.}
The module gets its own crate (LED on both variants; IMU and microphone behind
\codeinline{sense}), the Expansion Base gets its own crate that depends on the module's, and
the dependency arrow points one way. A user who does not own the base never compiles it, and
never sees an \codeinline{Option} that only exists because two boards were welded together.

\subsection{What Four Ecosystems Actually Do}

Read from the repositories, not from memory:

\begin{table}[H]
\centering
\footnotesize
\begin{tabular}{@{}>{\raggedright\arraybackslash}p{2.3cm}>{\raggedright\arraybackslash}p{2.6cm}>{\raggedright\arraybackslash}p{2.4cm}>{\raggedright\arraybackslash}p{3.4cm}>{\raggedright\arraybackslash}p{3.0cm}@{}}
\hline
\textbf{ecosystem} & \textbf{chip selection} & \textbf{BSP crates} & \textbf{where examples live} & \textbf{where wiring is written} \\
\hline
\textbf{rp-hal} (RP2040/RP2350) & one HAL per chip & \textbf{yes}, one per board,
  in a \emph{separate repository} (\codeinline{rp-hal-boards/boards/}) & in each BSP, plus
  separate example crates & in the BSP, once \\
\textbf{nrf-hal} & one HAL per chip (\codeinline{nrf52840-hal}, \codeinline{nrf9160-hal}, \dots)
  + \codeinline{nrf-hal-common} & no & in the HAL repo, board chosen by feature & in the
  examples \\
\textbf{stm32f4xx-hal} & one HAL per family & no & in the HAL crate &
  in the examples (pin numbers inline) \\
\textbf{esp-hal} & one crate, chip by feature & no & one \codeinline{examples/} grouped by
  capability & in the examples \\
\hline
\end{tabular}
\caption{How four Rust embedded ecosystems split PAC, HAL, BSP and examples. Verified from the
repositories in November 2026.}
\label{tab:ecosystem-layouts}
\end{table}

Two details in that table are worth more than the table.

First, \textbf{rp-hal puts its BSPs in a different repository from the HAL}. That is stronger
than a folder: the HAL's continuous integration cannot depend on a board, because the board is
not present. The BSP crates describe themselves as a HAL for one board that re-exports the
chip HAL — a facade, exactly as this book calls it.

Second, \textbf{nrf-hal ships a workspace template instead of a workspace}, and says why:
\emph{``Since Cargo resolves all features for all crates in the workspace together, it is not
possible to have a workspace with all the crates together, since they enable incompatible
features.''} That sentence is the argument against Layout B in the ecosystem's own words.

\subsection{The Decision, and the Rule Behind Each Part}

\begin{table}[H]
\centering
\footnotesize
\begin{tabular}{@{}>{\raggedright\arraybackslash}p{3.4cm}>{\raggedright\arraybackslash}p{5.2cm}>{\raggedright\arraybackslash}p{5.2cm}@{}}
\hline
\textbf{crate / folder} & \textbf{what goes in} & \textbf{what must never go in} \\
\hline
\codeinline{efr32mg24-pac} & registers, generated from the SVD & anything hand-written that
  a generator could produce \\
\codeinline{efr32mg24-hal} & drivers: which bytes, in which order, with which timing & pins,
  addresses, power gates, board names \\
\codeinline{xiao-mg24-bsp} & the module: LED; behind \codeinline{sense}, the IMU bus, its
  power gate and the microphone & any device that is not on the module \\
\codeinline{xiao-expansion-base-bsp} & the base: OLED, BME68x, RTC, buzzer, button, Grove,
  microSD, and the base's own bus & anything that exists without the base \\
\codeinline{libraries/} & domain policy with no hardware at all & any dependency on the
  hardware stack, direct or transitive \\
\hline
\end{tabular}
\caption{Ownership. The right-hand column is what the gate checks; the left two are what a
review checks.}
\label{tab:crate-ownership}
\end{table}

Three rules follow, and each one exists because of a specific failure:

\begin{keytakeaways}
    \item \textbf{A BSP describes one piece of plastic.} Two boards that stack are two
          crates, because a user may own only one of them.
    \item \textbf{A variant of one board is a feature, not a crate.} The Sense module is the
          same PCB with more parts fitted; a separate crate would duplicate the board init
          and the LED to add two sensors.
    \item \textbf{A dependency arrow never points at a board.} The base depends on the
          module; the module does not know the base exists. This is the property that makes
          ``does it work without the base?'' answerable by the compiler instead of by a test
          matrix.
\end{keytakeaways}

\subsection{Where the Algorithm Goes}

The classifier deserves its own sentence because it is where layering usually breaks. It has
no hardware, but the temptation is to put it ``near'' the driver that produces its input — in
the same crate as the IMU, or in the example that uses it.

Both are wrong for the same reason: the moment it shares a crate with anything that needs a
target triple, its tests need a target triple, and the tests are what keep it honest. In this
project it is a crate with \textbf{zero dependencies}, and the gate fails if its dependency
tree ever contains a hardware crate:

\begin{minted}{bash}
$ cargo tree -p motion-detection | grep efr32mg24
$ echo $?
1        # nothing from the hardware stack: the rule holds
\end{minted}

The contract is worth quoting, because it is the whole testability argument in one line:

\begin{minted}{rust}
// [i32; 3] of milligravity, and the interval between samples in microseconds.
pub fn update_us(&mut self, accel_mg: [i32; 3], dt_us: u32) -> Option<Event>
\end{minted}

No sensor type, no bus, no clock. Which is why nine tests run in under a second on a laptop,
and why the same code can be replayed against a recording at ten times real speed.

\subsection{And Where the Examples Go}

The layout of the examples follows the same principle as the crates, with one addition for
hardware that stacks: \textbf{an example lives in the crate of the largest hardware
requirement it has, and states that requirement in its first comment.}

\begin{minted}{text}
crates/xiao-mg24-bsp/examples/            runs with the module alone on the desk
crates/xiao-expansion-base-bsp/examples/  needs the module (Sense) and the base
crates/efr32mg24-pac/examples/            chip level: any pin is legitimate
\end{minted}

The alternative — every example in one folder — looks tidier and is worse: the crate of the
module would have no examples, and an API that nothing builds is an API that nothing checks.
A crate with hardware features and zero examples should fail a gate, not pass it quietly.

\subsection{A Decision Procedure for Hardware Not In This Book}

The four questions below decide the layout. They take minutes and they are the difference
between a structure that survives the second board and one that has to be refactored.

\begin{enumerate}
    \item \textbf{How many chips?} If more than one, the HAL is per chip (or per family, when
          the parts are close enough that one driver set serves them). This is what
          \codeinline{nrf-hal} does per chip and \codeinline{stm32f4xx-hal} does per family;
          \codeinline{esp-hal} instead keeps one crate and selects the chip by feature, which
          works when the vendor's crates already share their register code.
    \item \textbf{How many boards, and can one sit on top of another?} Boards that are
          alternatives become one BSP crate each. Boards that \emph{stack} become one crate
          each \emph{plus} a one-way dependency: the carrier depends on the module. A single
          crate that knows both cannot answer ``does this work without the carrier?''.
    \item \textbf{Are there variants of the same board?} A variant — same PCB, different
          parts fitted — is a \emph{feature}, not a crate. A separate crate duplicates the
          board init and everything the variants share, in exchange for a name that repeats
          what the feature flag already says.
    \item \textbf{Is the wiring obvious or hard-won?} Obvious wiring (an LED on a documented
          pin) may live in an example. Hard-won wiring — a swapped bus, a power gate, a pin
          that the vendor's own map gets wrong — belongs in a BSP, once, with the measurement
          recorded next to it. This is the question that decides between Layout A and
          Layout C for most hobby-scale projects.
\end{enumerate}

The same questions decide where an algorithm goes, with one addition: \textbf{if it can be
tested on a host, it must be}, so it cannot share a crate with anything that needs a target
triple.

\subsection{Adding Support for a New Chip or Board: the Order of Work}

Having decided the layout, the order below is what this project would follow again. Each step
produces something checkable before the next one starts.

\begin{enumerate}
    \item \textbf{Registers first.} Obtain the SVD, generate the PAC, and build one example
          that only writes a register and reads it back. If the SVD is wrong, everything
          downstream inherits the error, and a five-line example finds that in minutes.
    \item \textbf{Clock and pins.} Prove that the peripheral clock gate is on and that a GPIO
          changes state. This is the step where a missing
          \codeinline{CMU\_CLKEN0} bit appears as a bus fault at the first register access.
    \item \textbf{One transport.} Bring up the simplest bus the board has, and read a device
          identifier through it. A chip identifier or a WHO\_AM\_I is a free proof that
          addressing, timing and the byte order are all correct.
    \item \textbf{Drivers, in the HAL, with pins as arguments.} Nothing in this layer may name
          a port. If a driver needs a pin, take it as a parameter; the BSP will supply it.
    \item \textbf{The BSP for one board, and its examples.} Put the wiring in one place, in
          the crate named after the board whose documentation describes the feature, and
          write the example that proves the board works end to end.
    \item \textbf{A second board, if there is one, and the dependency direction.} Add the
          carrier crate only after the module crate stands alone, and make the dependency
          point one way. If it cannot, the split is wrong.
    \item \textbf{The gate.} Before more examples are written, make the machine check what
          the review would check: zero warnings, every example builds, every domain crate
          tests on the host, and no domain crate depends on the hardware stack.
\end{enumerate}

\begin{warningbox}
    \textbf{The step people skip is the first one.} Generating the PAC and immediately
    writing drivers feels productive, and it hides an SVD error until the driver is
    complicated enough that the error looks like a driver bug. The five-line register
    read-back is the cheapest test in the whole project.
\end{warningbox}

\subsection{What This Costs}

Honesty about a design decision includes its price:

\begin{itemize}
    \item \textbf{More manifests.} Two BSP crates mean two \filepath{Cargo.toml} files, two
          example folders and a dependency between them that has to be kept one-way.
    \item \textbf{Some knowledge is duplicated on purpose.} The clock setup lives in the
          module crate and the base depends on the module rather than re-initialising; a
          project with two unrelated boards could not share it this way.
    \item \textbf{The split is a judgement, not a theorem.} ``Which plastic is this pin on?''
          is unambiguous for an LED and less so for a Grove connector that reaches the base
          through the module's castellated pads. The rule that resolves it: the BSP is named
          after the board whose \emph{documentation} describes the feature.
\end{itemize}

\begin{notebox}
    \textbf{The test to apply to your own board.} Ask which of your examples would still make
    sense to someone holding only the chip. If the answer is ``none of them'', the board has
    swallowed the crate and a split is due. If the answer is ``all of them'', you may not need
    a BSP at all — Layout A is not a mistake, it is a choice with a cost that is sometimes
    worth paying.
\end{notebox}


\clearpage

% ============================================================================
% Chapter Summary and Key Takeaways
% ============================================================================

\begin{keytakeaways}
    \item Rust provides memory safety without sacrificing performance through its ownership system
    \item Embedded Rust programs use \codeinline{no_std} and \codeinline{no_main} for bare-metal operation
    \item The PAC/HAL/BSP architecture provides clean separation of concerns and promotes code reuse
    \item Cross-compilation requires the correct target specification (thumbv8m.main-none-eabihf for EFR32MG24)
    \item Building embedded binaries involves multiple steps: compile, link, and optionally convert to binary/hex format
    \item All necessary tools (rustup, probe-rs, cargo-binutils) are available cross-platform
    \item Our minimal program compiles to just 80 bytes in release mode
\end{keytakeaways}



% Chapter 2: SVD Files and PAC Generation (written)
% Chapter 2: Understanding SVD Files and the PAC Generation Process
% File: chapters/chapter-02-pac-generation/chapter-02-pac-generation.tex

\chapter{Understanding SVD Files and the PAC Generation Process}
\label{ch:chapter-02}
\chapthumb{2}

\begin{quote}
    \textit{"The PAC is the foundation of all embedded Rust development---get it right, and everything else follows naturally."}
    \\
    -- Embedded Rust Working Group, 2024
\end{quote}

This chapter dives deep into System View Description (SVD) files and the process of generating a Peripheral Access Crate (PAC). You'll learn how to obtain official SVD files from Silicon Labs, use \codeinline{svd2rust} to generate type-safe register access code, and verify the resulting 138,000+ lines of generated Rust code.

\minitoc

\clearpage

% ============================================================================
% SECTION 02.01: What are SVD Files?
% ============================================================================
\input{chapters/chapter-02-pac-generation/02.01-what-are-svd-files}

% ============================================================================
% SECTION 02.02: Obtaining SVD Files from Silicon Labs
% ============================================================================
\input{chapters/chapter-02-pac-generation/02.02-obtaining-svd-files}

% ============================================================================
% SECTION 02.03: svd2rust: The PAC Generation Tool
% ============================================================================
\input{chapters/chapter-02-pac-generation/02.03-svd2rust-tool}

% ============================================================================
% SECTION 02.04: Analyzing Generated PAC Output
% ============================================================================
\input{chapters/chapter-02-pac-generation/02.04-analyzing-pac-output}

% ============================================================================
% SECTION 02.04a: Progressive Example Part 2 - PAC-Based Blink
% ============================================================================
\input{chapters/chapter-02-pac-generation/02.04a-pac-blink-example}

% ============================================================================
% SECTION 02.05: Building and Verifying the PAC Crate
% ============================================================================
\input{chapters/chapter-02-pac-generation/02.05-building-verifying-pac}

% ============================================================================
% SECTION 02.06: PAC Variants and Device Selection
% ============================================================================
\input{chapters/chapter-02-pac-generation/02.06-pac-variants}

\clearpage

% ============================================================================
% Chapter Summary and Key Takeaways
% ============================================================================

\begin{keytakeaways}
    \item SVD files are XML-based hardware descriptions provided by chip vendors
    \item \codeinline{svd2rust} generates type-safe Rust code from SVD files automatically
    \item The \codeinline{--locked} flag is critical when installing \codeinline{svd2rust} to avoid dependency conflicts
    \item Generated PACs contain 138,000+ lines of register access code for the EFR32MG24
    \item PAC code should never be edited manually---always regenerate from SVD
    \item Device variant selection (A020 vs B220) must match your actual hardware
    \item The PAC provides zero-cost type-safe abstractions over hardware registers
    \item Build verification ensures PAC correctness before using in HAL development
\end{keytakeaways}


% Chapter 3: Build Systems (written)
% Chapter 3: Build Systems
% File: chapters/chapter-03-build/chapter-03-build.tex

\chapter{Build Systems and Cross-Compilation Setup}
\label{ch:chapter-03}
\chapthumb{3}

\begin{quote}
    \textit{``A build that succeeds is not evidence that it built what you wrote.''}
\end{quote}

Cross-compiling for a microcontroller means four things have to agree: the target triple, the memory
layout, the linker's idea of what a loadable segment is, and the image the programmer writes. When
they disagree, the build still succeeds — that is the nature of linker scripts — and the board either
does nothing or, in the worst case, no longer boots.

This chapter is the whole chain for this board, with the two build defects that produced a green log
and a wrong artefact.

\minitoc

\clearpage

\section{The Target Triple, Decoded}
\label{sec:build-target}

\codeinline{thumbv8m.main-none-eabihf} is four decisions in one string, and three of them can be
wrong in ways that still compile.

\begin{table}[htbp]
    \centering
    \footnotesize
    \begin{tabular}{@{}p{0.22\textwidth}p{0.30\textwidth}p{0.42\textwidth}@{}}
        \toprule
        \textbf{Part} & \textbf{Means} & \textbf{If you get it wrong} \\
        \midrule
        \codeinline{thumbv8m} & ARMv8-M | ARMv8-M Baseline & \codeinline{armv7m} links, then traps on an instruction the core has \\
        \codeinline{main} & mainline, not baseline & baseline drops the DSP and FPU instructions you may be using \\
        \codeinline{none} & no OS, no libc & \codeinline{std} becomes available and fails deep in linking \\
        \codeinline{eabihf} & hard-float ABI & \codeinline{eabi} links and quietly passes floats in the wrong registers \\
        \bottomrule
    \end{tabular}
    \caption{The four decisions. Only the last one fails \emph{quietly}.}
    \label{tab:target-triple}
\end{table}

\subsection{Hard Float, and Why It Matters Here}

The EFR32MG24 has a single-precision FPU (FPv5-SP). \codeinline{eabihf} tells the compiler it may
pass floating-point arguments in \codeinline{s0}--\codeinline{s15} and may emit FPU instructions
directly. With \codeinline{eabi} the same source compiles to calls into soft-float helpers, links
perfectly, and runs several times slower on exactly the arithmetic this project's classifier uses.

There is one trap on the way in: \codeinline{-mfpu} and \codeinline{-mfloat-abi} flags that
\emph{conflict} with the target's own choice produce a warning on every build — ``switch
\codeinline{-mcpu=cortex-m33} conflicts with switch \codeinline{-march=armv8-m.main+fp}''. That
warning is not harmless: it appeared once per C file in this project and buried the real warnings
under it. The flags belong to the target; a build script adding its own is duplicating a decision it
does not own.

\subsection{The Toolchain, in Files}

\begin{minted}{bash}
# the target must be installed once
rustup target add thumbv8m.main-none-eabihf

# .cargo/config.toml decides the default target and the runner
[build]
target = "thumbv8m.main-none-eabihf"

[target.thumbv8m.main-none-eabihf]
rustflags = ["-C", "link-arg=-Tlink.x"]
\end{minted}

\codeinline{link.x} comes from \codeinline{cortex-m-rt}, and it is the file that reads
\codeinline{memory.x} and lays the image out. That indirection is why a wrong \codeinline{memory.x}
produces no error: \codeinline{link.x} simply places the sections where it was told to.

\subsection{Per-Crate Features, Not Per-Crate Flags}

A workspace with several crates that must not duplicate the vector table uses a feature to decide
which \emph{one} of them pulls \codeinline{cortex-m-rt}:

\begin{minted}{toml}
# the HAL exposes the runtime but does not enable it by default
[features]
rt = ["dep:cortex-m-rt"]

# examples turn it on; there is exactly one #[entry] per binary
efr32mg24-hal = { path = "../efr32mg24-hal", features = ["rt"] }
\end{minted}

Without that split, two crates in one binary each define \codeinline{Reset} and the link fails with a
duplicate symbol that names neither crate. With it, the rule is explicit: a library may know about
the runtime; only a final binary may install it.

\section{The Bootloader Decides Your Memory Layout}
\label{sec:build-bootloader}

The single most dangerous line in an embedded project is the \codeinline{ORIGIN} of the flash region,
and it is usually copied from a tutorial for a board with no bootloader.

\subsection{What Is at the Base of Flash}

This board ships with a Gecko Bootloader occupying the first 24 KiB. On the EFR32MG24 flash begins at
\codeinline{0x08000000}, so:

\begin{table}[htbp]
    \centering
    \footnotesize
    \begin{tabular}{@{}p{0.26\textwidth}p{0.26\textwidth}p{0.42\textwidth}@{}}
        \toprule
        \textbf{Region} & \textbf{Addresses} & \textbf{Owner} \\
        \midrule
        Bootloader & \codeinline{0x08000000}--\codeinline{0x08005FFF} & factory; never written by this project \\
        Application & \codeinline{0x08006000}--end & every image in this book \\
        \bottomrule
    \end{tabular}
    \caption{24 KiB that a linker script must not touch.}
    \label{tab:flash-map}
\end{table}

\subsection{The Wrong memory.x, and What It Does}

The version that appeared in this project's first draft of the book — copied from a generic
\codeinline{cortex-m} example — read:

\begin{minted}{text}
MEMORY
{
  FLASH (rx) : ORIGIN = 0x00000000, LENGTH = 1536K
  RAM  (rwx) : ORIGIN = 0x20000000, LENGTH = 256K
}
\end{minted}

Every line is wrong for this board, and none of them produces an error. The linker happily lays the
image out from address zero, the programmer writes it exactly there, the write verifies — and the
board no longer boots, because the bootloader that was at zero has been replaced by an application
that expects to be at \codeinline{0x08006000}.

\subsection{The Version in Use}

\begin{minted}{text}
MEMORY
{
  /* The Gecko bootloader owns 0x08000000..0x08005FFF. An application that starts
     below 0x08006000 overwrites it. */
  FLASH (rx)  : ORIGIN = 0x08006000, LENGTH = 1536K - 24K
  /* The mailbox at 0x2003F000 is a convention shared by every example; the stack
     is placed below it so a deep call cannot scribble over the report. */
  RAM   (rwx) : ORIGIN = 0x20000000, LENGTH = 252K
}

_stack_start = ORIGIN(RAM) + LENGTH(RAM) - 4096;
\end{minted}

Both numbers are deliberate: 252K rather than 256K because the top 4 KiB are reserved, and
\codeinline{\_stack\_start} is below the mailbox so a stack overflow corrupts a report rather than the
data being reported.

\subsection{The Rule, Enforced in a Script}

A comment is not protection. The project has two mechanical guards:

\begin{minted}{bash}
# 1. where would this image land?
python3 scripts/elf-load-range.py target/thumbv8m.main-none-eabihf/release/examples/pac_led_blink
# load_min=0x08006000 load_max=0x08006523 size=1316 segments=3

# 2. refuse anything below the bootloader boundary, before writing
bash scripts/flash.sh "$ELF"          # dry run: prints the plan, refuses load_min < 0x08006000
bash scripts/flash.sh "$ELF" --confirm
\end{minted}

The first command is the one worth internalising: \textbf{read the load addresses before every flash},
not after a board stops answering. It takes one second and it is the difference between a mistake and
a brick.

\section{--nmagic, or the Linker's Idea of a Segment}
\label{sec:build-nmagic}

This section documents the least intuitive flag in the whole chain, because without it the image is
built, verified, and written to the wrong place.

\subsection{What the Linker Does by Default}

ELF loadable segments carry two addresses: a virtual address (\codeinline{VMA}, where the code expects
to run) and a physical address (\codeinline{LMA}, where the bytes must land in memory). For flash that
is always 1:1 — you cannot relocate code that runs from flash.

A linker aiming at an operating system still aligns segments to page boundaries, and to do that it
gives the segment an \emph{offset} that no longer matches its load address. An OS loader compensates;
a bare-metal programmer does not, and writes the file's bytes at the segment's file offset while the
code expects them at the segment's address.

\subsection{The Symptom, Measured}

\begin{minted}{text}
# with a default link: the loadable segment's file offset is 0x1000, its address is 0x08006000
probe-rs download image.elf        # writes bytes at the wrong offset
probe-rs verify image.elf          # verifies the same wrong offsets: it agrees with itself
                                   # the board does not run
\end{minted}

This is the failure that teaches the difference between ``verified'' and ``correct''. The programmer
compares the flash against the file it just wrote. If the file was written at the wrong offset by a
consistent rule, both sides agree and the check passes.

\subsection{The Fix, and Why It Is One Flag}

\begin{minted}{text}
# in .cargo/config.toml, or as a link argument
rustflags = ["-C", "link-arg=--nmagic"]
\end{minted}

\codeinline{--nmagic} tells the linker not to align segments and not to pad them — the image is flat,
VMA equals LMA for every loadable segment, and the programmer's naive write is correct. Every bare
metal Rust project needs it and most tutorials omit it, because their toolchain's runner happens to
compensate.

\subsection{Checking the Invariant Instead of Trusting the Flag}

The project's inspection script prints exactly what matters, so the flag is verified rather than
believed:

\begin{minted}{bash}
python3 scripts/elf-load-range.py "$ELF"
# for each PT_LOAD: paddr == vaddr, and load_min >= 0x08006000
#   seg paddr=0x08006158 vaddr=0x08006158 filesz=5400
#   seg paddr=0x08007670 vaddr=0x08007670 filesz=108
#   seg paddr=0x080076DC vaddr=0x20000000 filesz=4     <- the .data initialiser
\end{minted}

Three lines, and two facts read directly off them: every code segment has \codeinline{paddr == vaddr}
(the flag worked), and \codeinline{load\_min} is above the bootloader (the memory map is right). The
third segment is the initialiser for \codeinline{.data}, which correctly has a flash address and a RAM
address — that one is supposed to differ.

\section{Builds That Lie}
\label{sec:builds-that-lie}

Two defects in this project had the same shape, and it is the most dangerous shape a build system can
have: \textbf{the build succeeded, and the artefact was not what the sources said.}

\subsection{Case 1: A build.rs That Did Not Watch Its Own C File}

A C file compiled by the \codeinline{cc} crate, used to provide symbols the vendor blobs reference:

\begin{minted}{rust}
// build.rs — the version that lied
cc::Build::new().file("src/cmsis_stubs.c").compile("cmsis_stubs");
println!("cargo:rerun-if-changed=blobs/");
println!("cargo:rerun-if-changed=vendor/");
\end{minted}

Editing \codeinline{src/cmsis_stubs.c} changed nothing. Cargo had no reason to re-run the script, the
old object file stayed, and the link failed on a symbol that was visibly present in the source — with
a message that sends you looking at the blob instead of the build.

\begin{minted}{rust}
// the version that tells the truth
println!("cargo:rerun-if-changed=src/cmsis_stubs.c");
\end{minted}

The general rule: \textbf{any file a build script reads or compiles is an input, and an input it does
not declare is an input the build ignores.} Cargo does not scan your script for file paths.

\subsection{Case 2: A Book That Compiled Yesterday's Sources}

This book is built in a container on a different machine, from a tarball:

\begin{minted}{bash}
# the version that lied: extracts to efr32book/book/latex/
tar cf - book/latex | ssh host 'cd ~/efr32book && tar xf -'

# ...while the container compiles efr32book/latex/, so every build used the previous copy
podman exec --workdir /workspace/efr32book/latex tex xelatex main.tex
\end{minted}

The path \codeinline{book/latex} is preserved by \codeinline{tar}, the compile directory is not, and
the two never met. A whole section was written, ``built'' successfully, and did not appear in the PDF.

The fix is one flag plus one check:

\begin{minted}{bash}
tar cf - book/latex | ssh host 'cd ~/efr32book && tar xf - --strip-components=1'

# and then prove it, before compiling
md5sum book/latex/chapters/.../01.08-layer-boundaries.tex
ssh host 'md5sum ~/efr32book/latex/chapters/.../01.08-layer-boundaries.tex'   # must match
\end{minted}

\subsection{The Pattern Behind Both}

\begin{table}[htbp]
    \centering
    \footnotesize
    \begin{tabular}{@{}p{0.30\textwidth}p{0.28\textwidth}p{0.34\textwidth}@{}}
        \toprule
        \textbf{Lie} & \textbf{What you see} & \textbf{What catches it} \\
        \midrule
        unwatched input & link error naming a symbol you wrote & \codeinline{rerun-if-changed} on every input \\
        wrong extraction path & success, no new content & compare a hash before compiling \\
        cached artefact & old behaviour at new code & clean the package, or watch the log for the crate \\
        \bottomrule
    \end{tabular}
    \caption{Three ways to build the wrong thing, and the one-line check for each.}
    \label{tab:build-lies}
\end{table}

The common answer is the same in all three cases: \textbf{verify that the thing being built is the
thing you wrote.} It is a one-line check, and it is cheaper than any of the afternoons it saves.

\section{The Gate: What Runs Without a Board}
\label{sec:build-gate}

A project with hardware has a scheduling problem: the compiler is available constantly, the board is
available sometimes. Anything that needs the board gets deferred, and deferred checks do not happen.
The answer this project uses is a single script that runs everything that does not need hardware.

\subsection{What It Checks}

\begin{minted}{text}
$ bash scripts/gate.sh
PASS  cargo fmt --all --check
PASS  efr32mg24-hal clippy, 0 warnings
PASS  efr32mg24-pac clippy, 0 warnings
PASS  xiao-expansion-base-bsp clippy, 0 warnings
PASS  xiao-mg24-bsp clippy, 0 warnings
PASS  motion-detection clippy, 0 warnings
PASS  efr32mg24-hal: 5 examples of its own level
PASS  efr32mg24-pac: 19 examples of its own level
PASS  xiao-expansion-base-bsp: 9 examples of its own level
PASS  motion-detection: host tests (9 passed)
PASS  motion-detection: no hardware dependency
PASS  examples outside the PAC crate name no PAC type
PASS  no BSP re-exports the PAC
PASS  the HAL exposes no register block publicly
PASS  no example steals a core peripheral
VERDICT: PASS
\end{minted}

Every line is a property, not a command. That distinction matters: a script that runs
\codeinline{cargo clippy} tells you a tool ran; a line that says \emph{zero warnings in five crates}
tells you what is true.

\subsection{The Discovered-vs-Listed Rule}

The first version of this gate had a hard-coded list of crates, and the crate that most needed
checking was the one that was not on it. The version in use \emph{discovers} crates:

\begin{minted}{bash}
for dir in crates/*/ libraries/*/; do
  [ -f "${dir}Cargo.toml" ] || continue
  name=$(basename "$dir")
  ...
done
\end{minted}

A list is a snapshot of what existed when someone wrote it. Discovery cannot fall behind — which is
exactly what a check is for.

\subsection{The Rule Nobody Wrote Down}

The most valuable check in that output is the one that sounds odd: \textbf{every crate must have at
least one example of its own level}.

\begin{itemize}
    \item An API with zero examples passes every other check trivially — it compiles, it lints, and
          nothing calls it.
    \item It then rots silently: in this project a whole board module spent weeks documented and
          unbuildable, and the only symptom was that nobody used it.
    \item The check is cheap and it fails loudly, which is the entire point of a gate.
\end{itemize}

\subsection{What a Gate Cannot Do}

It cannot tell you the LED lit. This project's gate was green on a firmware that flashed successfully
and did not blink — four green steps, no light. That is not a failure of the gate; it is the boundary
of what a build-time check can know, and it is why the hardware checks in chapter~14 exist as a
separate, written procedure with its own measurements.


\clearpage

\begin{keytakeaways}
    \item The triple is \codeinline{thumbv8m.main-none-eabihf}: ARMv8-M mainline, no OS, hardware
          float. Getting the float ABI wrong links but produces a slow, subtly different binary.
    \item The bootloader owns the first 24 KiB. \codeinline{memory.x} starts flash at
          \codeinline{0x08006000}, and a script refuses any image that would start lower.
    \item \codeinline{--nmagic} is not optional: without it the linker emits a segment whose file
          offset disagrees with its load address, and the programmer writes to the wrong place.
    \item Two builds lied in this project: a \codeinline{build.rs} that did not watch its own C file,
          and a book build that compiled yesterday's sources because of a missing
          \codeinline{--strip-components}.
    \item A gate that runs without the board is worth more than a checklist, because it runs every
          time.
\end{keytakeaways}


% ========== PART II: HARDWARE ABSTRACTION LAYER ==========
\part{Hardware Abstraction Layer (HAL)}
\label{part:hal}

% Chapter 4: GPIO (written)
% Chapter 4: GPIO Implementation
% File: chapters/chapter-04-gpio/chapter-04-gpio.tex

\chapter{Phase A: GPIO Implementation}
\label{ch:chapter-04}
\chapthumb{4}

\begin{quote}
    \textit{``A pin is not an abstraction problem. It is a register problem that everybody thinks is
    an abstraction problem.''}
\end{quote}

A GPIO is the simplest peripheral on any microcontroller and the one that produces the most
over-engineered code. This chapter takes the opposite route: it starts from the registers on this
part, derives the smallest abstraction that is still safe, and shows the three defects that the
bigger abstraction — 605 lines of type-state pins — did not prevent.

Everything here is measured on the XIAO MG24 Sense (EFR32MG24B220F1536IM48-B), and the numbers are
read back over SWD rather than assumed.

\minitoc

\clearpage

\section{The Registers, and How Few There Are}
\label{sec:gpio-registers}

A GPIO port on the EFR32MG24 is not a black box. It is five things:

\begin{table}[htbp]
    \centering
    \footnotesize
    \begin{tabular}{@{}p{0.18\textwidth}p{0.42\textwidth}p{0.30\textwidth}@{}}
        \toprule
        \textbf{Register} & \textbf{What it decides} & \textbf{Access} \\
        \midrule
        \codeinline{MODEL} & mode of pins 0--7 & read-modify-write, per pin \\
        \codeinline{MODEH} & mode of pins 8--15 & only on PORTA and PORTC here \\
        \codeinline{DOUT} & the output latch, one bit per pin & read-modify-write \\
        \codeinline{DIN} & the pin as the outside world holds it & read only \\
        \codeinline{CMU\_CLKEN0} & whether the port answers at all & gate before the first access \\
        \bottomrule
    \end{tabular}
    \caption{The GPIO port, in full. Five registers, and only one of them is optional.}
    \label{tab:gpio-registers}
\end{table}

The modes are numbers, not types:

\begin{minted}{rust}
// Measured from the PAC's field encodings, and confirmed by the effect on the pin.
const DISABLED: u32 = 0;   // the pad is off; an analog peripheral may claim it
const INPUT:    u32 = 1;   // input buffer on, no pull
const PUSH_PULL:u32 = 4;   // driven high or low
const WIRED_AND: u32 = 8;  // open drain: 0 pulls low, 1 lets the pull-up win
const WIRED_OR:  u32 = 9;  // open drain, inverted sense
\end{minted}

Two of these decide most of the code in this book. \codeinline{PUSH_PULL} drives an LED and a
piezo. \codeinline{WIRED_AND} is what makes a bit-banged I2C bus possible, because ``output 1''
means ``release the line'' and the external pull-up makes it a real 1 — which is exactly what an
open-drain bus requires and exactly what a push-pull pin must never do.

\subsection{The Clock Gate Comes First}

The most expensive mistake in this whole project was not a wrong bit. It was touching a register
whose bus clock was off:

\begin{minted}{rust}
// Before: the write faulted, and the fault handler did not exist yet.
gpio.porta_dout().modify(|_, w| w.dout7().set_bit());

// After: gate the port, then talk to it.
dp.cmu_s.clken0().modify(|_r, w| w.gpio().set_bit());
gpio.porta_dout().modify(|_, w| w.dout7().set_bit());
\end{minted}

The measurement, taken from the fault status register after the fact:
\codeinline{CFSR = 0x8200} (a precise bus fault) with \codeinline{BFAR} pointing at the GPIO address
that was written. A GPIO register with no clock is not ``reads as zero'' — it is a bus fault, and if
the vector table still points at a default handler the board simply stops.

\subsection{Not Every Port Has Every Register}

\begin{minted}{rust}
// This package is a 48-pin QFN: PORTB and PORTD expose eight pins, not sixteen.
// MODEH covers 8..15, so on those two ports it does not exist at all.
match port {
    Port::A | Port::C => { /* MODEL and MODEH */ }
    Port::B | Port::D => { /* MODEL only */ }
}
\end{minted}

That is not a simplification invented for this book; it is what the SVD says, and a HAL that
pretends otherwise fails to compile \emph{after} someone has written the code. Two of the eight
compile errors during this project's pin primitive came from exactly that: the generated register
blocks differ per port, and a macro written for one does not typecheck for the other.

\section{A Pin Is Three Numbers}
\label{sec:pin-three-numbers}

The type-state approach to GPIO is seductive: make an output pin a different \emph{type} from an
input pin, make a pin claimed by one driver unclaimable by another, and let the compiler refuse the
mistakes. This book's repository had exactly that — 605 lines of it — and it is now retired, for a
reason worth stating precisely.

\subsection{What the Small Version Looks Like}

The whole abstraction is a port, a bit and a polarity:

\begin{minted}{rust}
pub enum Port { A, B, C, D }

pub struct Output<'a> {
    gpio: &'a pac::gpio_s::RegisterBlock,
    port: Port,
    bit: u8,
    active_low: bool,
}

impl<'a> Output<'a> {
    pub fn on(&self)  { self.write(!self.active_low) }
    pub fn off(&self) { self.write(self.active_low) }
    pub fn toggle(&self) {
        let current = read_dout(self.gpio, self.port);
        self.write(current & (1 << self.bit) == 0);
    }
}
\end{minted}

Everything else is a register write:

\begin{minted}{rust}
fn write_dout(gpio: &pac::gpio_s::RegisterBlock, port: Port, bit: u8, high: bool) {
    // Read-modify-write, always: DOUT is one register for sixteen pins.
    let set = |r: u32, w: &mut pac::W<_>| {
        let mask = 1u32 << bit;
        let value = if high { r | mask } else { r & !mask };
        unsafe { w.bits(value) }
    };
    match port {
        Port::A => gpio.porta_dout().modify(set),
        Port::B => gpio.portb_dout().modify(set),
        Port::C => gpio.portc_dout().modify(set),
        Port::D => gpio.portd_dout().modify(set),
    }
}
\end{minted}

That is the entire output path: 40 lines including the four-way match. The input path is shorter
still — \codeinline{DIN} needs no mask, only a shift — and the mode setter is the read-modify-write
of \codeinline{MODEL} or \codeinline{MODEH} shown in the previous section.

\subsection{Why the Larger Version Was Retired}

The type-state layer was not wrong in principle. It was wrong here, in three ways that only show up
in a real project:

\begin{enumerate}
    \item \textbf{Nothing used it.} Its only consumer was a board module that a newer facade
          replaced. An API with zero users compiles, lints clean, and rots: at the time it was
          retired, its README documented the user LED as \codeinline{PB2} active-\emph{high}, when
          the board's LED is \codeinline{PA7} active-\emph{low}.
    \item \textbf{It moved the wiring into the HAL.} To hand out a typed \codeinline{LED} the HAL had
          to know which pin the LED was on — a board fact in the wrong layer. The layered rule from
          section~\ref{sec:layer-boundaries} is that the HAL names no board wiring, and the typed
          pins were the mechanism by which it did.
    \item \textbf{It cost more than it prevented.} The three real defects in this project's GPIO
          code — a blind \codeinline{MODEL} write, a mode value with the wrong width, and an
          open-drain pin configured push-pull — were all invisible to the type system, because all
          three are legal code with the wrong \emph{number} in it. Types catch the mistakes you can
          name; registers catch the rest.
\end{enumerate}

\subsection{What Replaced the Type System}

An ownership rule that is checkable without generics: the board takes the peripherals once
(\codeinline{Board::take()}), and every pin comes out of it with a port and a bit named at the call
site:

\begin{minted}{rust}
let board = Board::take().expect("peripherals taken twice");
let led     = board.output_on(Port::A, 7, true);      // active low
let button  = board.input_on(Port::C, 1);             // pull-up fitted on the base
let buzzer  = board.output_on(Port::C, 3, false);
\end{minted}

The exclusivity that the type system used to provide is now the PAC singleton's job — one
\codeinline{take()} per program — and the wiring that the HAL used to hide lives in whichever crate
owns the board. Fewer lines, fewer layers, and the two facts a reader needs (which pin, which
polarity) are in the code they are reading.

\section{Read-Modify-Write, or Why the Bus Stopped Working}
\label{sec:gpio-rmw}

One register, sixteen pins. That sentence is the whole section, and it caused one of the most
confusing failures in this project.

\subsection{The Defect}

Setting a pin's mode looks like a write. It is a write to a register that also holds fifteen other
pins:

\begin{minted}{rust}
// Wrong: this configures *every* pin on the low half of PORTC, not just PC3.
gpio.portc_model().write(|w| w.mode3().pushpull());

// Right: read what is there, change one field, write it back.
gpio.portc_model().modify(|_r, w| w.mode3().pushpull());
\end{minted}

The wrong version does not fail. It reconfigures the fifteen other pins to their reset value — and on
this board the other pins on that half-port include the two that carry an I2C bus. A bus that had
just been proved working stops answering, with no error anywhere, because the pins that carried it
went back to \codeinline{DISABLED}.

\subsection{How It Was Found}

The failure looked like an I2C problem: the device stopped acknowledging, a scan returned nothing,
and reading the bus's own registers showed the peripheral configured exactly as intended. The
registers that mattered were the \emph{pin} ones:

\begin{minted}{text}
PORTC_MODEL = 0x00000000      # every pin on bits 0..7 disabled
PORTC_DOUT  = 0x00000030      # the bus lines still latched high
\end{minted}

Both lines at once: the latch said high, the pads were off. That combination can only come from a
write that reset the mode register, and a search for \codeinline{.write(} on that register found it
in one place — the function that was setting a buzzer pin.

\subsection{The Rule}

Every write to a GPIO configuration register in this project now goes through
\codeinline{.modify()}. The only \codeinline{.write()} calls are to registers that genuinely hold one
field and no other pin — and even then, the code says so in a comment, because the next person will
copy the pattern somewhere it does not hold.

The same rule applies to \codeinline{DOUT}: a blind \codeinline{write} there sets every pin on the
port high or low at once. It is the reason the pin primitives in this repository all take a port and
a bit, and never a whole register value:

\begin{minted}{rust}
// The signature is the rule: there is no way to write DOUT without saying which bit.
fn write_dout(gpio: &RegisterBlock, port: Port, bit: u8, high: bool)
\end{minted}

\subsection{A Second, Quieter Version of the Same Bug}

\codeinline{modify()} on a register whose clock is off is a bus fault, but \codeinline{modify()} on
a register with a \emph{write-only field} is worse: it reads back a value the field does not
contribute to, so the read-modify-write silently clears it. This is why the mode values in this
project are written as an explicit numeric field of the right width rather than as a chain of
boolean writers, and why the IADC's configuration registers (chapter 7's sibling problem) are written
once, in one function, in a documented order — conditional writes into a synchronising peripheral
lose the field you did not intend to touch.

\section{Open Drain, and the Line That Has Two Masters}
\label{sec:gpio-open-drain}

An I2C bus has two wires and, on each of them, two possible masters, each of which may want to pull
the line low while the other wants it high. The electrical answer is the reason the mode list has a
\codeinline{WIRED\_AND} in it.

\subsection{What WIRED\_AND Does}

\begin{minted}{text}
mode = WIRED_AND   (8)
  DOUT bit = 0  ->  the pad pulls the line to ground
  DOUT bit = 1  ->  the pad releases the line (high impedance)
\end{minted}

The line is held high by an external resistor. So ``1'' is not something the pin produces — it is
something the pin \emph{stops preventing}. Two devices can then share the wire: any one of them can
pull it low, and the only way it goes high is if nobody does. That is what makes arbitration
possible, and it is why a push-pull pin must never be connected to a shared bus: driving it high
against another device driving low is a short circuit through two transistors.

\subsection{The Measurement That Proved the Pull-Ups}

The IMU could not be found for two sessions. Before blaming the pin map, this project measured
whether the bus lines had pull-ups at all:

\begin{minted}{text}
# every pin on PORTB set to input, then read DIN
PORTB_MODEL = 0x11111111     # all inputs
PORTB_DIN   = 0x0000000C     # bits 2 and 3 high
\end{minted}

Bits 2 and 3 are \codeinline{PB2} and \codeinline{PB3}: the two pins of the IMU's bus, reading high
with nothing driving them. That is an external pull-up, present and working — which retired the
``the board has no pull-ups'' hypothesis in one instruction and moved the investigation to the pin
\emph{order}.

\subsection{Bit-Banging an Open-Drain Bus}

With \codeinline{WIRED\_AND} the whole I2C protocol reduces to two writes and one read:

\begin{minted}{rust}
fn sda(&self, high: bool) { write_dout(self.gpio, self.bus.port, self.bus.sda_bit, high) }
fn scl(&self, high: bool) { write_dout(self.gpio, self.bus.port, self.bus.scl_bit, high) }
fn read_sda(&self) -> bool { read_din(self.gpio, self.bus.port) & (1 << self.bus.sda_bit) != 0 }
\end{minted}

\textbf{There is no direction switch.} On an open-drain bus a pin used as an input with the output
latch high is the same electrical state as a released output — which is why a bit-banged I2C master
can be four functions and no state machine, and why the same pin can be written and read in the same
transfer.

\subsection{Half a Bit Period, and Why It Is a Number in Cycles}

Timing on a bit-banged bus is a delay per half period, and the delay is not a duration — it is a
count of loop iterations whose real length depends on the compiler:

\begin{minted}{rust}
// ~5 us at 19 MHz, measured: 60 iterations of this loop on this compiler.
const HALF_BIT_CYCLES: u32 = 60;
\end{minted}

That constant is a measurement, not a datasheet value, and it is the reason chapter~14 exists: the
same class of mistake — \codeinline{asm::delay(n)} assumed to be \codeinline{n} cycles when it is
about three per iteration — turned a 2.7 kHz tone into something else entirely.

\section{What the Vendor's Own Documents Get Wrong}
\label{sec:gpio-vendor-errors}

Three documents describe this board's pins, and two of them disagree with the board. This section is
short and uncomfortable, and it exists because a wrong pin in a comment is worse than no comment: it
is believed.

\subsection{The User LED}

\begin{table}[htbp]
    \centering
    \footnotesize
    \begin{tabular}{@{}p{0.34\textwidth}p{0.26\textwidth}p{0.32\textwidth}@{}}
        \toprule
        \textbf{Source} & \textbf{Says} & \textbf{What it is} \\
        \midrule
        The HAL's retired GPIO README & \codeinline{PB2}, active high & wrong pin, wrong polarity \\
        An earlier board module doc   & \codeinline{PB02/D10}       & wrong pin \\
        Register read-back (measured) & \codeinline{PA7}, active low & \textbf{this one} \\
        \bottomrule
    \end{tabular}
    \caption{The user LED, as documented and as measured. The measurement is a write to
    PORTA\_DOUT bit 7 and a read of PORTA\_DOUT to confirm the latch moved.}
    \label{tab:led-truth}
\end{table}

Why the error survived: nothing in the project ever \emph{used} the wrong pin in a way that failed.
A pin that does not exist toggles a latch nobody looks at, and the code still returns
\codeinline{Ok}. Only the blink example — the one that must be seen — exposed it.

\subsection{The IMU's Two Wires}

The Arduino variant for this board declares:

\begin{minted}{text}
PIN_WIRE_SDA1 = PB3
PIN_WIRE_SCL1 = PB2
\end{minted}

The board is wired the other way round. Measured, in this order:

\begin{enumerate}
    \item a bit-banged scan of all 128 addresses on the pair \codeinline{(PB2, PB3)} with
          SDA=\codeinline{PB2} found a device at \codeinline{0x6A};
    \item the same scan with the roles swapped found nothing;
    \item the device at \codeinline{0x6A} answered \codeinline{WHO\_AM\_I = 0x6A}, which is the
          LSM6DS3TR-C's identity register.
\end{enumerate}

Three independent facts — the scan, the negative control, and the identity register — agree. The
variant file does not, and it is the file a reader would trust first.

\subsection{The Lesson, Which Is Not ``the Vendor Is Wrong''}

The vendor's documents are mostly right and are the only reason this works at all: the SVD, the
reference manual, the bootloader layout and the radio headers are all correct and all essential. But
the pin map of a \emph{board} is not the same document as the pin map of a \emph{chip}, and it is
maintained by people assembling a product, not by the silicon team.

So the rule this project adopted: \textbf{every pin fact is measured on the board, and the
measurement is written next to the claim}. Two of the three corrections above cost an afternoon
each. Both would have been found in ten minutes by reading a register instead of a comment.

\section{Proving a Pin Without Trusting Your Eyes}
\label{sec:gpio-proving}

The blink is the first program every embedded developer writes, and it is the first one that lies. A
blink that does not blink has four different causes that look identical: the code never ran, the pin
is wrong, the polarity is wrong, or the LED is not connected. This section is about telling them
apart with instruments that are already there.

\subsection{The Case That Made This Necessary}

This project had a blink that flashed successfully and did not blink. The write verified, the
read-back matched, the board was silent. Every step of the toolchain reported success:

\begin{minted}{text}
cargo build            Finished `release` profile [optimized] target(s)
python3 scripts/elf-load-range.py   load_min=0x08006000 load_max=0x08006813
probe-rs download      Finished in 0.41s
probe-rs verify        Finished in 0.28s
\end{minted}

Four green steps and no light. The only way forward was to stop asking the board to show and start
asking the register to tell.

\subsection{Reading the Latch Over SWD}

The probe can read any address while the program runs. The output latch is one word, and it toggles
or it does not:

\begin{minted}{bash}
# halt, read DOUT twice, resume
openocd -c "init; halt" -c "mdw 0x5003C040 1" -c "sleep 500" -c "mdw 0x5003C040 1" -c "resume; shutdown"

# 0x5003c040: 00000080      <- bit 7 high
# 0x5003c040: 00000000      <- bit 7 low, 500 ms later
\end{minted}

That pair of readings is a proof of the \emph{program}: the latch is moving at the rate the code
intended. What it is not is a proof of the \emph{board}: a latch that toggles and an LED that stays
dark means the pin is not the LED's pin, or the LED is active-low and the code drives it backwards.

The second reading of \codeinline{PORTA\_MODEL} closes the other half:

\begin{minted}{text}
PORTA_MODEL = 0x00000040   # mode 4 (push-pull) in the field for pin 7
\end{minted}

Latch moving, mode correct, and still no light: that combination is now a board question, not a
software one — which is exactly the division of labour a measurement is for.

\subsection{Three Instruments, Three Questions}

\begin{table}[htbp]
    \centering
    \footnotesize
    \begin{tabular}{@{}p{0.26\textwidth}p{0.34\textwidth}p{0.34\textwidth}@{}}
        \toprule
        \textbf{Instrument} & \textbf{Question it answers} & \textbf{How it is read} \\
        \midrule
        The LED & is the whole chain alive? & eyes, and only after the two below \\
        \codeinline{DOUT}/\codeinline{MODEL} over SWD & is the program right? & \codeinline{mdw} on the port \\
        A mailbox in RAM & did the code reach this point? & \codeinline{mdw} on a magic word \\
        \bottomrule
    \end{tabular}
    \caption{Order matters: the LED answers last, because it answers slowly and ambiguously.}
    \label{tab:gpio-instruments}
\end{table}

\subsection{The Mailbox, Which Is the Real Lesson}

The instrument that made the rest of this project possible is not a pin at all. It is a small block
of RAM — a magic word, a sequence number, a status and a message — that a program writes before it
does anything it might not survive:

\begin{minted}{rust}
// 0x2003F000: magic "MBOX", seq, status (0 info / 1 pass / 2 fail), len, message bytes
mailbox::write(Status::Pass, b"cap=0x20000000 4000@400Hz real=clean");
\end{minted}

Read with one \codeinline{mdw} over SWD, it answers questions that no LED can: did the classifier
load its calibration, how many samples did the ring hold, which of the two I2C aliases answered. It
costs 300 bytes of RAM and it replaced an entire class of guesswork — which is why every example in
this book reports through it rather than through a serial port that may not exist.


\clearpage

\begin{keytakeaways}
    \item A GPIO port on this part is five registers: \codeinline{MODEL}, \codeinline{MODEH},
          \codeinline{DOUT}, \codeinline{DIN} and the port's clock gate.
    \item \codeinline{MODEH} does not exist on every port: on PORTB and PORTD this package has eight
          pins, so those pins live in \codeinline{MODEL} alone.
    \item A pin is a port, a bit and a polarity. That is enough to drive an LED, read a button and
          open-drain a bus, and it is 189 lines — not 605.
    \item Writing \codeinline{MODEL} or \codeinline{MODEH} without reading first reconfigures every
          other pin on the same half-port, which is how a working bus stops working.
    \item The LED on this board is \codeinline{PA7}, active \emph{low}. Two of the vendor's own
          documents disagree with that.
    \item A blink is proven by the register toggling, not by the LED looking on.
\end{keytakeaways}


% Chapter 5: Clocks (written)
% Chapter 5: Clock Management
% File: chapters/chapter-05-clocks/chapter-05-clocks.tex

\chapter{Phase A: Clock Management and System Configuration}
\label{ch:chapter-05}
\chapthumb{5}

\begin{quote}
    \textit{``Every timing bug is a clock bug wearing a disguise.''}
\end{quote}

Nothing on this part runs without a clock, and almost nothing that goes wrong is blamed on one. This
chapter is the record of four separate defects in one project that were all clocks: a delay that was
three times too short because a loop is not an instruction count, a peripheral that never started
because its oscillator was selected but not enabled, a converter that hung because it was clocked
from the wrong source, and a sampling loop that claimed 400 Hz and ran at 158.

The numbers here are measured with the cycle counter, not read from a datasheet.

\minitoc

\clearpage

\section{The Tree, in the Order You Meet It}
\label{sec:clocks-tree}

A Series 2 clock tree is not complicated once it is drawn as three questions. This section answers
them for this board only, because the full diagram in the reference manual has forty boxes and no
project needs forty.

\subsection{Question 1: Which Oscillator}

\begin{table}[htbp]
    \centering
    \footnotesize
    \begin{tabular}{@{}p{0.20\textwidth}p{0.16\textwidth}p{0.52\textwidth}@{}}
        \toprule
        \textbf{Oscillator} & \textbf{Rate} & \textbf{When it is the right answer} \\
        \midrule
        \codeinline{HFRCO0} & 19 MHz default & the reset state; enough for GPIO, timers and I2C \\
        \codeinline{HFXO0}  & 38.4 MHz crystal & radio, USB, anything needing ppm accuracy \\
        \codeinline{FSRCO}  & 20 MHz fast RC & cheap and fast, but see the IADC below \\
        \codeinline{HFRCOEM23} & 19 MHz & low-energy domain; the IADC and TIMER0 like it \\
        \codeinline{DPLL0}  & up to 78 MHz & the CPU's fast mode, locked to HFRCO or HFXO \\
        \codeinline{LFRCO}/\codeinline{LFXO} & 32.768 kHz & RTC, sleep, wake-up \\
        \bottomrule
    \end{tabular}
    \caption{Six oscillators, and the four that matter for a board with an I2C bus, a buzzer and a
    microphone.}
    \label{tab:clocks-oscillators}
\end{table}

\subsection{Question 2: Which Branch}

Each oscillator feeds branches, and each branch has a selector and a divider:

\begin{minted}{rust}
// The low-energy peripherals (IADC, TIMER0 in this project) are clocked from EM01GRPACLK.
dp.cmu_s.em01grpaclkctrl().modify(|_r, w| w.clksel().hfrcoem23());
\end{minted}

The name matters when reading the manual: \codeinline{EM01GRPACLK} is one branch,
\codeinline{EM23GRPACLK} another, and a peripheral belongs to exactly one. Choosing the right
oscillator on the wrong branch gives a peripheral clocked at zero — which looks identical to a
peripheral whose registers are not being written.

\subsection{Question 3: Which Gate}

Every peripheral has a bus-clock gate in \codeinline{CMU\_CLKEN0} or \codeinline{CMU\_CLKEN1}. On
this part the GPIO bank is bit 0 of \codeinline{CLKEN0} and the MVP's is bit 30 of
\codeinline{CLKEN1}. A register access with the gate closed is a bus fault, precisely reported:

\begin{minted}{text}
CFSR = 0x00008200   # PRECISERR set: a precise data bus fault
BFAR = 0x5003C040   # ...at PORTA_DOUT, whose clock was off
\end{minted}

That pair of registers is the fastest diagnosis in this book. If a peripheral reads as zero, check
the gate; if a write faults, read \codeinline{CFSR} and \codeinline{BFAR} before anything else.

\subsection{The Three Questions, Together}

\begin{minted}{rust}
// Oscillator: force HFRCOEM23 on and wait for it.
dp.cmu_s.clken0().modify(|_r, w| w.hfrcoem23().set_bit());
dp.hfrcoem23_s.ctrl().modify(|_r, w| w.forceen().set_bit());
while !dp.hfrcoem23_s.status().read().rdy().bit_is_set() {}

// Branch: the low-energy domain comes from it.
dp.cmu_s.em01grpaclkctrl().modify(|_r, w| w.clksel().hfrcoem23());

// Gate: the peripheral's own bus clock.
dp.cmu_s.clken0().modify(|_r, w| w.timer0().set_bit());
\end{minted}

Three decisions, three writes, and one wait. The wait is the subject of the next-but-one section: it
is the only one of the four steps that can hang.

\section{The Clock You Forgot You Were Using}
\label{sec:clocks-sysclk}

Every HAL has a constant called something like \codeinline{DEFAULT\_SYSCLK\_HZ}, and it is almost
always wrong for the board it runs on. This section is about the one on this board, because it
silently defines the length of every delay in the project.

\subsection{19 MHz, Not 78}

The EFR32MG24 can run at 78 MHz. It does not, out of reset, and this board's firmware never raises
it. The default is \codeinline{HFRCO0} at 19 MHz, and the measurement that settled it was a stopwatch
rather than a datasheet:

\begin{minted}{text}
# a blink with a 500 ms half-period, counted against a phone stopwatch
# 20 blinks took 20.0 s  ->  1.000 s per cycle  ->  the delay was right at 19 MHz
# the same code timed as if SYSCLK were 78 MHz produced a 4x fast blink
\end{minted}

That is a crude instrument, and it is the right one: a stopwatch over twenty cycles has an error far
below the 4× discrepancy being tested.

\subsection{Why the Error Is Invisible}

A wrong \codeinline{SYSCLK\_HZ} does not produce a wrong-looking program. It produces a program whose
delays are all scaled by the same factor — blinks that are too fast, tones that are too high, buses
that are too slow or too fast. Each of those looks like a bug in that peripheral:

\begin{table}[htbp]
    \centering
    \footnotesize
    \begin{tabular}{@{}p{0.30\textwidth}p{0.30\textwidth}p{0.32\textwidth}@{}}
        \toprule
        \textbf{Symptom} & \textbf{Blamed on} & \textbf{Actually} \\
        \midrule
        LED blinks at 4× the rate & the delay loop & \codeinline{SYSCLK\_HZ} too high \\
        Tone is 4× the intended pitch & the timer's TOP & same constant \\
        I2C bus works on one board, not another & the devices & same constant, other clock \\
        Buzzer sounds, then the beep pattern is wrong & the pattern code & same constant \\
        \bottomrule
    \end{tabular}
    \caption{One wrong constant, four different suspects. This is why it is measured once and written
    down next to the definition.}
    \label{tab:sysclk-symptoms}
\end{table}

\subsection{The Right Shape for the Constant}

The version that survives contact with a second board is not a number in the HAL. It is a number the
board passes in:

\begin{minted}{rust}
// The HAL knows how to delay; it does not know how fast this board's clock runs.
let mut delay = Delay::with_sysclk(syst, 19_000_000);
\end{minted}

and in this project's board layer, one place says it out loud:

\begin{minted}{rust}
pub const DEFAULT_SYSCLK_HZ: u32 = 19_000_000;   // measured, see 05.02
\end{minted}

A second board with a 38.4 MHz crystal changes one line, and every delay, tone and bus timing
follows it. A HAL with the constant baked in cannot be reused, and — worse — fails in a way that
looks like the new board's fault.

\subsection{Measuring It Properly}

The stopwatch is enough to catch a factor of four. For anything finer, use the cycle counter, which
counts CPU cycles and therefore cannot disagree with the CPU:

\begin{minted}{rust}
let start = dwt::cycles();
delay.delay_ms(1000);
let elapsed = dwt::since(start);
// elapsed / 19_000_000 must be ~1.0 s; a factor of 3.0 here means the delay's
// internal arithmetic is wrong, not the constant.
\end{minted}

That distinction — constant wrong \emph{versus} arithmetic wrong — is the difference between
recalibrating and debugging, and the cycle counter is what separates them.

\section{Selecting an Oscillator Is Not Starting It}
\label{sec:clocks-starting}

The clock tree has a trap that costs an afternoon the first time and ten minutes every time after:
choosing a source does not turn it on, and a peripheral clocked by a stopped source behaves like a
peripheral that ignores you.

\subsection{The Two-Step That Is Actually Three}

\begin{minted}{rust}
// 1. Ask for it.
dp.cmu_s.clken0().modify(|_r, w| w.hfrcoem23().set_bit());
// 2. Force it on: without this the oscillator idles when nothing demands it,
//    and "nothing" includes the branch you are about to select.
dp.hfrcoem23_s.ctrl().modify(|_r, w| w.forceen().set_bit());
// 3. Wait. This is the step that hangs if the oscillator never comes up.
while !dp.hfrcoem23_s.status().read().rdy().bit_is_set() { /* spin */ }
\end{minted}

Step 2 is the one people miss. It is also the one that makes step 3 a real wait instead of an
instant pass — and an instant pass is worse than a hang, because it looks like success.

\subsection{The Converter That Hung on FSRCO}

The microphone's converter (IADC0) had two separate problems in this project, and one of them was
purely a clock choice. Clocked from \codeinline{FSRCO}, the converter's disable sequence never
completed:

\begin{minted}{text}
# after writing EN = 0 and waiting for the disabling bit to clear
IADC0_STATUS  = ...DISABLING = 1     (still, after 100 ms)
# the next configuration write then faulted, because the peripheral was mid-synchronisation
\end{minted}

The same code, clocked from \codeinline{HFRCOEM23}, disabled in microseconds and never faulted. No
register value changed; only \codeinline{em01grpaclkctrl}. That is the entire difference between a
working microphone driver and an afternoon of bus faults, and it was found by trying the other
oscillator rather than by reading harder.

\subsection{The Rule That Came Out of It}

Before blaming a peripheral, confirm that it is \emph{clocked}, and confirm it in the way the
hardware reports it:

\begin{enumerate}
    \item read the gate bit back — not the value you wrote, the value the register holds;
    \item read the oscillator's ready status, not just its enable bit;
    \item if the peripheral has a synchronisation status, read it: \codeinline{DISABLING},
          \codeinline{SYNCBUSY} and \codeinline{RESETTING} are the hardware telling you the
          difference between ``configured'' and ``ready''.
\end{enumerate}

A peripheral that reads as zero after a correct-looking write has failed one of those three, and the
three have different fixes. Without the read-backs they are indistinguishable — which is why every
driver in this repository returns an error instead of a value when a readiness wait expires.

\section{Delay Without Lying}
\label{sec:clocks-delay}

There are three ways to wait on a Cortex-M and this project has been wrong about all three.

\subsection{asm::delay Is Not a Cycle Count}

\begin{minted}{rust}
// "1000 cycles", according to the name and to every example on the internet.
cortex_m::asm::delay(1000);
\end{minted}

Measured against the cycle counter, one iteration of the loop that \codeinline{asm::delay} expands
to takes about \textbf{three} cycles on this part with this compiler. So
\codeinline{delay(1\_000\_000)} waits three million cycles, not one — a 3× error, and the reason a
2.7 kHz tone computed with it came out as something else entirely.

Two consequences, both practical:

\begin{itemize}
    \item \textbf{Never derive a period from it.} A tone, a baud rate or a bus half-period computed
          with \codeinline{asm::delay} is wrong by a factor nobody wrote down.
    \item \textbf{Do not multiply by three and call it fixed.} The factor is a property of the
          compiler's code generation, not of the hardware; it changed between optimisation levels in
          this project.
\end{itemize}

\subsection{SysTick: Real Milliseconds}

SysTick counts CPU cycles and can therefore produce real time:

\begin{minted}{rust}
pub fn delay_us(&mut self, us: u32) {
    let cycles = (u64::from(us) * u64::from(self.sysclk_hz)) / 1_000_000;
    // SysTick's reload is 24 bits: split anything larger into chunks.
    let mut remaining = cycles;
    while remaining > 0 {
        let chunk = remaining.min(0x00FF_FFFF);
        self.syst.set_reload(chunk as u32);
        self.syst.clear_current();
        self.syst.enable_counter();
        while !self.syst.has_wrapped() {}
        self.syst.disable_counter();
        remaining -= chunk;
    }
}
\end{minted}

It has one owner — the board — and it must never be shared with a scheduler that also wants it. In
this project that rule is enforced by the type system in the smallest possible way: \codeinline{SYST}
is taken once by \codeinline{Board::take()}, and \codeinline{board.delay()} hands it out.

\subsection{Ever the Millisecond: Do Not Divide Early}

A defect worth its own paragraph, because it looks like arithmetic hygiene:

\begin{minted}{rust}
// The sample interval arrives as microseconds: 2500 us, i.e. 416 Hz.
let dt_ms = dt_us / 1000;        // 2 ms
// Every timeout in the state machine is now 20% short.
\end{minted}

Integer division by 1000 truncated 2500 to 2, and every threshold derived from it inherited a 20%
error. The classifier's free-fall window, its landing window and its settle timer were all wrong at
once, and every one of them was measured against real drops before anyone suspected the unit.

The rule: \textbf{carry the finest unit the hardware gives you and divide only where a human reads
it}. This project's classifier now takes microseconds everywhere and converts once, at the boundary,
for display.

\subsection{The Cycle Counter: Measurement, Not Delay}

The third instrument is not for waiting at all:

\begin{minted}{rust}
let start = dwt::cycles();
// ... work ...
let elapsed_us = dwt::cycles_to_us(dwt::since(start), 19_000_000);
\end{minted}

Two registers and no ownership questions: enable \codeinline{TRCENA} in \codeinline{DEMCR}, then
\codeinline{CYCCNTENA} in \codeinline{DWT->CTRL}, then read \codeinline{CYCCNT}. It is how every
timing number in this book was obtained, and it is the only instrument here that cannot lie about
the clock because the CPU \emph{is} the clock.

One trap, measured: a counter that was never enabled reads as zero, and zero elapsed cycles looks
exactly like a loop that ran impossibly fast. Compare two consecutive reads before trusting either.

\section{What to Write Down, and Where}
\label{sec:clocks-write-down}

A clock fact has three places it can live, and only one of them survives the project.

\subsection{The Three Places}

\begin{table}[htbp]
    \centering
    \footnotesize
    \begin{tabular}{@{}p{0.28\textwidth}p{0.30\textwidth}p{0.34\textwidth}@{}}
        \toprule
        \textbf{Place} & \textbf{Lifetime} & \textbf{What happens to it} \\
        \midrule
        In the chat where you found it & minutes & lost \\
        In a commit message & months, if you search & nobody searches \\
        In the source, next to the constant it explains & the life of the code & \textbf{read by the next person} \\
        \bottomrule
    \end{tabular}
    \caption{The only durable location is the one somebody will read while editing the thing it
    describes.}
    \label{tab:where-facts-live}
\end{table}

\subsection{How the Numbers Here Are Written}

Every measured constant in this project carries three things: the number, the unit, and how it was
obtained. For example:

\begin{minted}{rust}
/// Half a bit period on the bit-banged I2C bus, in loop iterations.
///
/// Measured: 60 iterations take ~5 us at SYSCLK = 19 MHz on this board, so ~95 kHz.
/// This is a measurement, not a datasheet value -- `asm::delay` is about three cycles
/// per iteration here, so the divisor is not the clock.
const HALF_BIT_CYCLES: u32 = 60;
\end{minted}

and for the value that everything else depends on:

\begin{minted}{rust}
/// SYSCLK as measured with a stopwatch over 20 blink cycles: 19 MHz (the HFRCO default).
/// The part supports 78 MHz; this board's firmware never selects it.
pub const DEFAULT_SYSCLK_HZ: u32 = 19_000_000;
\end{minted}

\subsection{The Table That Belongs in Every Project}

\begin{table}[htbp]
    \centering
    \footnotesize
    \begin{tabular}{@{}p{0.28\textwidth}p{0.24\textwidth}p{0.40\textwidth}@{}}
        \toprule
        \textbf{Quantity} & \textbf{Value here} & \textbf{How it was measured} \\
        \midrule
        SYSCLK & 19 MHz & stopwatch over 20 blink cycles \\
        \codeinline{asm::delay} cost & ~3 cycles/iteration & cycle counter, two build profiles \\
        I2C half period & ~5 $\mu$s (60 iterations) & scope-free: 95 kHz against a known device \\
        Buzzer tone & 2.7 kHz, TOP = 3518 & TIMER0 TOP read back over SWD \\
        IMU sample rate & 416 Hz configured & \codeinline{CTRL1\_XL} register + measured intervals \\
        Sampling loop & 158 Hz when claiming 400 & cycle counter, per-sample interval recorded \\
        \bottomrule
    \end{tabular}
    \caption{The clock facts of this project. The last line is the one that mattered: a loop that
    reports a rate it does not achieve poisons every timing test downstream.}
    \label{tab:clock-facts}
\end{table}

\subsection{The One That Generalises}

The last row of that table is the lesson of this chapter. The firmware said 400 Hz because the
\emph{intended} interval was 2500 $\mu$s. It ran at 158 Hz because the display flush blocked for tens
of milliseconds inside the loop. Every threshold in the state machine was therefore 2.5× fast, and
three rounds of drop tests were interpreted against a wrong time base before anyone noticed that four
different drops showed the same 0.62 s of calm.

The fix was not a better delay. It was to \textbf{measure the interval between samples and pass the
measurement on}:

\begin{minted}{rust}
let now = dwt::cycles();
let dt_us = dwt::cycles_to_us(dwt::since(self.last), sysclk_hz);
self.last = now;
// The classifier is told what actually happened, not what was planned.
detector.update(&accel, dt_us);
\end{minted}

Every timing defect in this project — and there were four — would have been caught on the day it was
introduced by that three-line habit.


\clearpage

\begin{keytakeaways}
    \item A clock tree is three separate decisions: which oscillator, which branch, and which
          peripheral gate. Getting one of the three wrong produces the same symptom — a peripheral
          that reads as zero.
    \item SYSCLK on this board is the HFRCO default, measured at 19 MHz. The part can do 78 MHz; the
          board does not, until someone configures it.
    \item Selecting an oscillator is not starting it. Force it on, then wait for the ready bit.
    \item \codeinline{asm::delay(n)} is not \codeinline{n} cycles — it measured about three per
          iteration here, which is a 3× error in every derived period.
    \item A loop's \emph{intended} rate is not its actual rate. Measure the interval between samples
          and pass the measurement on, never the intention.
\end{keytakeaways}


% Chapter 6: USART
\chapter{Phase B: Communication Peripherals - USART}
\label{ch:chapter-06}
\chapthumb{6}
\section*{This chapter content will be added}
\clearpage

% Chapter 7: I2C
% Chapter 7: I2C Master Mode (written; the rest are placeholders until they are too)
% Chapter 7: I2C Master Mode Implementation
% File: chapters/chapter-07-i2c/chapter-07-i2c.tex

\chapter[I2C Master Mode]{I2C Master Mode: Two Buses, One That Did Not Work, and Two Lines That Were Swapped}
\label{ch:chapter-07}
\chapthumb{7}

\begin{quote}
    \textit{``The datasheet is a description of the hardware. It is not the hardware.''}
    \\
    -- the lesson of this chapter
\end{quote}

I2C looks like the easiest peripheral in the machine: two open-drain wires, an address, a
byte, an acknowledge. It is also the peripheral where this project lost the most time, and
every hour was lost for the same reason — trusting a document over a measurement.

This chapter builds an I2C master from the register level up, and then does something most
texts skip: it shows what happens when the on-chip controller does not work, how to prove
that it is the controller and not your code, and how to build a transport that does work
instead. Along the way it solves a real failure: a sensor that never answered because the
two bus lines were the other way round, and every scan of the ``correct'' pins returned an
empty bus.

\minitoc

\clearpage

% ============================================================================
% SECTION 07.01
% ============================================================================
\section{I2C Protocol Fundamentals}
\label{sec:i2c-fundamentals}

I2C is two wires and a clock, and everything else follows from one electrical decision: the
lines are \structure{open drain}. A device never drives a line high; it releases it and a
pull-up resistor does the work. Two consequences drive every design decision that follows.

\begin{keytakeaways}
    \item A device can only pull a line \emph{low}. ``High'' is the absence of anybody
          pulling.
    \item Because releasing is not a conflict, two devices can never damage each other by
          talking at once — the bus resolves it as a wired-AND.
    \item But because a ``1'' comes from a resistor, the rise time is set by
          $R \times C$: the bus capacitance is what limits the clock rate and the wire length.
\end{keytakeaways}

\subsection{The Transaction, in Voltage}

A transfer is a sequence of edges, and each of them means something specific:

\begin{enumerate}
    \item \textbf{START}: SDA falls while SCL is high. Nothing else on the bus produces that
          pattern, which is why it can be used as a marker without a separate wire.
    \item \textbf{Address}: seven bits, most significant first, then a direction bit
          ($0$ = write, $1$ = read).
    \item \textbf{ACK}: the addressed device pulls SDA low for the ninth clock. If nobody
          pulls, the byte is unacknowledged — the single most useful diagnostic on the bus.
    \item \textbf{Data}: bytes, each followed by an ACK, in the direction agreed.
    \item \textbf{STOP}: SDA rises while SCL is high.
\end{enumerate}

The detail that catches everybody the first time: \textbf{SDA may only change while SCL is
low}. Both the START and the STOP are defined as the violation of that rule, which is
exactly why they are unambiguous.

\subsection{The Electrical Reality on This Board}

The EFR32MG24 does not have a dedicated I2C pin pair in the sense a beginner expects. The
peripheral observes ordinary GPIO pins, so the same pin can be used as a GPIO, as an I2C
line, or as something else — and the routing is a separate configuration step.

That gives two ways to implement a bus on this part:

\begin{description}
    \item[The on-chip controller (I2C0/I2C1).] The hardware generates the edges, handles
          addressing, ACK and clock stretching, and interrupts the core when done.
    \item[A bit-banged master.] Software drives the pins. Slower, more code, and completely
          under your control — which matters enormously when the first option misbehaves.
\end{description}

This chapter uses both, and the reason is not stylistic. On this part the on-chip
controller could not complete a read, and the proof of that is the subject of the next
section.

\subsection{Open Drain in the Registers}

An open-drain pin is not configured ``as output and set to 1''. It is configured in a mode
where the output driver can only sink current, and the value written has the opposite
meaning to the one intuition suggests:

\begin{minted}{rust}
// The line is released (reads high through the pull-up) when the output latch is 1,
// and pulled low when it is 0. This is not a convention we chose: it is what WIREDAND
// means, and it is why a bit-banged master writes `1` to let go.
porta_dout().write(|w| unsafe { w.dout(1 << pin) }); // release
porta_dout().write(|w| unsafe { w.dout(0) });        // drive low
\end{minted}

Getting this backwards produces the most confusing failure in the whole chapter: a bus that
appears dead, while the code looks correct.

\begin{warningbox}
    \textbf{The mistake that cost the most time in this project.} The GPIO was configured for
    the bus, the address was right, the device was powered, and every scan returned nothing —
    because the pin was in a mode where writing \codeinline{1} \emph{drove the line high}
    instead of releasing it. Two open-drain lines both driven high is a short through the
    drivers, and the resulting waveform is not an I2C bus at all.
\end{warningbox}


% ============================================================================
% SECTION 07.02
% ============================================================================
\section{The Controller That Could Not Read}
\label{sec:i2c-controller-cannot-read}

This section is a worked example of a specific and unglamorous engineering skill: proving
that a peripheral, and not your code, is at fault. It matters because the alternative —
guessing — is unbounded. You can spend days rewriting a driver that was never the problem.

\subsection{The Symptom}

An I2C0 master was configured with the PAC: bus clock enabled, pins routed, clock divider
set, address written, and a read requested. Writes completed. Reads did not: the transfer
never finished, and the state register reported the controller waiting in the transmit path
with the bus already acknowledged as finished for the write half.

The measurement that matters, taken from \codeinline{I2C0->STATE} after a write-then-read
sequence:

\begin{minted}{text}
I2C0 STATE = 0xd7
    BUSY         = 1
    TRANSMITTER  = 1
    NACKED       = 1
    RXDATAV      = 0     <- no receive data will ever arrive
\end{minted}

\codeinline{RXDATAV = 0} is the whole story: the controller accepted the read command and
never received a byte into its buffer. There is nothing to read out, so any polling loop
written against \codeinline{RXDATAV} waits forever.

\subsection{What Was Ruled Out, and How}

A diagnosis is only worth reading if it says what it eliminated and on what evidence. The
table below is the actual list, in the order it was worked through.

\begin{table}[H]
\centering
\begin{tabular}{p{4.6cm}p{9.2cm}}
\hline
\textbf{Hypothesis} & \textbf{Evidence against} \\
\hline
Wrong address & The same address acknowledged during the write half of the transfer. \\
Wrong pins & A bit-banged master on the same two pins read the device immediately. \\
Missing pull-ups & The pins read high with every GPIO configured as input, no internal
                  pull: the pull-ups are on the board. \\
Missing bus clock & \codeinline{CMU\_CLKEN0.I2C0} read back set; with it cleared the
                    peripheral faults instead of running. \\
Bad clock divider & Two dividers, including the slowest available, produced the same
                    \codeinline{STATE}. \\
The device itself & The device answered a bit-banged transaction on the same bus, at the
                   same address, in the same second. \\
\hline
\end{tabular}
\caption{Hypotheses eliminated before concluding that the on-chip controller was at fault.}
\label{tab:i2c-elimination}
\end{table}

Fourteen measurements across these hypotheses produced the same \codeinline{STATE} every
time. That repetition is the point: a single failing attempt is consistent with a bug, a
dozen identical failures along different paths are consistent with the hardware.

\begin{notebox}
    \textbf{The control experiment is the one that decided it.} The bit-banged master used
    the same pins, the same pull-ups and the same device. It read the correct identifier on
    the first attempt. That single result moved the fault from ``our driver'' to ``this
    controller, on this part, in this configuration'', and it took two minutes to run.
\end{notebox}

\subsection{What To Do About It}

There are three honest responses to a peripheral that does not work, and only one of them is
acceptable in a project that has to ship.

\begin{enumerate}
    \item \textbf{Rewrite the driver.} Cheap to try, unbounded to continue. Only justified
          while the control experiment has not been run.
    \item \textbf{Build a transport that works.} Slower, but fully understood, and it turns
          the bus into something you can debug byte by byte. This is what this project did,
          and it is the subject of the next section.
    \item \textbf{Find the missing configuration.} The correct answer when it exists. It does
          not always: this chapter does not claim the controller is broken in general, only
          that it could not complete a read under every configuration that was tried, and
          that a working alternative was available.
\end{enumerate}

\begin{warningbox}
    \textbf{Write down which one you chose, and why.} A project that silently switches to a
    workaround loses the knowledge that it did. This project keeps the fourteen measurements
    in \filepath{docs/04-reference/I2C0-READ-DEFECT.md}, so that a later engineer — or a
    later silicon revision — can pick up the investigation instead of repeating it.
\end{warningbox}

\subsection{Why This Belongs in a Book}

Because it is the normal experience of driving new hardware, and textbooks rarely show it.
The registers are documented, the sequence is correct, and the thing still does not work.
The skill being taught here is not I2C: it is how to reduce a hardware failure to a
statement you can defend, and then how to keep shipping while it is unresolved.


% ============================================================================
% SECTION 07.03
% ============================================================================
\section{A Transport That Works: The Bit-Banged Master}
\label{sec:i2c-bitbang}

Once the on-chip controller is out of the picture, the bus becomes software, and the
software can be made small enough to hold in your head.

\subsection{The Whole Protocol in Four Primitives}

Every I2C transfer — a register write, a 12-byte burst, a scan of 128 addresses — reduces
to four operations on the two lines:

\begin{minted}{rust}
/// Let SDA float (the pull-up makes it a 1) or pull it to ground.
fn sda(&mut self, high: bool);
/// Same for the clock.
fn scl(&mut self, high: bool);
/// Let the bus settle for half a bit period, then sample SDA back.
fn read_sda(&self) -> bool;
/// Half a bit period of wall time.
fn half_bit(&self);
\end{minted}

With those, a byte is a loop and a START is two writes:

\begin{minted}{rust}
fn start(&mut self) -> bool {
    self.sda(true);
    self.scl(true);
    self.half_bit();
    self.sda(false);        // SDA falls while SCL is high: START
    self.half_bit();
    self.scl(false);
    self.half_bit();
    true
}

fn write_byte(&mut self, byte: u8) -> bool {
    for bit in (0..8).rev() {
        self.sda(byte & (1 << bit) != 0);
        self.half_bit();
        self.scl(true);
        self.half_bit();
        self.scl(false);    // data is sampled by the device on the rising edge
    }
    // Ninth clock: release SDA and see whether anybody pulls it down.
    self.sda(true);
    self.half_bit();
    self.scl(true);
    self.half_bit();
    let acked = !self.read_sda();
    self.scl(false);
    acked
}
\end{minted}

\begin{tipbox}
    The acknowledge is the only place where the master \emph{reads}. Reading SDA during a
    write is what allows the driver to report ``nobody is home'' instead of silently
    succeeding — and that distinction is what makes a bus scan possible at all.
\end{tipbox}

\subsection{Timing Without a Timer}

The half-bit period does not need an interrupt or a timer peripheral. A counted loop is
enough, and its length is calibrated against something known: the CPU clock.

\begin{minted}{rust}
/// ~5 us at the 19 MHz this board runs at, which is a 100 kHz bus.
const HALF_BIT_CYCLES: u32 = 60;

fn half_bit(&self) {
    cortex_m::asm::delay(HALF_BIT_CYCLES);
}
\end{minted}

\begin{warningbox}
    \codeinline{cortex\_m::asm::delay(n)} does \textbf{not} execute exactly $n$ CPU cycles.
    On this part it ran about three times faster than the name suggests, which turned a
    nominal 100 kHz bus into a 300 kHz one — and turned a buzzer calibration into a 3$\times$
    error elsewhere in the same project. Anything that claims a frequency should be measured
    with the cycle counter (\codeinline{DWT->CYCCNT}), not assumed from a loop count.
\end{warningbox}

\subsection{The \texttt{embedded-hal} Face}

The transport is more useful if the rest of the ecosystem can use it. \codeinline{embedded-hal}
v1.0 defines the trait, and the implementation is a thin wrapper over the primitives above:

\begin{minted}{rust}
impl<B: I2cBus> embedded_hal::i2c::I2c for SharedBus<'_, B> {
    fn transaction(&mut self, address: u8, operations: &mut [Operation<'_>])
        -> Result<(), Error>
    {
        // One START, the operations in order, one STOP. A register read is a write of the
        // register number followed by a repeated START and a read: getting the repeated
        // START wrong is the second most common cause of a silent bus.
        ...
    }
}
\end{minted}

Designing the driver against the trait rather than against the pin-level functions is what
allowed the same SSD1306, BME68x, PCF8563 and LSM6DS3TR-C drivers to be written once and
used on two different buses — one on PC04/PC05, another on PB2/PB3 — without changing a line
of driver code.

\subsection{Two Buses on One Board}

The board this book was written against has two I2C buses, and the reason is worth recording:

\begin{table}[H]
\centering
\begin{tabular}{llll}
\hline
\textbf{Bus} & \textbf{Pins} & \textbf{Devices} & \textbf{Notes} \\
\hline
I2C on the base & PC04 (SDA), PC05 (SCL) & OLED 0x3C, BME68x 0x76, RTC 0x51 & pull-ups on the base \\
IMU bus         & PB2 (SDA), PB3 (SCL)  & LSM6DS3TR-C 0x6A                & external pull-ups, own power gate \\
\hline
\end{tabular}
\caption{The two I2C buses of the XIAO MG24 Sense on the XIAO Expansion Base.}
\label{tab:i2c-two-buses}
\end{table}

That second row is where the chapter's hardest lesson lives, and it is the subject of the
next section.


% ============================================================================
% SECTION 07.04
% ============================================================================
\section{Two Lines, Swapped: A Case Study in Not Trusting a Pin Map}
\label{sec:i2c-swapped-lines}

A six-axis IMU was supposed to be on this board at address \codeinline{0x6A}. The vendor's
documentation said so, the vendor's example code said so, and the board's own pin map said
which two pins carried the bus. Every scan of those pins returned an empty bus: not a
mismatch, not garbage — nothing at all, as though no device existed.

It existed. The two lines were the other way round.

\subsection{What the Documentation Said}

Three sources, all consistent and all wrong in the same direction:

\begin{itemize}
    \item The vendor's product page: the IMU answers at \codeinline{0x6A}.
    \item The vendor's example: \codeinline{LSM6DS3 myIMU(I2C\_MODE, 0x6A)}, and a library
          header that selects \codeinline{Wire1} rather than the default bus.
    \item The board support package's pin map: \codeinline{SDA1 = PB3},
          \codeinline{SCL1 = PB2}.
\end{itemize}

The third line is the one that matters, and it is the one that is wrong. On this board the
data line is \textbf{PB2} and the clock is \textbf{PB3}.

\begin{warningbox}
    \textbf{Why a swapped pair produces silence rather than an error.} Both lines idle high
    through their pull-ups. If the master drives what it believes to be the clock while the
    device is watching it as data, the START condition — SDA falling while SCL is high —
    never occurs. The device never wakes up, never acknowledges, and never reports anything,
    because no transaction ever began. Every failure mode of a wrong pin assignment looks
    like an empty bus, not like a bad byte.
\end{warningbox}

\subsection{How It Was Found}

The useful part of this story is not the answer; it is the method, which is reusable on any
board with an unknown bus.

\paragraph{Step one: prove the pins are a bus at all.}
Before talking to anything, configure every GPIO as an input with no internal pull, and read
the port. Two pins that read high with no internal help are held up by resistances on the
board — which is what an I2C bus looks like, and nothing else on a general GPIO header does.

\begin{minted}{text}
PORTB_DIN = 0x0C    ->  bit 2 (PB2) and bit 3 (PB3) are both high, no internal pull
\end{minted}

\begin{notebox}
    The first version of this probe was useless: it configured the pins as \emph{outputs}
    (\codeinline{mode 0}), which disables the input buffer, so the result carried no
    information at all. \codeinline{mode 1} — input, no pull — is the measurement.
\end{notebox}

\paragraph{Step two: stop assuming which pin is which.}
Instead of testing one assignment, test the matrix. Every plausible pin pair, in both
orders, against both addresses the device could answer on:

\begin{minted}{text}
8 pin pairs x 2 orders x 2 addresses = 32 probes
result: sda=p1.2 scl=p1.3 addr=0x6a id=0x6a
\end{minted}

That single line is the whole fix. The device had been there the entire time, on the pair
the documentation listed — with the two wires the other way round.

\paragraph{Step three: make the driver's configuration a value, not a hardcoded write.}
The first attempt after the fix still failed, for a second and much more embarrassing reason:
the pin-abstraction's write function always addressed the same port, so changing ``which
bus'' changed nothing about the electrical output while appearing to work. The lesson is
small and general:

\begin{minted}{rust}
// Wrong: the selector looked configurable but every branch wrote PORTB.
fn set(&mut self, pin: u8, high: bool) {
    self.portb_dout().write(|w| unsafe { w.dout(if high { 1 << pin } else { 0 }) });
}

// Right: the port is part of the state, so choosing a bus chooses the registers.
fn write(&mut self, high: bool) {
    match self.pins {
        Pins::Base => self.portc_dout().write(|w| unsafe { w.dout(high as u32) }),
        Pins::Imu  => self.portb_dout().write(|w| unsafe { w.dout(high as u32) }),
    }
}
\end{minted}

\begin{remark}
    \textbf{The pattern to notice.} A parameter that does not reach the hardware is worse than
    a missing parameter: it produces a plausible-looking experiment whose result is meaningless,
    and you will believe the result.
\end{remark}

\subsection{Verification}

The fix was accepted only when two independent implementations agreed:

\begin{table}[H]
\centering
\begin{tabular}{lll}
\hline
\textbf{Implementation} & \textbf{Reading} & \textbf{Expected} \\
\hline
PAC example   & \codeinline{id=6a a=985,81,139 mg} & WHO\_AM\_I = 0x6A, $|a| \approx 1$ g \\
HAL example   & \codeinline{id=6a a=970,79,142 mg}  & same device, same registers \\
\hline
\end{tabular}
\caption{Two layers, one device: the numbers differ only by sensor noise and the physical
orientation at the moment of the read.}
\label{tab:i2c-imu-parity}
\end{table}

A magnitude of about 1000\,mg is not a detail to skip: a three-axis accelerometer at rest
measures gravity, so the vector length is a free self-check on both the bus and the
scale-factor arithmetic. A dead sensor reports zeros; a mis-scaled one reports something
that is not 1\,g.

\subsection{What This Teaches}

\begin{keytakeaways}
    \item \textbf{A pin map is a claim, not a measurement.} Check the electrical behaviour of
          the pins before believing a table in a repository.
    \item \textbf{Test the matrix, not the hypothesis.} Eight pairs in two orders took less
          time than a day of reasoning about which document was right.
    \item \textbf{Silence is informative.} An empty bus narrows the problem to ``no
          transaction started'', which points at wiring, mode or timing — never at the
          protocol layer.
    \item \textbf{Verify with a physical invariant.} Gravity, a chip identifier, a checksum:
          something that would be wrong if the bus were returning plausible garbage.
\end{keytakeaways}


% ============================================================================
% SECTION 07.05
% ============================================================================
\section{What To Take Away}
\label{sec:i2c-takeaways}

This chapter covered one peripheral, but three of its lessons apply to every peripheral in
this book.

\begin{keytakeaways}
    \item \textbf{The order of diagnosis is electrical, then protocol, then driver.} Prove the
          pins, prove the timing, prove the bytes \emph{in that order}. Skipping to the driver
          is how a project loses a week to a mode bit.
    \item \textbf{Controls turn opinions into statements.} ``The controller does not complete
          reads'' is an opinion until the same pins, pull-ups and device are read successfully
          by a bit-banged master. Then it is a finding.
    \item \textbf{A workaround is a decision, and decisions get written down.} The
          bit-banged transport exists because the on-chip one could not be made to work in
          any configuration that was tried. That is a legitimate engineering outcome, and it
          is only legitimate while the measurements behind it are kept.
    \item \textbf{Every abstraction must be tested against the hardware it claims to
          abstract.} The pin abstraction whose ``which bus'' parameter never reached a
          register is the canonical example: it compiled, it ran, and it measured nothing.
\end{keytakeaways}

\subsection{Exercises}

\begin{exercises}
    \item \textbf{(Easy)} Write a bus scan that reports, for every address that acknowledges,
          the byte read back from register \codeinline{0x00}. On a bus with three devices,
          how do you tell a device that answered from one that is not there?
    \item \textbf{(Easy)} Configure a GPIO as an input with no internal pull and read a pin
          that has an external pull-up, then one that does not. Compare the port values.
          What does the difference prove about the board?
    \item \textbf{(Medium)} Add a repeated-START to the bit-banged master and use it to read a
          register in a single transaction. Then implement the same read as two transactions
          and explain, with the device's datasheet, which devices tolerate the difference.
    \item \textbf{(Medium)} Calibrate \codeinline{HALF\_BIT\_CYCLES} with the DWT cycle
          counter instead of trusting the constant, and report the bus frequency you actually
          get. Compare with the value the name suggests.
    \item \textbf{(Hard)} Reproduce the fault-isolation table of this chapter for a peripheral
          of your choice: list every hypothesis, the measurement that eliminates it, and the
          control experiment that decides between ``our code'' and ``the hardware''.
\end{exercises}

\begin{notebox}
    \textbf{Where this leads.} The transport built here is what every other device driver in
    this book sits on: the OLED, the environmental sensor and the real-time clock all use the
    same bit-banged master, on their own bus, through the same \codeinline{embedded-hal}
    trait. The next chapter takes the first of them — the display — and shows what changes
    when the payload is a framebuffer rather than a register.
\end{notebox}



% Chapter 8: SPI
\chapter{Phase B: SPI Master Mode (Three Variants)}
\label{ch:chapter-08}
\chapthumb{8}
\section*{This chapter content will be added}
\clearpage

% Chapter 9: Advanced Peripherals
\chapter{Phase C: Advanced Peripherals - DMA, ADC, and Power Management}
\label{ch:chapter-09}
\chapthumb{9}
\section*{This chapter content will be added}
\clearpage

% ========== PART III: WIRELESS PROTOCOL STACKS ==========
\part{Wireless Protocol Stacks}
\label{part:wireless}

% Chapter 11: RAIL
\chapter{Phase D: Wireless Protocol Stacks - RAIL Foundation}
\label{ch:chapter-10}
\chapthumb{10}
\section*{This chapter content will be added}
\clearpage

% Chapter 12: Bluetooth (written)
% Chapter 12: Bluetooth 5.3 Low Energy
% File: chapters/chapter-12-bluetooth/chapter-12-bluetooth.tex

\chapter{Phase D2: Bluetooth 5.3 Low Energy}
\label{ch:chapter-11}
\chapthumb{11}

\begin{quote}
    \textit{``A radio stack is the one part of an embedded project you do not write.''}
\end{quote}

Bluetooth on this part is not a peripheral. It is a second processor's worth of protocol logic,
delivered as three binary libraries with a C API, and the job of a Rust project is to make that API
usable without pretending to reimplement it.

This chapter is that job, end to end: what the three blobs are, how the bindings are generated and
where they lie, the symbols nobody but you can provide, and what a beacon and a GATT peripheral
actually cost in code. Every claim here was measured on this board's SDK, and the places where it does
not work yet are marked as such.

\minitoc

\clearpage

\section{What the Stack Actually Is}
\label{sec:ble-what}

Before writing a line of binding code, three questions: what is in the binary, how do you talk to it,
and what does it expect you to do.

\subsection{Three Libraries, One Protocol}

\begin{table}[htbp]
    \centering
    \footnotesize
    \begin{tabular}{@{}p{0.34\textwidth}p{0.26\textwidth}p{0.34\textwidth}@{}}
        \toprule
        \textbf{Library} & \textbf{What it is} & \textbf{Size on this part} \\
        \midrule
        \codeinline{libbluetooth\_controller} & link layer: radio, timing, packets & ~600 KB \\
        \codeinline{libbluetooth\_host} & GAP, GATT, SM: the protocol logic & ~700 KB \\
        \codeinline{libbgcommon} & shared buffers, BGAPI parsing & ~40 KB \\
        \bottomrule
    \end{tabular}
    \caption{The three blobs of a Gecko SDK 4.x Bluetooth build. There is no fourth: the profiles and
    services you need are your code.}
    \label{tab:ble-blobs}
\end{table}

Together they are larger than most applications on this part, which is why a project chooses
Bluetooth deliberately rather than by default. All three are MSLA licensed and therefore not in this
book's repository; a clean checkout compiles the bindings and fails at link with a list of missing
symbols, which is the correct behaviour.

\subsection{BGAPI: Commands In, Events Out}

The API is asynchronous and message-based, and both halves matter:

\begin{minted}{text}
command:  sl_bt_advertiser_set_data(handle, ...)   ->  sl_status_t (immediate)
reply:    sl_bt_evt_advertiser_set_data_complete   (later, in the event queue)
\end{minted}

Every command returns a status that says ``accepted'' or ``refused'' — never ``done''. The completion
arrives as an event, and an application that does not drain the event queue eventually stops being
able to send commands, because the queue is finite and the stack does not grow it.

\begin{minted}{rust}
loop {
    if let Some(event) = ble.poll_event() {
        match event { /* Connected, Disconnected, GattWrite, ... */ }
    }
}
\end{minted}

That loop is not optional and it is not decoration: it is the other half of every call.

\subsection{What the Stack Assumes You Have Done}

The blobs are written for a vendor project, and they assume the surrounding C world exists:

\begin{itemize}
    \item \textbf{Clocks.} The radio needs the HFXO; the stack enables it, but if the crystal never
          settles the boot event never arrives.
    \item \textbf{Critical sections.} The CMSIS \codeinline{CORE\_*} helpers, which in a vendor
          project come from headers as \codeinline{static inline}.
    \item \textbf{An event loop.} Yours.
    \item \textbf{Memory.} Fixed buffers inside the blobs, sized by the stack configuration in the
          header. A Rust project cannot change them without regenerating that configuration.
    \item \textbf{Crypto} — for any connectable, bondable device, a PSA implementation the blob calls
          but does not contain.
\end{itemize}

None of that is in the API documentation, and all of it appears as a link error or a timeout on first
use. Section~\ref{sec:ble-symbols} is about the link errors; section~\ref{sec:ble-beacon} is about the
timeouts.

\section{The Bridge: bindgen Over a Vendor Header}
\label{sec:ble-bridge}

Generating Rust bindings from \codeinline{sl_bt_api.h} is one command and three surprises. All three
of this project's were compile errors in the wrapper, and all three are the kind that look like
mistakes in your own code.

\subsection{Surprise 1: The Union Is Emulated}

The header declares its event payload as an \emph{anonymous} union inside a struct. bindgen
represents anonymous unions with \codeinline{\_\_BindgenUnionField<T>}, regardless of the Rust target
requested — measured: adding \codeinline{.rust\_target(Stable\_1\_71)} left the generated file
byte-identical.

The consequence is 18 errors of the form:

\begin{minted}{text}
error[E0609]: no field `connection` on type
              `&__BindgenUnionField<sl_bt_evt_connection_opened_s>`
\end{minted}

The fix is not a bindgen flag. It is to read the payload the way the emulation requires:

\begin{minted}{rust}
// before: reads the wrapper as if it were the union
let data = &evt.data.evt_connection_opened;

// after: says out loud that the active member is the caller's responsibility
let data = unsafe { evt.data.evt_connection_opened.as_ref() };
\end{minted}

That \codeinline{unsafe} is the correct amount of friction. Reading a union is unsafe because nothing
checks that the member you chose is the one the sender wrote — and the event ID is the only thing that
tells you.

\subsection{Surprise 2: A Flexible Array Member Is Not a Slice}

\begin{minted}{text}
pub struct uint8array {
    pub len: u8,                        // "< Number of bytes stored in data"
    pub data: __IncompleteArrayField<u8>,
}
\end{minted}

\codeinline{data} has no length and no \codeinline{Index}, so \codeinline{data.value.data[i]} does
not compile. The bytes are read through the pointer, with the length taken from the header's own
field:

\begin{minted}{rust}
let byte = unsafe { data.value.data.as_ptr().add(i).read() };
\end{minted}

Two facts worth keeping: the length is authoritative (it was written by the sender), and the pointer
arithmetic is the only way to reach a flexible array — which is exactly why C APIs that return them
need a total-length parameter somewhere.

\subsection{Surprise 3: One Field, Two Names}

\begin{minted}{text}
sl_bt_evt_gatt_server_attribute_value_t      ->  attribute: u16
sl_bt_evt_gatt_server_user_read_request_t    ->  characteristic: u16
\end{minted}

The same 16-bit attribute handle, named differently in two events of the same SDK version. A wrapper
that assumes one name fails to compile on the other, and a wrapper that ``fixes'' it by renaming the
Rust type breaks the one that was right.

The project's rule: convert at the boundary, once, with a comment naming both spellings. The
application-facing type has one field, \codeinline{attribute}, and the FFI layer knows about the
vendor's inconsistency so nothing above it has to.

\subsection{What the Generated Code Is Good For}

With those three handled, the bindings are genuinely useful in the way the SVD's are: every field
carries the header's doc comment, so \codeinline{cargo doc} shows the vendor's own description of the
value you are about to write. That is the strongest argument for generating bindings rather than
hand-writing them — not the function signatures, which are mechanical, but the documentation that
travels with them.

\section{The Symbols Nobody Provides}
\label{sec:ble-symbols}

The bindings compile, the blobs are in the link path, and the linker still fails — on functions that
exist nowhere in the SDK's public surface. This section is the complete answer for this board, because
it cost a full session to assemble.

\subsection{The CMSIS Helpers}

\begin{minted}{text}
undefined symbol: CORE_EnterAtomic
  >>> referenced by bgbuf.c.obj:(bgbuf_free)
undefined symbol: CORE_ExitAtomic
  >>> referenced by ll_hci.c.obj:(ll_hciDataFromHost)
\end{minted}

In a vendor project these are \codeinline{static inline} in the CMSIS headers, so the C compiler
inlines them and no symbol is ever needed. The blob was compiled against those headers, so its object
code contains \emph{references}, and the application is expected to provide the definitions.

The definitions are four lines each, and the two that matter differ:

\begin{minted}{c}
uint32_t CORE_EnterCritical(void) { /* masks with BASEPRI */ }
uint32_t CORE_EnterAtomic(void)   { /* masks with PRIMASK */ }
\end{minted}

\textbf{This difference is not cosmetic.} \codeinline{BASEPRI} masks interrupts at or below a chosen
priority and leaves priority~0 running; \codeinline{PRIMASK} masks everything. A stack that enters a
buffer-manipulating critical section with the wrong one still links, still works at idle, and corrupts
its attribute buffer once a connection is active. Implement them from the CMSIS semantics, not from a
guess about what ``critical'' means.

\subsection{Provided Once, for Every Stack}

This board links two blob families — Bluetooth and RAIL — and both want the same helpers. There are
two ways to get that wrong:

\begin{itemize}
    \item \textbf{Compile the C file in each crate.} Two archives with the same symbols: either a
          duplicate-symbol link error, or one silently wins and the other's semantics are lost.
    \item \textbf{Give each its own copy.} Works, until the two copies drift and only one is updated.
\end{itemize}

This project's answer is a build-support crate that compiles the file once, which both stacks depend
on. A crate with no Rust API, whose entire contract is ``linking me gives you these symbols'':

\begin{minted}{toml}
[dependencies]
cmsis-stubs = { path = "../cmsis-stubs" }
\end{minted}

\subsection{The \#[link] Directive That Is Not Obvious}

Having the crate as a dependency is not enough: nothing in Rust \emph{calls} those symbols, so the
linker never pulls the archive in, and every symbol comes back undefined while the archive sits in the
search path. The missing piece is an explicit directive:

\begin{minted}{rust}
#[link(name = "cmsis_stubs", kind = "static")]
extern "C" {}
\end{minted}

Five lines that do nothing at runtime and everything at link time. They belong in the crate whose
objects need the symbols — the bindings crate — not in the helper crate.

\subsection{The Rule That Generalises}

\begin{enumerate}
    \item Read the link errors in full, once, and list every symbol. The linker names the object file
          that wanted each one (\codeinline{bgbuf.c.obj}, \codeinline{ll_hci.c.obj}), which tells you
          which layer is asking.
    \item Decide, per symbol, whether it is \emph{data you must supply} (CMSIS, emlib) or \emph{a
          library you must vendor} (PSA crypto). They have completely different fixes.
    \item Supply data once for the whole binary; vendor libraries with their licence in mind.
    \item Declare the link directive where the reference is, and comment why — because the next reader
          will see an empty \codeinline{extern} block and delete it.
\end{enumerate}

\section{A Beacon and a GATT Peripheral}
\label{sec:ble-beacon}

Two examples, two rungs. The beacon proves the stack boots and the antenna works; the GATT peripheral
proves a connection, a subscription and a notification path. Everything else in Bluetooth is a
variation on those.

\subsection{The Beacon: Nine Lines and a Phone}

\begin{minted}{rust}
let mut ble = Bluetooth::beacon(&mut delay, "xiao-mg24")?;
loop {
    if ble.poll_event().is_some() {
        led.on();  delay.delay_ms(20);  led.off();
    }
}
\end{minted}

What has to be true for a phone to see the name:

\begin{enumerate}
    \item \textbf{The stack boots.} \codeinline{Bluetooth::new} resets the controller and waits for
          \codeinline{system\_boot}. If the HFXO never settles, this times out — a board fault, not a
          radio one.
    \item \textbf{The payload fits.} Legacy advertising carries \textbf{31 bytes} per PDU. A complete
          name costs \codeinline{2 + len}; a 128-bit UUID costs 18; the flags cost 3. A name of 20
          characters and a 16-bit service UUID is already 27 of the 31.
    \item \textbf{Advertising actually starts.} \codeinline{start\_advertising} returns a status; a
          non-zero status is a refusal, and the reason is in the event queue, not in the return value.
\end{enumerate}

\subsection{The GATT Peripheral: Four Steps and Four Traps}

\begin{table}[htbp]
    \centering
    \footnotesize
    \begin{tabular}{@{}p{0.24\textwidth}p{0.34\textwidth}p{0.38\textwidth}@{}}
        \toprule
        \textbf{Step} & \textbf{What it does} & \textbf{The trap} \\
        \midrule
        advertise connectable & sets the connectable flag in the PDU & a beacon's PDU is
        non-connectable; a phone times out with no explanation \\
        wait for \codeinline{Connected} & gives you the connection handle & the handle is not the
        advertising set handle, and they are both small integers \\
        wait for a write on the CCC & the client subscribed & a notification sent before this is
        dropped by the stack, not queued \\
        notify & pushes a value & the transmit queue is finite: refusals must be counted, not ignored \\
        \bottomrule
    \end{tabular}
    \caption{Four steps, four traps. Three of the four produce no error — the value simply does not
    arrive.}
    \label{tab:ble-gatt}
\end{table}

\subsection{Counting Refusals Is the Whole Discipline}

\begin{minted}{rust}
match ble.notify(handle, &payload) {
    Ok(()) => notified += 1,
    Err(_) => refused += 1,
}
\end{minted}

A product that silently loses measurements is worse than one that reports them late. The counter is
the difference between ``the phone did not get them'' and ``the board could not send them'', and it is
the first thing to read when a client reports gaps.

\subsection{The Beacon Is the Debugging Instrument}

Before writing the GATT path, prove the beacon. It exercises the same stack, the same blob, the same
antenna and the same boot sequence, and it fails for a strictly smaller set of reasons. A project
that debugs connectable advertising and beacon advertising at the same time is debugging two things
at once with one symptom.

\section{What It Costs, and What Is Still Missing}
\label{sec:ble-costs}

A chapter about a radio stack that only reports successes would be useless. This one ends with the
bill.

\subsection{The Numbers}

\begin{table}[htbp]
    \centering
    \footnotesize
    \begin{tabular}{@{}p{0.30\textwidth}p{0.22\textwidth}p{0.40\textwidth}@{}}
        \toprule
        \textbf{Example} & \textbf{Image size} & \textbf{What it needs} \\
        \midrule
        \codeinline{ble\_beacon} & ~5.9 KB & three blobs, the stubs, a delay \\
        \codeinline{rail\_version} & ~1.3 KB & one blob, the stubs \\
        \codeinline{ble\_gatt} & does not link & the above plus PSA crypto \\
        \bottomrule
    \end{tabular}
    \caption{Measured with elf-load-range.py. The application is small; the blobs are
    not, and they are what fills flash.}
    \label{tab:ble-sizes}
\end{table}

The three Bluetooth blobs are on the order of a megabyte together. On a part with 1.5 MB of flash and
a 24 KiB bootloader, that is affordable once — and it is the reason a project chooses BLE for a
product rather than adding it ``just in case''.

\subsection{The Gap That Must Be Named: PSA Crypto}

A connectable stack references the PSA crypto API, and this SDK build does not ship it in the blobs.
The eight symbols are:

\begin{minted}{text}
psa_crypto_init      psa_destroy_key       psa_import_key    psa_mac_abort
psa_mac_sign_finish  psa_mac_sign_setup    psa_mac_update    psa_reset_key_attributes
\end{minted}

They come from Silicon Labs' \codeinline{libpsa\_crypto} and \codeinline{libmbedtls} builds, which are
separate artefacts with their own licence terms. Until they are vendored, a GATT example is a staged
file rather than a working one — and this book says so instead of shipping it as an example that
``should work''.

A project that does not need bonding can sometimes avoid the crypto path, but not by accident: the
link tells you whether it has.

\subsection{The Three Things to Verify Before Believing a BLE Build}

\begin{enumerate}
    \item \textbf{Version match.} The header and the blobs are from the same SDK. Check a symbol the
          blob exports (\codeinline{arm-none-eabi-nm libbluetooth_host*.a | grep sl\_bt\_advertiser})
          against the generated bindings.
    \item \textbf{Link, not compile.} A \codeinline{-sys} crate that compiles proves nothing about
          the blobs. Build one example that links.
    \item \textbf{One real advertisement.} A phone that sees the name is the only proof that the
          antenna, the crystal and the stack configuration agree. Everything below that is a build
          check.
\end{enumerate}

\subsection{Where This Book Stands}

\begin{minted}{text}
bindings compile                    yes, with three documented workarounds
beacon links                        yes
beacon flashed and seen by a phone  not yet -- the next measurement
GATT links                          no, PSA crypto missing
RAIL links for version only         yes; RAIL_Init needs emlib sources
\end{minted}

Stating that plainly is the point. A chapter that claimed the last four were done would be worse than
no chapter, because the next person would trust it.

\section{The API Is Not the Air}
\label{sec:ble-api-not-air}

This section could not have been written before the board was flashing, because everything in it is a
measurement that contradicts a success message. It is the part of Bluetooth bring-up that no header
documents, because the header is describing the API and the failure is between the API and the antenna.

\subsection{Four Successes, One Silence}

The first working attempt reported, on every link of the chain, that it was fine:

\begin{minted}{text}
RAIL_Init                 handle returned
RAIL_BLE_Init             void
RAIL_BLE_ConfigPhy1Mbps   status 0
RAIL_ConfigTxPower        status 0
RAIL_SetTxPowerDbm(100)   status 0
RAIL_StartTx              status 0
RAIL_EVENT_TX_PACKET_SENT 39529 times   <- counted in the radio's own interrupt
\end{minted}

and no receiver heard anything. The counter in the last line is the one that makes this worth writing
down: it is not the program counting its own calls, it is the radio's interrupt handler reporting that
the modem completed a packet. The packet was built correctly too, and read back off the chip as

\begin{minted}{text}
00 14 | 00 11 22 33 44 55 | 02 01 06 | 0A 09 "xiao-mg24"
\end{minted}

which is a well-formed \codeinline{ADV\_IND}: public address, general-discoverable flags, complete local
name. A scanner at one hundred percent duty cycle heard four other advertisers in the same room, up to
262 reports each, and this board zero times.

The API was not lying about what it was asked. It was reporting completion of a transmission that had
no power behind it.

\subsection{The Crystal Has To Be Selected, Not Merely Started}

The radio is clocked from the 39 MHz crystal. Starting it is half the job:

\begin{minted}{c}
HFXO0->CTRL |= HFXO_CTRL_FORCEEN;                       /* the oscillator runs */
CMU_ClockSelectSet(cmuClock_SYSCLK, cmuSelect_HFXO);   /* and now it is the clock */
\end{minted}

With only the first line, \codeinline{RAIL\_GetRadioClockFreqHz()} reports 19000 kHz --- the internal
oscillator --- and the BLE channel configuration is refused with
\codeinline{RAIL\_ASSERT\_INVALID\_XTAL\_FREQUENCY}. That assert is not a complaint about the crystal; it
is the library saying that the frequency it was given and the frequency it is running at are different
numbers, and it is right.

Two smaller traps live in the same paragraph. \codeinline{CMU\_OscillatorEnable} in this emlib is
documented as a dummy function kept for backward compatibility: it compiles, it is callable, and it
does nothing. And \codeinline{CMU\_ClockFreqGet(cmuClock\_HFXO)} returns zero because the function's
switch has no case for that clock, so a zero there says nothing at all about the oscillator.

\subsection{The Amplifier Needs Its Curves, and the Call Is One Line}

The gap between the modem and the antenna was a missing initialisation, and the library documents it in
a comment inside its own conversion function: without \codeinline{RAIL\_InitTxPowerCurves} the
dBm-to-raw conversion finds no curve and returns 255, its ``invalid'' value. Silicon Labs' entry point
for that is one call:

\begin{minted}{c}
sl_rail_util_pa_init();                  /* loads the PA curves into RAIL */
RAIL_ConfigTxPower(handle, sl_rail_util_pa_get_tx_power_config_2p4ghz());
RAIL_SetTxPowerDbm(handle, 100);
\end{minted}

The measurement that identifies this failure is a read-back, not a status:

\begin{minted}{text}
                            before                      after sl_rail_util_pa_init()
RAIL_SetTxPowerDbm(100)     0 (success)                 0 (success)
RAIL_GetTxPowerDbm()        -32768  (invalid power)     99   (10 dBm)
RAIL_GetTxPower()           0       (no level)          0x24
\end{minted}

A radio with no valid amplifier level completes every packet it is asked to send. The success in the
first row is true and worthless; the two read-backs are the only rows that describe the air.

\subsection{The Vector Table Had No Radio In It}

Before any of that could be measured, the first transmission faulted: \codeinline{HardFault},
\codeinline{CFSR=0x00000100} (imprecise bus error), at \codeinline{RAIL\_StartTx}. The cause was in the
PAC, not in the radio: it declared 46 of the device's 74 interrupts, and every radio one was among the
missing --- \codeinline{AGC}, \codeinline{BUFC}, \codeinline{FRC\_PRI}, \codeinline{FRC},
\codeinline{MODEM}, \codeinline{PROTIMER}, \codeinline{RAC\_RSM}, \codeinline{RAC\_SEQ},
\codeinline{SYNTH}, \codeinline{RFECA0/1}. Those slots held zeroes in the vector table, so the first
radio interrupt jumped to address zero.

Two further layers hide the same defect, and both were measured:

\begin{itemize}
    \item The vendor names its handlers \codeinline{<NAME>\_IRQHandler}. A vector entry saying
          \codeinline{FRC} falls short of the blob's \codeinline{FRC\_IRQHandler} by the suffix.
    \item A \codeinline{PROVIDE(FRC = DefaultHandler)} \emph{satisfies} the vector's reference, so the
          linker never pulls the archive member that defines the real handler. With the aliases in
          place the program halted inside \codeinline{DefaultHandler}. Linking the RAIL blob with
          \codeinline{+whole-archive} keeps both: a PAC example with no radio still links, and a binary
          with the radio gets RAIL's handlers.
\end{itemize}

\subsection{The Instrument Matters More Than the Hypothesis}

Two instruments were needed and neither was the obvious one.

\codeinline{hcitool lescan} and \codeinline{bluetoothctl} printed nothing on this host's adapter,
including while the host itself was advertising, so both would have been read as ``nothing on air''.
The scan that works is set up through raw HCI with the window equal to the interval --- one hundred
percent duty cycle --- and the reports are read back out of a \codeinline{btmon} capture, which sees
other advertisers the command-line tools never printed.

\codeinline{RAIL\_GetRadioState} returning idle after each transmission proves only that the state
machine moved. The value that identified the fault was the power read-back; the value that proved the
modem was working all along was the send-complete event, counted in an interrupt.

\subsection{What the Working Sequence Is}

\begin{minted}{c}
emlib_hfxo_start(); emlib_clock_select_hfxo();
RAIL_Init(&cfg, NULL);                        /* with an events callback */
RAIL_BLE_Init(handle);
RAIL_BLE_ConfigPhy1MbpsViterbi(handle);       /* only the Viterbi 1M config exists here */
sl_rail_util_pa_init();                       /* the PA curves: without this, silence */
RAIL_ConfigTxPower(handle, sl_rail_util_pa_get_tx_power_config_2p4ghz());
RAIL_SetTxPowerDbm(handle, 100);
RAIL_ConfigData(&packetMode); RAIL_SetTxFifo(txBuf, 0, 64);
for (;;) {
  channel = 37 + (n % 3);                     /* the three advertising channels */
  RAIL_BLE_ConfigChannelRadioParams(handle, 0x555555, 0x8E89BED6, channel, false);
  RAIL_PrepareChannel(handle, channel);
  RAIL_WriteTxFifo(pdu, pduLen, true);        /* the full PDU, header included */
  RAIL_StartTx(handle, channel, RAIL_TX_OPTIONS_DEFAULT, NULL);
}
\end{minted}

The radio parameters are re-applied per channel because the whitening seed is the channel number.
\end{document} is not needed; the loop is the whole application.

With that sequence the board advertises, and the proof is a capture from another machine:

\begin{minted}{text}
Address: 55:44:33:22:11:00 (OUI 55-44-33)
Flags: 0x06   LE General Discoverable Mode, BR/EDR Not Supported
Name (complete): xiao-mg24
RSSI: -83 dBm
\end{minted}

\subsection{The Rule This Chapter Adds to the Book}

A radio that reports success is not a radio that radiates. Every link in this chain --- the channel
configuration, the transmission, the modem, the interrupt count --- can be green while the antenna is
silent, and the only measurement that distinguishes them is one taken \emph{outside the chip}: a
receiver, and a read-back of the amplifier's own idea of its power. The build was correct for a long
time before it was true.


\clearpage

\begin{keytakeaways}
    \item A Bluetooth stack arrives as binary blobs plus a C API. The Rust work is bindings, symbols
          and a safe surface — not a protocol implementation.
    \item The header and the blobs must be the same SDK version. A mismatch produces calls that exist
          in the blob and not in the bindings, which reads like a typo and is not.
    \item The blobs reference \codeinline{CORE\_*} helpers that the application must provide, and they
          must be provided once for the whole binary, not once per stack.
    \item Legacy advertising carries 31 bytes; a name plus a service UUID uses more than half of it,
          and a builder that refuses an oversized payload is better than one that truncates.
    \item A connectable stack pulls the bonding path, which needs a crypto library. That is a license
          and a link problem, not an API problem.
    \item A radio can report every transmission complete and radiate nothing: the amplifier needs its
          power curves loaded, and the API's own success says nothing about the air.
    \item RAIL's channel argument is the physical channel, not the logical one, and the advertising
          channels are (0,37), (12,38), (39,39).
\end{keytakeaways}


% Chapter 13: Zigbee
\chapter{Phase D3: Zigbee 3.0 Protocol Stack}
\label{ch:chapter-12}
\chapthumb{12}
\section*{This chapter content will be added}
\clearpage

% Chapter 14: Thread and Matter
\chapter{Phase D4: Thread Protocol and Matter Integration}
\label{ch:chapter-13}
\chapthumb{13}
\section*{This chapter content will be added}
\clearpage

% ========== PART IV: TESTING, BSP, AND PRODUCTION ==========
\part{Testing, BSP, and Production}
\label{part:testing-bsp-production}

% Chapter 14: Testing (physically 14th: Part III added four chapters after this was planned; written: the fall detector, nine real drops, and an instrument that lied)
% Chapter 10: Test-Driven Development and Quality Assurance
% File: chapters/chapter-14-testing/chapter-14-testing.tex

\chapter[Test-Driven Development]{Test-Driven Development and Quality Assurance: Nine Real Drops}
\label{ch:chapter-14}
\chapthumb{10}

\begin{quote}
    \textit{``A test that passes while the thing it tests is broken is worse than no test at
    all: it converts an open defect into a certified one.''}
    \\
    -- the rule this chapter was written to satisfy
\end{quote}

Embedded software is usually tested in two places: on a host, where the compiler is friendly
and the hardware is absent, and on the target, where the opposite is true. Most projects do
only the second, call it ``testing on hardware'', and never learn that the numbers they
trusted were produced by an instrument that was lying.

This chapter follows one feature from first line to validated behaviour: a fall detector for
a wrist-worn or carried device, built on a six-axis IMU. It is a good case study because it
cannot be judged by inspection — a fall is an event in time, and the only way to know whether
you detect it is to drop the thing and look. Along the way, three of the measurements used to
validate it turn out to be wrong in ways that would have certified a defect as fixed.

\minitoc

\clearpage

\section{What Can Be Tested Off Target}
\label{sec:off-target-testing}

The first design decision of a testable embedded feature is not a test at all: it is where
the code lives.

\subsection{Split the Policy From the Peripheral}

A fall detector is two things wearing one name:

\begin{itemize}
    \item \textbf{Acquisition} — power the IMU, configure its scale and output rate, read a
          12-byte burst, convert to milligees. This only exists on the target.
    \item \textbf{Classification} — decide, from a stream of magnitudes and the time between
          them, that what just happened was a fall. This is arithmetic and state, and it does
          not know that a sensor exists.
\end{itemize}

If both live in the driver crate, the second becomes untestable: building the tests drags in
a PAC, a linker script and a target triple, so nobody writes them. This project keeps the
classifier in its own crate with \textbf{no dependencies at all}:

\begin{minted}{toml}
[package]
name = "motion-detection"
version = "0.1.0"
edition = "2021"

[dependencies]
# deliberately empty: no PAC, no HAL, no panic handler
\end{minted}

\begin{tipbox}
    A dependency-free classifier crate is not an aesthetic choice. It is what makes
    \codeinline{cargo test} work on the developer's machine in under a second, which is what
    makes writing the tenth test cheap enough to bother.
\end{tipbox}

\subsection{The API That Makes Testing Possible}

The crate takes samples and returns events, with the sample interval as a parameter rather
than something it looks up:

\begin{minted}{rust}
pub struct FallDetector { /* phase, peak, timers */ }

impl FallDetector {
    pub const fn new() -> Self;
    /// Feed one sample; the interval is microseconds.
    pub fn update_us(&mut self, accel_mg: [i32; 3], dt_us: u32) -> Option<Event>;
    /// True while the device is believed to be in the air.
    pub fn in_free_fall(&self) -> bool;
}

pub enum Event {
    FreeFall { ms: u32 },
    Impact { peak_mg: i32 },
    Fall { peak_mg: i32, freefall_ms: u32 },
}
\end{minted}

Two details are worth copying into any similar design:

\begin{enumerate}
    \item \textbf{Time is an argument.} A detector that calls a clock cannot be fed a
          recording at ten times real speed, and cannot be tested for how it behaves when a
          sample arrives late. Passing the interval makes both possible and costs nothing.
    \item \textbf{Microseconds, not milliseconds.} At 400\,Hz a sample is 2.5\,ms; rounding to
          whole milliseconds is a 20\,\% error in every timeout. This is not a theoretical
          concern — see Section~\ref{sec:instrument-lied}.
\end{enumerate}

\subsection{The Shape of the Whole Test Suite}

Nine host tests cover the classifier. They are not nine variations of the same case; each one
exists because a specific wrong implementation would pass the others:

\begin{table}[H]
\centering
\begin{tabular}{p{6.6cm}p{7.4cm}}
\hline
\textbf{Test} & \textbf{What it would catch} \\
\hline
\codeinline{quiet_board_reports_nothing} & an implementation that reports on any movement \\
\codeinline{a_dip_too_short_is_not_a_fall} & counting a single low sample as free fall \\
\codeinline{impact_without_free_fall_is_reported_as_impact} & calling every knock a fall \\
\codeinline{a_fall_caught_before_landing_is_not_reported} & ignoring the landing entirely \\
\codeinline{handling_is_not_a_fall} & thresholds tuned only to the positive case \\
\codeinline{real_drop_with_impact_after_a_one_g_interlude_is_detected} &
  requiring the impact at the transition \\
\codeinline{real_drop_with_a_long_interlude_is_detected} & the same, with a longer gap \\
\codeinline{a_soft_landing_after_a_short_free_fall_is_detected} &
  a settle rule that is too strict \\
\codeinline{isqrt_is_exact} & square roots that drift and shift every threshold \\
\hline
\end{tabular}
\caption{The classifier's host tests, each justified by the wrong implementation it rejects.}
\label{tab:classifier-tests}
\end{table}

\begin{remark}
    \textbf{Every one of those tests was added after a real failure, not before.} The suite
    started with five and grew as the recordings exposed assumptions: the interlude tests came
    from two drops that were missed, the soft-landing test from a third. That is normal. What
    is not acceptable is deleting a failing test because the new code is ``better''.
\end{remark}

\section{Tests That Pin a Threshold to a Measurement}
\label{sec:tests-pin-thresholds}

A threshold with no test is a guess that compiles. The discipline that changes it is simple:
\textbf{every constant is justified by a number, and every number has a test on the other side
of it.}

\subsection{The Near-Miss Rule}

Any rule that suppresses or classifies something needs a case that comes \emph{close} and
must still behave correctly. For a fall detector the near misses are the interesting ones:

\begin{table}[H]
\centering
\begin{tabular}{llll}
\hline
\textbf{Boundary} & \textbf{Value} & \textbf{Must trigger} & \textbf{Must not trigger} \\
\hline
free fall & $<$ 400\,mg & real drops: 32--50\,mg & shaking: never below 630\,mg \\
free-fall length & $\geq$ 40\,ms & 42, 70, 80, 100, 138, 145, 176--252\,ms & dips of 10--25\,ms \\
landing & $\geq$ 1100\,mg & 1148--12823\,mg & caught in the hand: $\approx$ 990\,mg \\
stillness & $<$ 1500\,mg for 600\,ms & 1.04--1.17\,s observed & bouncing: $>$1500\,mg \\
\hline
\end{tabular}
\caption{Every constant, the measurements on both sides of it, and the control that must not fire.}
\label{tab:thresholds-near-miss}
\end{table}

Two of these margins are uncomfortably thin, and saying so is part of the test:

\begin{warningbox}
    The landing threshold is 1100\,mg with the softest real landing at 1148\,mg — an 8\,\%
    margin, on nine drops, performed by one person on one desk. It is enough to ship a rule
    and not enough to claim the rule is right. The test encodes the measurements, not a
    belief about all landings everywhere.
\end{warningbox}

\subsection{Writing the Test From the Recording, Not From the Idea}

The failing pattern to avoid is a test that restates the implementation:

\begin{minted}{rust}
// Useless: this asserts that the code does what the code does.
#[test]
fn detects_a_fall() {
    let mut det = FallDetector::new();
    assert!(det.update([0, 0, 0], 5).is_some());   // whatever the code says
}
\end{minted}

The useful version uses a sequence of samples taken from a recording, and asserts on the
number a person measured:

\begin{minted}{rust}
#[test]
fn a_soft_landing_after_a_short_free_fall_is_detected() {
    // Drop 3 of the 400 Hz round: 42 ms of free fall, landing of only 1493 mg,
    // then quiet. This is the case an 800 ms settle rule lost.
    let mut det = FallDetector::new();
    feed(&mut det, [990, 40, 60], 300);      // resting
    feed(&mut det, [30, 25, 20], 42);        // in the air
    feed(&mut det, [1200, 700, 400], 20);    // 1493 mg landing
    match feed(&mut det, [980, 40, 60], 700) {
        Some(Event::Fall { peak_mg, .. }) => {
            assert!((1200..1700).contains(&peak_mg), "peak {peak_mg}");
        }
        other => panic!("expected a fall, got {other:?}"),
    }
}
\end{minted}

Three properties make this test worth its lines: the numbers come from hardware, the
assertion range is wider than the measured value (sensor noise is real), and it fails loudly
with the actual event when the behaviour changes.

\subsection{Embed the Recording in the Image}

Host tests prove the logic. They do not prove that the code on the device behaves the same,
because the device may run a different build, a different clock, a different optimisation
level. The cheap bridge between the two is to embed the recordings in the firmware and replay
them at boot:

\begin{minted}{text}
=== self test at boot, no host attached ===
real=clean  synth=ok  200Hz=2/2  400Hz=4/4
\end{minted}

That line is the whole argument: the handling recording produces nothing, the synthetic fall
is detected, and six of the real drops — replayed through the same binary that runs live — are
all found. A regression in the shipped image shows up as a different line, on a device that
needs no laptop.

\begin{tipbox}
    Embed the recordings as \codeinline{const} arrays (a few kilobytes) and replay them
    \emph{before} the main loop. It costs one flash write and turns every power-up into a
    regression test.
\end{tipbox}

\section{Nine Real Drops}
\label{sec:nine-real-drops}

The detector was validated by dropping the board, by hand, and reading what it recorded. Not
because that is rigorous — it is not, and the chapter says so below — but because it is the
only measurement that can distinguish a plausible rule from a working one.

\subsection{The Protocol}

Three rounds, each recorded by the board itself, plus one control:

\begin{enumerate}
    \item Configure the IMU for a range that will not clip ($\pm$8\,g) and the fastest output
          rate it offers (416\,Hz).
    \item Sample as fast as the loop allows, and record every sample \emph{with its interval}.
    \item Drop the board from about 40\,cm, several times, two seconds apart.
    \item Then perform the negative case: shake hard, set it down suddenly, shake again.
    \item Read the log over the debug probe, with a script, and analyse it on the host.
\end{enumerate}

\begin{notebox}
    \textbf{The control is not optional.} A detector that fires on everything detects every
    fall. Every round of drops in this chapter was recorded together with a handling round
    that had to stay silent, and the thresholds were chosen to satisfy both at once.
\end{notebox}

\subsection{What the Recordings Showed}

\begin{table}[H]
\centering
\begin{tabular}{llll}
\hline
\textbf{Round} & \textbf{Rate} & \textbf{Drops} & \textbf{Detected} \\
\hline
1 & 200\,Hz & 2 & 2 \\
2 & 200\,Hz & 5 & 3 (two lost their free fall between samples) \\
3 & 400\,Hz & 4 onto a bed & 4 \\
control & 200\,Hz & handling & 0 false alarms \\
\hline
\end{tabular}
\caption{Three rounds of drops and one control, all recorded on the board.}
\label{tab:drop-rounds}
\end{table}

Round 3 is the one worth reading in full, because it is the round that corrected the design:

\begin{table}[H]
\centering
\begin{tabular}{llll}
\hline
\textbf{Free fall} & \textbf{Landing} & \textbf{Quiet before the verdict} & \textbf{Verdict} \\
\hline
176\,ms & 11699\,mg & 1.04\,s & fall \\
223\,ms & 13840\,mg & 1.05\,s & fall \\
248\,ms & 13149\,mg & 1.14\,s & fall \\
252\,ms & 13742\,mg & 1.17\,s & fall \\
\hline
\end{tabular}
\caption{Four drops onto a bed, measured. The landings are three to four times harder than
the prediction that a mattress would absorb them.}
\label{tab:bed-round}
\end{table}

\subsection{The Three Assumptions the Data Destroyed}

\paragraph{``A fall is a big impact.''} Real landings measured 1148 to 12823\,mg. Handling
peaked at 1801\,mg. The ranges overlap: a board laid down flat lands more gently than a board
shaken hard, so \textbf{no impact threshold can separate a fall from a knock}. What separates
them is the free-fall phase. The impact threshold only has to distinguish ``landed'' from
``returned to one g because it was caught''.

\paragraph{``The impact follows the free fall immediately.''} In the first recordings it came
120\,ms later, and in the second 400\,ms later. A state machine that required the landing at
the transition detected \textbf{zero of two} drops. Opening a 600\,ms window took the same data
to two of two.

\paragraph{``A bed is soft.''} Predicted: a mattress landing would fall below the 1100\,mg
threshold and be missed. Measured: 11.7 to 13.8\,g. The prediction was wrong, the rule was
right, and the honest conclusion is that nine drops on one bed say nothing about mattresses
in general.

\begin{warningbox}
    \textbf{What this validation is not.} Nine drops, one board, one desk, one bed, one
    person. It proves the detector works on this setup and that the thresholds are not
    arbitrary. It does not prove a fall-detection rate for a product, and a claim that it did
    would be the kind of overreach that gets someone hurt.
\end{warningbox}

\subsection{What Is Still Missing}

Two things are absent from this validation and both are visible in the code:

\begin{itemize}
    \item \textbf{The gyroscope is read and ignored.} Rotation would separate ``dropped while
          flat'' from ``tipped over on a table'', and the data is already in the same burst.
    \item \textbf{A person falling is not a board being dropped.} A carried device falls with
          its owner: shorter free fall, cushioned landing, and a body that moves afterwards.
          The classifier is tuned to the harsher case for the impact threshold and the easier
          one for the free-fall phase, which is a defensible starting point and not an answer.
\end{itemize}

\section{The Instrument That Lied}
\label{sec:instrument-lied}

Everything in the previous section was, at one point, wrong in the same way: the measurements
were taken on a time axis that did not exist. This section is the most important one in the
chapter, because the failure was silent, it survived a code review, and it certified a defect
as fixed.

\subsection{The Sampler That Claimed 400\,Hz and Ran at 158}

The loop was written to sample at 400\,Hz and handed the detector a constant:

\begin{minted}{rust}
const SAMPLE_PERIOD_CYCLES: u32 = 47_500;  // 400 Hz at 19 MHz
const DT_US: u32 = 2_500;                  // 400 Hz, in microseconds
...
application.sample_us(motion.accel_mg, DT_US);
\end{minted}

The loop was paced correctly — it never sampled \emph{faster} than 400\,Hz. But each iteration
also did work: a 12-byte I2C burst, a detector update, a record write, and periodically a
1090-byte display flush over the same bit-banged bus. The true interval was set by the
\textbf{slowest} part of the body, not by the period:

\begin{minted}{text}
measured: 317 samples in 2 s  ->  158 Hz, mean interval 6.3 ms
\end{minted}

Every timeout in the detector was therefore 2.5 times too short, and every duration in the
first analysis — ``40 ms of free fall'' — was 2.5 times too small. The measurements were
self-consistent, which is exactly why nobody noticed.

\subsection{What Gave It Away}

Not a test. A pattern in the output:

\begin{minted}{text}
AIRBORNE  free 85ms   ->  FALL   after 0.62 s
AIRBORNE  free 52ms   ->  FALL   after 0.62 s
AIRBORNE  free 120ms  ->  FALL   after 0.62 s
AIRBORNE  free 40ms   ->  FALL   after 0.53 s
\end{minted}

Four different drops, four different free falls, and the settle time is the same to the
millisecond. \textbf{The number was describing the code, not the world.} That is the signature
of a broken instrument: when an output is constant across inputs that should vary, you are
reading your own parameter back.

\subsection{Three Fixes, and One of Them Was Deleted by the Compiler}

\begin{enumerate}
    \item \textbf{Measure the interval.} The loop takes the difference of the DWT cycle
          counter between samples and passes that everywhere — detector, recorder, alarm.
          \begin{minted}{rust}
let dt_us = now.wrapping_sub(last_sample_cycles) / CYCLES_PER_US;
last_sample_cycles = now;
application.sample_us(motion.accel_mg, dt_us);
          \end{minted}
    \item \textbf{Record the interval with every sample.} A dataset whose timing is unknown
          can only be analysed dishonestly. One byte per sample, in 250\,$\mu$s units.
    \item \textbf{Get the display out of the critical path.} Refresh every 800 samples instead
          of 80, and never while the device is airborne or within 1.2\,s of an event: a flush
          in the middle of a fall stalls the sampler exactly when the samples matter. Measured
          effect: 158\,Hz to 235\,Hz without touching the sensor.
\end{enumerate}

And then the second-order failure, which is the reason this section exists:

\begin{warningbox}
    \textbf{The compiler deleted the per-sample timing array.} The recording was written to
    ``carry its own timing'', and nothing in the crate ever read that array — so the optimiser
    removed it from the image entirely. The code compiled, the comment claimed a property, the
    symbol was absent from the binary, and the recording carried nothing. It was found by
    running \codeinline{arm-none-eabi-nm} and looking for the symbol instead of trusting the
    comment. The writes are volatile now, and the check is part of the habit: \emph{a property
    you cannot see in the artefact is a property you do not have.}
\end{warningbox}

\subsection{The Instrument Rules}

\begin{keytakeaways}
    \item \textbf{Never pass the interval you asked for; pass the one you measured.} The
          intended rate is a wish. The DWT counter is the fact.
    \item \textbf{A constant output across varying inputs means you are measuring yourself.}
          The 0.62\,s settle was a parameter echoed back.
    \item \textbf{A recording must carry its timing.} Otherwise every later analysis inherits
          an assumption nobody remembers making.
    \item \textbf{Verify the artefact, not the intent.} Check that the symbol is in the binary,
          that the counter advanced, that the file grew. Comments describe intentions; the
          binary is what runs.
    \item \textbf{Calibrate the instrument before calibrating the product.} The rate, the
          units and the resolution have to be right before any threshold on top of them means
          anything.
\end{keytakeaways}

\section{What To Take Away}
\label{sec:testing-takeaways}

\begin{keytakeaways}
    \item \textbf{Put the policy in a crate with no dependencies.} If the classifier needs a
          PAC to compile, its tests will not be written.
    \item \textbf{Pass time in as an argument.} It makes fast replay, late samples and
          deterministic tests free.
    \item \textbf{Every threshold is a measurement with a test on both sides.} A rule without
          a near-miss case is a guess with a comment.
    \item \textbf{Embed the recordings and replay them at boot.} Every power-up becomes a
          regression test on the exact binary that ships.
    \item \textbf{Calibrate the instrument first.} A wrong sample rate invalidates every
          number downstream, silently and consistently.
    \item \textbf{Say what the validation is not.} Nine drops prove a rule is not arbitrary;
          they do not prove a detection rate for a product.
\end{keytakeaways}

\subsection{Exercises}

\begin{exercises}
    \item \textbf{(Easy)} Take the classifier out of the driver crate in a project of your own
          and make its tests run with \codeinline{cargo test} on the host. Report how long the
          suite takes before and after.
    \item \textbf{(Easy)} Add a test for the boundary on each side of one threshold: a sample
          exactly at the limit, and one a single unit below it. Which side does your
          implementation treat as inclusive, and is that what the documentation says?
    \item \textbf{(Medium)} Instrument your own sampling loop with \codeinline{DWT->CYCCNT} and
          report the real rate and the worst-case interval. Compare with the rate you intended.
    \item \textbf{(Medium)} Add the gyroscope to the classifier: compute the rotation during the
          free fall and use it to separate ``dropped flat'' from ``tipped over''. What
          additional recording do you need to set the threshold, and can you get it without a
          second person?
    \item \textbf{(Hard)} Design the validation for a detector of a different event of your
          choice (a door being forced, a machine starting, a container being opened) following
          the protocol of this chapter: rounds, a control that must stay silent, near-miss
          cases, and an explicit statement of what the sample does not prove.
\end{exercises}

\begin{notebox}
    \textbf{Where this leads.} The next chapter takes the same evidence-driven approach to the
    layer above the drivers: the board support package, where the question stops being ``does
    this device work'' and becomes ``does this board work, and how do I prove that nobody has
    to know which pin it is on''.
\end{notebox}



% Chapter 15: BSP (written)
% Chapter 15: BSP
% File: chapters/chapter-15-bsp/chapter-15-bsp.tex

\chapter{Board Support Packages and Hardware Abstraction}
\label{ch:chapter-15}
\chapthumb{15}

\begin{quote}
    \textit{``A BSP is the layer that knows which wire is which — and it is the only layer allowed
    to.''}
\end{quote}

The PAC knows registers. The HAL knows protocols. Neither knows that a particular board puts an LED on
\codeinline{PA7} with the cathode to the pin, and that is the BSP's job: the last layer between the
user's intent and the silicon, and the only one that should contain a pin number.

This chapter builds a BSP for hardware that arrives in two pieces — a module and a carrier — which is
the arrangement that exposes every mistake in the design.

\minitoc

\clearpage

\section{Two Pieces of Hardware, Not One}
\label{sec:bsp-two-pieces}

This board is not one board. It is a module — the XIAO MG24 Sense, which runs on its own — plus a
carrier that the module plugs into. The distinction is physical, it decides who owns which feature, and
a single BSP crate erases it.

\subsection{What Each Piece Owns}

\begin{table}[htbp]
    \centering
    \footnotesize
    \begin{tabular}{@{}p{0.22\textwidth}p{0.30\textwidth}p{0.40\textwidth}@{}}
        \toprule
        \textbf{Piece} & \textbf{Components} & \textbf{Can be used alone?} \\
        \midrule
        XIAO MG24 & LED \codeinline{PA7}, USB, radio & yes \\
        XIAO MG24 \emph{Sense} & the above + IMU \codeinline{LSM6DS3TR-C}, analog mic & yes \\
        Expansion Base & OLED, BME68x, RTC, buzzer, button, Grove, microSD & no; needs a XIAO \\
        \bottomrule
    \end{tabular}
    \caption{Three hardware identities and one board on the desk. The Sense is a variant of the module,
    not a different module.}
    \label{tab:bsp-pieces}
\end{table}

\subsection{The Same PCB With More Parts}

The Sense is a XIAO with two extra components fitted. That single fact decides the crate layout:

\begin{itemize}
    \item \textbf{A feature, not a crate.} Two crates would each need the boot sequence, the clock
          tree, the LED and the pin map — four things to keep in step, for two products that are
          physically the same PCB.
    \item \textbf{A feature that is opt-in.} A user with a plain XIAO must not compile the IMU driver
          or its bus. \codeinline{--features sense} adds exactly those.
    \item \textbf{The feature names the thing that is fitted,} not the price bracket or the marketing
          name: \codeinline{sense} means ``the two sensors are populated''.
\end{itemize}

\begin{minted}{toml}
[features]
default = []
# The Sense variant is the same PCB with more components fitted: the IMU and the
# analog microphone. Same board, fewer parts -- so a feature, not a crate.
sense = []
\end{minted}

\subsection{The Carrier Is a Different Crate, and It Points One Way}

The Expansion Base adds devices the module has never heard of, and a user may own neither, one or both.
So it is its own crate, and the dependency has a direction:

\begin{minted}{text}
xiao-expansion-base-bsp  ──depends on──▶  xiao-mg24-bsp  ──depends on──▶  efr32mg24-hal
\end{minted}

The opposite arrow — the module crate knowing that a carrier exists — is the defect this book's
section~\ref{sec:layer-boundaries} describes: it makes the module unusable without the carrier's code,
and it puts carrier pins in a crate whose whole claim is ``the XIAO on its own''.

\subsection{What the Direction Buys}

\begin{itemize}
    \item A user with only the XIAO compiles one small crate, not the carrier's drivers.
    \item Carrier examples live where the carrier is, so the path of a file tells you the hardware it
          needs before you open it.
    \item The module crate cannot drift into knowing the carrier, because the manifest forbids it —
          and \codeinline{cargo} enforces manifests even when reviewers do not.
\end{itemize}

\section{One Crate per Board, and the Arrow Between Them}
\label{sec:bsp-one-crate}

The crate layout is not a matter of taste. It follows from one question: \emph{which of these things
can be owned separately?}

\subsection{The Two Schools}

\begin{table}[htbp]
    \centering
    \footnotesize
    \begin{tabular}{@{}p{0.24\textwidth}p{0.32\textwidth}p{0.36\textwidth}@{}}
        \toprule
        \textbf{Project} & \textbf{Layout} & \textbf{Where the wiring lives} \\
        \midrule
        \codeinline{stm32f4xx-hal} & one HAL per family, no BSP & in every example \\
        \codeinline{esp-hal} & one HAL, no BSP, chip by feature & in every example \\
        \codeinline{rp-hal} / \codeinline{rp-hal-boards} & HAL per chip + a BSP crate per board & in the BSP crate \\
        \bottomrule
    \end{tabular}
    \caption{Two schools. The difference is not sophistication: it is whether the board's wiring is a
    fact worth writing down once.}
    \label{tab:bsp-schools}
\end{table}

Both schools are legitimate. The first optimises for a HAL whose users all have a schematic in front of
them; the second optimises for a board whose wiring was expensive to establish. This project is firmly
in the second camp, for a reason that is in chapter~4's table of contradictions: the pin map of this
board was measured, and a measurement that is repeated in every example is a measurement that will be
repeated wrongly.

\subsection{The Four Crate Questions}

\begin{enumerate}
    \item \textbf{How many chips?} One HAL per chip or per family. This part gets one.
    \item \textbf{Do the boards stack?} If a carrier holds a module, one crate per piece, arrow carrier
          $\rightarrow$ module.
    \item \textbf{Are there variants?} A variant of the \emph{same} PCB is a feature. A different PCB is
          a crate.
    \item \textbf{Was the wiring obvious or was it conquered?} Obvious wiring can live in examples;
          conquered wiring belongs in a BSP with the measurement next to it.
\end{enumerate}

\subsection{What Ends Up in the Workspace}

\begin{minted}{text}
crates/
  efr32mg24-pac/            registers
  efr32mg24-hal/            drivers + pin and bus primitives
  xiao-mg24-bsp/            Module: led(), and with --features sense: imu(), microphone()
  xiao-expansion-base-bsp/  ExpansionBase: oled(), environment(), rtc(), buzzer(), button()
libraries/
  motion-detection/         domain policy: no hardware at all
\end{minted}

Two rules hold that shape together, and both are checks in this project's gate:

\begin{itemize}
    \item \textbf{Nothing in libraries/ may depend on the hardware stack.} A classifier
          that needs a HAL is a classifier that cannot be tested on a laptop, and the test suite is the
          only thing that makes it trustworthy.
    \item \textbf{Every crate that exposes hardware carries at least one example of its own level.} An
          API with no examples compiles, lints, and rots — this project had a whole module do exactly
          that for weeks.
\end{itemize}

\subsection{The Defect the Layout Prevented}

Before the split, the HAL itself contained \codeinline{board.rs} with the carrier's wiring — the OLED
bus, the button pin, the buzzer channel — because it was convenient. Two consequences followed
immediately:

\begin{enumerate}
    \item The HAL became unusable without the carrier's assumptions, so a user with a bare XIAO had to
          read carrier pins to blink an LED.
    \item Two BSPs existed at once: the HAL's board module and the new facade. Half the examples proved
          one, half the other, and the older one stopped compiling without anyone noticing.
\end{enumerate}

The fix was not better documentation. It was deleting the wiring from the HAL and making the
dependency direction a manifest fact.

\section{What a BSP Should Expose}
\label{sec:bsp-api}

A BSP's public API is a vocabulary, and the test of a vocabulary is what a user can say with it without
knowing anything else.

\subsection{The Right Shape}

\begin{minted}{rust}
let module = Module::take().expect("peripherals taken twice");
let mut led = module.led();

let base = ExpansionBase::take().expect("peripherals taken twice");
let mut oled = base.oled();
let env = base.environment()?;      // BME68x
let rtc = base.rtc()?;              // PCF8563
let mut buzzer = base.buzzer();
\end{minted}

Every one of those names is a thing on the desk. None of them is an address, a bus, a pin, an alias or
a register — and that is the whole specification.

\subsection{What the Signature Hides, and What It Must Not}

\begin{table}[htbp]
    \centering
    \footnotesize
    \begin{tabular}{@{}p{0.26\textwidth}p{0.30\textwidth}p{0.40\textwidth}@{}}
        \toprule
        \textbf{Hidden} & \textbf{Example} & \textbf{Why hiding it is right} \\
        \midrule
        addresses & \codeinline{0x3C}, \codeinline{0x76}, \codeinline{0x51} & a user cannot check them and must not need to \\
        pins & \codeinline{PC04}/\codeinline{PC05}, \codeinline{PB2}/\codeinline{PB3} & measured facts; wrong in half the documents \\
        power gates & \codeinline{PD5}, \codeinline{PC8} & sequencing the user would get wrong \\
        aliases & which peripheral window answered & a security-attribution detail \\
        baud and clock & I2C half-period, PWM frequency & driver knowledge \\
        \bottomrule
    \end{tabular}
    \caption{What a facade should hide. Note that every row is a \emph{fact}, not a \emph{choice} — the
    user has no decision to make about any of them.}
    \label{tab:bsp-hides}
\end{table}

\subsection{What It Must Not Hide: Structure}

The temptation is to hide the two buses, because ``two I2C buses'' is an implementation detail. It is
not — it is a fact about the hardware that changes what the user can do:

\begin{itemize}
    \item The panel, the BME68x and the RTC share \codeinline{PC04}/\codeinline{PC05}.
    \item The IMU has its own pair, \codeinline{PB2}/\codeinline{PB3}, because the module's sensors are
          wired to the module's pins.
    \item A single \codeinline{i2c()} would have to take a ``which bus'' argument anyway, so hiding it
          only moves the choice to a place with less information.
\end{itemize}

The rule that came out of this: \textbf{hide the addresses, keep the structure}. A facade that hides
structure has to guess, and guessing is what the layer below exists to prevent.

\subsection{Conflicts Should Be Refused, Not Resolved}

\begin{minted}{rust}
// The microSD card's chip select shares a pin on this board. When that driver lands,
// this facade is where the conflict is refused -- not silently arbitrated.
pub fn sd(&mut self) -> Result<SdCard<'_>, BoardError>
\end{minted}

A facade that resolves a real conflict silently is a facade that produces a bug the user cannot see.
Returning an error — or not offering the operation at all — is honest, and it is the difference between
an API that documents the hardware and one that pretends the hardware is simpler than it is.

\section{Examples Belong to the Hardware They Need}
\label{sec:bsp-examples}

The most practical BSP rule in this book is about where example files live, because that is where a
layout either documents itself or lies.

\subsection{The Rule}

\begin{quote}
    An example lives in the crate whose \textbf{smallest} hardware requirement it satisfies.
\end{quote}

\begin{table}[htbp]
    \centering
    \footnotesize
    \begin{tabular}{@{}p{0.36\textwidth}p{0.28\textwidth}p{0.32\textwidth}@{}}
        \toprule
        \textbf{Example} & \textbf{Lives in} & \textbf{Needs} \\
        \midrule
        \codeinline{module\_led\_blink} & \codeinline{xiao-mg24-bsp} & the module \\
        \codeinline{module\_imu} & \codeinline{xiao-mg24-bsp} (\codeinline{sense}) & the module, Sense variant \\
        \codeinline{base\_env\_clock} & \codeinline{xiao-expansion-base-bsp} & module + carrier \\
        \codeinline{base\_fall\_detector} & \codeinline{xiao-expansion-base-bsp} & module + carrier (panel, buzzer) \\
        \bottomrule
    \end{tabular}
    \caption{The path of the file is the bill of materials.}
    \label{tab:bsp-examples}
\end{table}

\subsection{The Three Ways to Get It Wrong}

\begin{enumerate}
    \item \textbf{Everything in the carrier crate.} Convenient, and it makes the module crate's API
          unexercised: this project's module facade had no examples for weeks, and its documentation
          described a pin that did not exist.
    \item \textbf{Everything in the module crate.} Impossible without the carrier's devices, so the
          examples would need \codeinline{cfg} guards and would prove neither level.
    \item \textbf{A flat examples/ directory at the workspace root.} The path stops telling you
          anything, and \codeinline{cargo} stops checking that the crate still builds its own examples.
\end{enumerate}

\subsection{The Check That Keeps It True}

\begin{minted}{bash}
# every crate that exposes hardware must carry at least one example of its own level
for dir in crates/*/ libraries/*/; do
  [ -f "${dir}Cargo.toml" ] || continue
  n=$(find "$dir/examples" -maxdepth 1 -name '*.rs' | wc -l)
  [ "$n" -gt 0 ] || fail "$(basename "$dir"): no examples -- its API is not built by anything"
done
\end{minted}

That check is why this project's module facade has a \codeinline{module\_led\_blink} that uses only the
module, and why the carrier's examples are in the carrier. A rule enforced by the file system is a rule
nobody has to remember.

\subsection{The Half-Example That Taught the Rule}

A fall detector was written in the HAL's examples directory. It used the IMU, the panel and the buzzer:
two boards' worth of hardware, in the crate that knows neither. It compiled, because the HAL happened to
contain the wiring at the time — which was itself the defect the split fixed.

After the split, the same program had an obvious home: the carrier crate, which depends on the module
and therefore has both the IMU (through the module) and the panel (its own). Moving it changed no code.
That is the sign of a layout that was right: the program did not have to change to be correctly placed.


\clearpage

\begin{keytakeaways}
    \item Hardware that can be owned in pieces needs a crate per piece, with the dependency arrow
          pointing one way: carrier depends on module, module depends on HAL.
    \item A variant of the same PCB is a feature, not a crate. Two crates would duplicate the boot,
          the clock tree and the LED, and one of them would drift.
    \item A BSP's API is what a user asks for — \codeinline{led()}, \codeinline{imu()},
          \codeinline{oled()} — never an address, an alias or a register.
    \item A facade that hides a real conflict (two buses, a shared pin) hides it until the user hits
          it. Hide the addresses; keep the structure visible.
    \item Every example belongs to the crate whose hardware it needs, and the gate fails a crate that
          has hardware and no example.
\end{keytakeaways}


% Chapter 16: Documentation
\chapter{Documentation Strategies and Project Management}
\label{ch:chapter-16}
\chapthumb{16}
\section*{This chapter content will be added}
\clearpage

% Chapter 17: Production
\chapter{Production Deployment and Ecosystem Integration}
\label{ch:chapter-17}
\chapthumb{17}
\section*{This chapter content will be added}
\clearpage

% Chapter 18: Complete Walkthrough
\chapter{Complete Project Walkthrough: From Zero to Production}
\label{ch:chapter-18}
\chapthumb{18}
\section*{This chapter content will be added}
\clearpage

% ========== APPENDICES ==========
\appendix

\chapter{EFR32MG24 Reference Manual Quick Guide}
\label{ap:appendix-a}
This appendix will be added.
\clearpage

\chapter{Rust Embedded Ecosystem Overview}
\label{ap:appendix-b}
This appendix will be added.
\clearpage

\chapter{Troubleshooting Common Issues}
\label{ap:appendix-c}
This appendix will be added.
\clearpage

\chapter{Complete Code Examples Repository}
\label{ap:appendix-d}
This appendix will be added.
\clearpage

\chapter{Wireless Protocol Comparison}
\label{ap:appendix-e}
This appendix will be added.
\clearpage

% Back Matter
\backmatter

% Print glossary
\printglossary[title=Glossary,toctitle={Glossary of Terms}]
\addcontentsline{toc}{chapter}{Glossary}
\clearpage

% Print bibliography
\printbibliography[heading=bibintoc]
\addcontentsline{toc}{chapter}{Bibliography}
\clearpage

% An index would need \usepackage{imakeidx} and \index{} entries; neither exists,
% and an undefined \printindex is a hard error, so the page is simply absent.

% Back cover
\thispagestyle{empty}
\begin{center}
    \vspace*{\fill}
    {\huge\bfseries End of Book}

    \vspace{2em}
    {\large For more information and updates:}

    \url{https://github.com/bitscrafts/EFR32MG2X-RS}

    \vspace{2em}
    Document built: \today

    Build version: 1.0.0
    \vspace*{\fill}
\end{center}

\end{document}
